\documentclass[11pt,letterpaper]{article}
\usepackage{geometry}
\geometry{verbose,tmargin=1in,bmargin=1in,lmargin=1in,rmargin=1in,headheight=1cm,headsep=1cm,footskip=1cm}
\usepackage{amsmath,amsthm,amssymb}
\allowdisplaybreaks
\usepackage{graphicx}
\usepackage{booktabs}
\usepackage{natbib}
\usepackage{setspace}
\usepackage{chngcntr}
\doublespacing
\usepackage{tikz}
\usepackage{pgfplots}   
\usepackage{pgfplotstable}
\pgfplotsset{compat=1.16}
\usepgfplotslibrary{dateplot}
\usepgfplotslibrary{groupplots}
\usepgfplotslibrary{fillbetween}
\usepgfplotslibrary{dateplot}
\usepgfplotslibrary{statistics}
\usepackage[colorlinks=true, linkcolor=black, citecolor=black, urlcolor=blue]{hyperref}
\usepackage{comment}
\usepackage{etoolbox} 
\usepackage{sectsty}
\usepackage{xcolor}
\usepackage[format=plain,
            justification=centering]{caption}
\usepackage{placeins}
\setcounter{topnumber}{5}
\setcounter{bottomnumber}{2}
\setcounter{totalnumber}{7}
\renewcommand{\topfraction}{0.95}
\renewcommand{\bottomfraction}{0.80}
\renewcommand{\textfraction}{0.05}
\renewcommand{\floatpagefraction}{0.70}
\usepackage{threeparttable}
\usepackage{datatool}
% Guard against silently redefining a datatool macro.  Every paper number is
% pulled with \DTLgetvalueforkey{\Macro}{...}; if two calls target the same
% \Macro the second wins and a table quietly prints the wrong series.  Wrap the
% command so a repeat target stops the build instead.
\makeatletter
\let\TR@DTLgetvalueforkey\DTLgetvalueforkey
\def\DTLgetvalueforkey#1#2#3#4#5{%
  \expandafter\ifx\csname\expandafter\@gobble\string#1\endcsname\relax\else
    \PackageError{tcja}{Duplicate datatool macro \string#1 (from db `#3')}{%
      \string#1 is already defined, so this call would silently overwrite the
      earlier value and every table using it would print the wrong series.
      Give one of the two targets a distinct name.}%
  \fi
  \TR@DTLgetvalueforkey{#1}{#2}{#3}{#4}{#5}%
}
\makeatother
\usepackage{siunitx}
\sisetup{
    group-digits = integer,
    group-separator = {,},
    group-minimum-digits = 4,
  }
\newcommand{\CountNum}[1]{\num[round-mode=places,round-precision=0]{#1}}
\newcommand{\PctNum}[1]{\num[round-mode=places,round-precision=1]{#1}\%}
\newcommand{\ShareNum}[1]{\num[round-mode=places,round-precision=1]{#1}}
\newcommand{\CoefNum}[1]{\num[round-mode=places,round-precision=3]{#1}}
\newcommand{\SENum}[1]{(\num[round-mode=places,round-precision=3]{#1})}
% Significance stars sit in a box of fixed width, so a starred coefficient and
% an unstarred one right-align on the same digit.  Padding with \phantom{***}
% instead adds the stars' width on top of the stars themselves, which pushes
% the significant rows left of the insignificant ones.
\newlength{\StarWidth}
\settowidth{\StarWidth}{***}
\newcommand{\StarBox}[1]{\makebox[\StarWidth][l]{#1}}
\newcommand{\RatioNum}[1]{\num[round-mode=places,round-precision=2]{#1}}
\newcommand{\RoundNum}[1]{\num[round-mode=places,round-precision=0]{#1}}
\newcommand{\RatioNumAbs}[1]{\num[round-mode=places,round-precision=2]{#1}}

\definecolor{IntRed}{RGB}{220, 53, 69}
\definecolor{OneOrange}{RGB}{247, 127, 0}
\definecolor{FedGold}{RGB}{245, 190, 44}
\definecolor{StateGreen}{RGB}{40, 167, 69}
\definecolor{ValTeal}{RGB}{23, 162, 184}
\definecolor{BookBlue}{RGB}{56, 175, 237}
\definecolor{OthViolet}{RGB}{155, 93, 181}
\definecolor{Indpre}{RGB}{120,90,170} 
\definecolor{Indpost}{RGB}{230,150,60}  

\newcommand{\Inputs}{Inputs}

\DTLloaddb{samplesanddescriptive}{\Inputs/slide_data/samplesanddescriptive.csv}
\DTLgetvalueforkey{\ProfitPerPpClosure}{value}{samplesanddescriptive}{key}{profitPerPpClosure}
\DTLgetvalueforkey{\FdiiClaimantsCount}{value}{samplesanddescriptive}{key}{fdiiClaimantsCount}
\DTLgetvalueforkey{\FdiiNonClaimantsCount}{value}{samplesanddescriptive}{key}{fdiiNonClaimantsCount}
\DTLgetvalueforkey{\WorkedExampleEtrPp}{value}{samplesanddescriptive}{key}{workedExampleEtrPp}
\DTLgetvalueforkey{\WorkedExampleIntlPp}{value}{samplesanddescriptive}{key}{workedExampleIntlPp}
\DTLgetvalueforkey{\WorkedExampleGapPp}{value}{samplesanddescriptive}{key}{workedExampleGapPp}
\DTLgetvalueforkey{\WorkedExampleFromPp}{value}{samplesanddescriptive}{key}{workedExampleFromPp}
\DTLgetvalueforkey{\WorkedExampleToPp}{value}{samplesanddescriptive}{key}{workedExampleToPp}
\DTLgetvalueforkey{\WorkedExampleOtherPp}{value}{samplesanddescriptive}{key}{workedExampleOtherPp}
\DTLgetvalueforkey{\PassthroughEtrPct}{value}{samplesanddescriptive}{key}{passthroughEtrPct}
\DTLgetvalueforkey{\PassthroughLeakagePct}{value}{samplesanddescriptive}{key}{passthroughLeakagePct}

\DTLloaddb{cashandcapital}{\Inputs/slide_data/cashandcapital.csv}
\DTLgetvalueforkey{\CapLocCapxPerPp}{value}{cashandcapital}{key}{capLocCapxPerPp}

\DTLloaddb{pipelinecounts}{\Inputs/slide_data/pipelinecounts.csv}
\DTLgetvalueforkey{\InputUrlsTotal}{value}{pipelinecounts}{key}{inputUrlsTotal}
\DTLgetvalueforkey{\InputUrlsSkipNonHtml}{value}{pipelinecounts}{key}{inputUrlsSkipNonHtml}
\DTLgetvalueforkey{\InputUrlsFail}{value}{pipelinecounts}{key}{inputUrlsFail}
\DTLgetvalueforkey{\InputUrlsDuplicates}{value}{pipelinecounts}{key}{inputUrlsDuplicates}
\DTLgetvalueforkey{\PhaseOneUniqueFilings}{value}{pipelinecounts}{key}{phaseOneUniqueFilings}
\DTLgetvalueforkey{\PhaseOneStatusOK}{value}{pipelinecounts}{key}{phaseOneStatusOk}
\DTLgetvalueforkey{\PhaseOneBucketHighConf}{value}{pipelinecounts}{key}{phaseOneBucketHighConf}
\DTLgetvalueforkey{\PhaseOneBucketLikely}{value}{pipelinecounts}{key}{phaseOneBucketLikely}
\DTLgetvalueforkey{\PhaseOneBucketAmbiguous}{value}{pipelinecounts}{key}{phaseOneBucketAmbiguous}
\DTLgetvalueforkey{\PhaseOneBucketNarrative}{value}{pipelinecounts}{key}{phaseOneBucketNarrative}
\DTLgetvalueforkey{\PhaseOneBucketNoTable}{value}{pipelinecounts}{key}{phaseOneBucketNoTable}
\DTLgetvalueforkey{\PhaseOneBucketWrongTable}{value}{pipelinecounts}{key}{phaseOneBucketWrongTable}
\DTLgetvalueforkey{\PhaseOneExtractableSum}{value}{pipelinecounts}{key}{phaseOneExtractableSum}
\DTLgetvalueforkey{\PhaseTwoTotalRows}{value}{pipelinecounts}{key}{phaseTwoTotalRows}
\DTLgetvalueforkey{\PhaseTwoCheckTrue}{value}{pipelinecounts}{key}{phaseTwoCheckTrue}
\DTLgetvalueforkey{\PhaseTwoCheckFalse}{value}{pipelinecounts}{key}{phaseTwoCheckFalse}
\DTLgetvalueforkey{\PhaseTwoDeterministic}{value}{pipelinecounts}{key}{phaseTwoDeterministic}
\DTLgetvalueforkey{\PhaseTwoLlmFallback}{value}{pipelinecounts}{key}{phaseTwoLlmFallback}
\DTLgetvalueforkey{\PhaseTwoRateForm}{value}{pipelinecounts}{key}{phaseTwoRateForm}
\DTLgetvalueforkey{\PhaseTwoAmountForm}{value}{pipelinecounts}{key}{phaseTwoAmountForm}
\DTLgetvalueforkey{\PhaseTwoUniqueFilings}{value}{pipelinecounts}{key}{phaseTwoUniqueFilings}
\DTLgetvalueforkey{\PacketsProcessed}{value}{pipelinecounts}{key}{packetsProcessed}
\DTLgetvalueforkey{\ProcessedWithRows}{value}{pipelinecounts}{key}{processedWithRows}
\DTLgetvalueforkey{\ProcessedNoRows}{value}{pipelinecounts}{key}{processedNoRows}
\DTLgetvalueforkey{\FilingsOutsideMainRun}{value}{pipelinecounts}{key}{filingsOutsideMainRun}
\DTLgetvalueforkey{\BasePanelFilingYears}{value}{pipelinecounts}{key}{basePanelFilingYears}
\DTLgetvalueforkey{\NetFilingMatch}{value}{pipelinecounts}{key}{netFilingMatch}
\DTLgetvalueforkey{\DropYearRange}{value}{pipelinecounts}{key}{dropYearRange}
\DTLgetvalueforkey{\DropFilingWindow}{value}{pipelinecounts}{key}{dropFilingWindow}
\DTLgetvalueforkey{\DropNoCompustat}{value}{pipelinecounts}{key}{dropNoCompustat}
\DTLgetvalueforkey{\DropCheckFail}{value}{pipelinecounts}{key}{dropCheckFail}
\DTLgetvalueforkey{\DropDuplicateFirmYears}{value}{pipelinecounts}{key}{dropDuplicateFirmYears}
\DTLgetvalueforkey{\PanelCanonicalRows}{value}{pipelinecounts}{key}{panelCanonicalRows}
\DTLgetvalueforkey{\AnalysisPanelFirmYears}{value}{pipelinecounts}{key}{analysisPanelFirmYears}


\DTLloaddb{regressionheadlines}{\Inputs/slide_data/regressionheadlines.csv}
\DTLgetvalueforkey{\WindfallThetaCapex}{value}{regressionheadlines}{key}{windfallThetaCapex}
\DTLgetvalueforkey{\WindfallThetaCapexSe}{value}{regressionheadlines}{key}{windfallThetaCapexSe}
\DTLgetvalueforkey{\WindfallThetaRd}{value}{regressionheadlines}{key}{windfallThetaRd}
\DTLgetvalueforkey{\WindfallThetaRdSe}{value}{regressionheadlines}{key}{windfallThetaRdSe}


\newcommand{\PhaseOneParsedPct}{%
    \pgfmathparse{100*\PhaseOneStatusOK/\PhaseOneUniqueFilings}%
    \pgfmathresult%
}
\newcommand{\NoPhaseOneExtraction}{%
    \number\numexpr\PhaseOneStatusOK-\PhaseOneExtractableSum\relax
}
\newcommand{\NoReconRows}{%
    \number\numexpr\PhaseOneExtractableSum-\PhaseTwoUniqueFilings\relax
}
\newcommand{\DedupPanelDropRows}{%
    \number\numexpr\PhaseTwoCheckTrue-\PanelCanonicalRows\relax
}
\newcommand{\OtherExclusions}{%
    \number\numexpr\PanelCanonicalRows-\AnalysisPanelFirmYears\relax
}
\newcommand{\AbsCoefNum}[1]{%
    \pgfmathparse{abs(#1)}%
    \num[round-mode=places,round-precision=3]{\pgfmathresult}%
}
\newcommand{\AbsRatioNum}[1]{%
    \pgfmathparse{abs(#1)}%
    \num[round-mode=places,round-precision=2]{\pgfmathresult}%
}


\DTLloaddb{reconextract}{\Inputs/tables/reconciliation_extraction_headlines.csv}
\DTLgetvalueforkey{\ReconUniqueFilings}{value}{reconextract}{row_key}{unique_filings}
\DTLgetvalueforkey{\ReconParsedFilings}{value}{reconextract}{row_key}{parsed_filings}
\DTLgetvalueforkey{\ReconParsedPct}{display_1dp}{reconextract}{row_key}{parsed_share_pct}
\DTLgetvalueforkey{\ReconExtractablePackets}{value}{reconextract}{row_key}{extractable_packets}
\DTLgetvalueforkey{\ReconPhaseTwoFilings}{value}{reconextract}{row_key}{phase2_extracted_filings}
\DTLgetvalueforkey{\ReconPhaseTwoFirmYears}{value}{reconextract}{row_key}{phase2_firmyears}

\DTLloaddb{coverage}{\Inputs/tables/coverage_table.csv}
\DTLgetvalueforkey{\TopHundredTotal}{compustat_firm_years}{coverage}{row_key}{top100}
\DTLgetvalueforkey{\TopTwoFiftyTotal}{compustat_firm_years}{coverage}{row_key}{top250}
\DTLgetvalueforkey{\TopFiveHundredTotal}{compustat_firm_years}{coverage}{row_key}{top500}
\DTLgetvalueforkey{\FullSampleTotal}{compustat_firm_years}{coverage}{row_key}{full_sample}
\DTLgetvalueforkey{\TopHundredCovered}{covered_firm_years}{coverage}{row_key}{top100}
\DTLgetvalueforkey{\TopHundredEq}{equal_weight_share_pct}{coverage}{row_key}{top100}
\DTLgetvalueforkey{\TopHundredSales}{sales_weight_share_pct}{coverage}{row_key}{top100}
\DTLgetvalueforkey{\TopTwoFiftyCovered}{covered_firm_years}{coverage}{row_key}{top250}
\DTLgetvalueforkey{\TopTwoFiftyEq}{equal_weight_share_pct}{coverage}{row_key}{top250}
\DTLgetvalueforkey{\TopTwoFiftySales}{sales_weight_share_pct}{coverage}{row_key}{top250}
\DTLgetvalueforkey{\TopFiveHundredCovered}{covered_firm_years}{coverage}{row_key}{top500}
\DTLgetvalueforkey{\TopFiveHundredEq}{equal_weight_share_pct}{coverage}{row_key}{top500}
\DTLgetvalueforkey{\TopFiveHundredSales}{sales_weight_share_pct}{coverage}{row_key}{top500}
\DTLgetvalueforkey{\FullSampleCovered}{covered_firm_years}{coverage}{row_key}{full_sample}
\DTLgetvalueforkey{\FullSampleEq}{equal_weight_share_pct}{coverage}{row_key}{full_sample}
\DTLgetvalueforkey{\FullSampleSales}{sales_weight_share_pct}{coverage}{row_key}{full_sample}

\DTLloaddb{reconcovtime}{\Inputs/tables/reconciliation_coverage_time_headlines.csv}
\DTLgetvalueforkey{\ReconCovInit}{coverage_count_pct}{reconcovtime}{row_key}{all_1998}
\DTLgetvalueforkey{\ReconCovGR}{coverage_count_pct}{reconcovtime}{row_key}{all_2008}

\DTLloaddb{contribdist}{\Inputs/tables/contribution_gap_distribution_lookup.csv}
\DTLgetvalueforkey{\PreTotalCounts}{counts}{contribdist}{lookup_key}{pre_total_gap}
\DTLgetvalueforkey{\PreTotalAvg}{average}{contribdist}{lookup_key}{pre_total_gap}
\DTLgetvalueforkey{\PreTotalPTF}{p25}{contribdist}{lookup_key}{pre_total_gap}
\DTLgetvalueforkey{\PreTotalPFZ}{median}{contribdist}{lookup_key}{pre_total_gap}
\DTLgetvalueforkey{\PreTotalPSF}{p75}{contribdist}{lookup_key}{pre_total_gap}
\DTLgetvalueforkey{\PreIntlCounts}{counts}{contribdist}{lookup_key}{pre_international}
\DTLgetvalueforkey{\PreIntlAvg}{average}{contribdist}{lookup_key}{pre_international}
\DTLgetvalueforkey{\PreIntlPTF}{p25}{contribdist}{lookup_key}{pre_international}
\DTLgetvalueforkey{\PreIntlPFZ}{median}{contribdist}{lookup_key}{pre_international}
\DTLgetvalueforkey{\PreIntlPSF}{p75}{contribdist}{lookup_key}{pre_international}
\DTLgetvalueforkey{\PreOneTimeCounts}{counts}{contribdist}{lookup_key}{pre_one_time}
\DTLgetvalueforkey{\PreOneTimeAvg}{average}{contribdist}{lookup_key}{pre_one_time}
\DTLgetvalueforkey{\PreOneTimePTF}{p25}{contribdist}{lookup_key}{pre_one_time}
\DTLgetvalueforkey{\PreOneTimePFZ}{median}{contribdist}{lookup_key}{pre_one_time}
\DTLgetvalueforkey{\PreOneTimePSF}{p75}{contribdist}{lookup_key}{pre_one_time}
\DTLgetvalueforkey{\PreFedCredCounts}{counts}{contribdist}{lookup_key}{pre_federal_credits}
\DTLgetvalueforkey{\PreFedCredAvg}{average}{contribdist}{lookup_key}{pre_federal_credits}
\DTLgetvalueforkey{\PreFedCredPTF}{p25}{contribdist}{lookup_key}{pre_federal_credits}
\DTLgetvalueforkey{\PreFedCredPFZ}{median}{contribdist}{lookup_key}{pre_federal_credits}
\DTLgetvalueforkey{\PreFedCredPSF}{p75}{contribdist}{lookup_key}{pre_federal_credits}
\DTLgetvalueforkey{\PreStateCounts}{counts}{contribdist}{lookup_key}{pre_state_local}
\DTLgetvalueforkey{\PreStateAvg}{average}{contribdist}{lookup_key}{pre_state_local}
\DTLgetvalueforkey{\PreStatePTF}{p25}{contribdist}{lookup_key}{pre_state_local}
\DTLgetvalueforkey{\PreStatePFZ}{median}{contribdist}{lookup_key}{pre_state_local}
\DTLgetvalueforkey{\PreStatePSF}{p75}{contribdist}{lookup_key}{pre_state_local}
\DTLgetvalueforkey{\PreValCounts}{counts}{contribdist}{lookup_key}{pre_valuation}
\DTLgetvalueforkey{\PreValAvg}{average}{contribdist}{lookup_key}{pre_valuation}
\DTLgetvalueforkey{\PreValPTF}{p25}{contribdist}{lookup_key}{pre_valuation}
\DTLgetvalueforkey{\PreValPFZ}{median}{contribdist}{lookup_key}{pre_valuation}
\DTLgetvalueforkey{\PreValPSF}{p75}{contribdist}{lookup_key}{pre_valuation}
\DTLgetvalueforkey{\PreBookTaxCounts}{counts}{contribdist}{lookup_key}{pre_book_v_tax}
\DTLgetvalueforkey{\PreBookTaxAvg}{average}{contribdist}{lookup_key}{pre_book_v_tax}
\DTLgetvalueforkey{\PreBookTaxPTF}{p25}{contribdist}{lookup_key}{pre_book_v_tax}
\DTLgetvalueforkey{\PreBookTaxPFZ}{median}{contribdist}{lookup_key}{pre_book_v_tax}
\DTLgetvalueforkey{\PreBookTaxPSF}{p75}{contribdist}{lookup_key}{pre_book_v_tax}
\DTLgetvalueforkey{\PreOtherCounts}{counts}{contribdist}{lookup_key}{pre_other}
\DTLgetvalueforkey{\PreOtherAvg}{average}{contribdist}{lookup_key}{pre_other}
\DTLgetvalueforkey{\PreOtherPTF}{p25}{contribdist}{lookup_key}{pre_other}
\DTLgetvalueforkey{\PreOtherPFZ}{median}{contribdist}{lookup_key}{pre_other}
\DTLgetvalueforkey{\PreOtherPSF}{p75}{contribdist}{lookup_key}{pre_other}
\DTLgetvalueforkey{\PostTotalCounts}{counts}{contribdist}{lookup_key}{post_total_gap}
\DTLgetvalueforkey{\PostTotalAvg}{average}{contribdist}{lookup_key}{post_total_gap}
\DTLgetvalueforkey{\PostTotalPTF}{p25}{contribdist}{lookup_key}{post_total_gap}
\DTLgetvalueforkey{\PostTotalPFZ}{median}{contribdist}{lookup_key}{post_total_gap}
\DTLgetvalueforkey{\PostTotalPSF}{p75}{contribdist}{lookup_key}{post_total_gap}
\DTLgetvalueforkey{\PostIntlCounts}{counts}{contribdist}{lookup_key}{post_international}
\DTLgetvalueforkey{\PostIntlAvg}{average}{contribdist}{lookup_key}{post_international}
\DTLgetvalueforkey{\PostIntlPTF}{p25}{contribdist}{lookup_key}{post_international}
\DTLgetvalueforkey{\PostIntlPFZ}{median}{contribdist}{lookup_key}{post_international}
\DTLgetvalueforkey{\PostIntlPSF}{p75}{contribdist}{lookup_key}{post_international}
\DTLgetvalueforkey{\PostOneTimeCounts}{counts}{contribdist}{lookup_key}{post_one_time}
\DTLgetvalueforkey{\PostOneTimeAvg}{average}{contribdist}{lookup_key}{post_one_time}
\DTLgetvalueforkey{\PostOneTimePTF}{p25}{contribdist}{lookup_key}{post_one_time}
\DTLgetvalueforkey{\PostOneTimePFZ}{median}{contribdist}{lookup_key}{post_one_time}
\DTLgetvalueforkey{\PostOneTimePSF}{p75}{contribdist}{lookup_key}{post_one_time}
\DTLgetvalueforkey{\PostFedCredCounts}{counts}{contribdist}{lookup_key}{post_federal_credits}
\DTLgetvalueforkey{\PostFedCredAvg}{average}{contribdist}{lookup_key}{post_federal_credits}
\DTLgetvalueforkey{\PostFedCredPTF}{p25}{contribdist}{lookup_key}{post_federal_credits}
\DTLgetvalueforkey{\PostFedCredPFZ}{median}{contribdist}{lookup_key}{post_federal_credits}
\DTLgetvalueforkey{\PostFedCredPSF}{p75}{contribdist}{lookup_key}{post_federal_credits}
\DTLgetvalueforkey{\PostStateCounts}{counts}{contribdist}{lookup_key}{post_state_local}
\DTLgetvalueforkey{\PostStateAvg}{average}{contribdist}{lookup_key}{post_state_local}
\DTLgetvalueforkey{\PostStatePTF}{p25}{contribdist}{lookup_key}{post_state_local}
\DTLgetvalueforkey{\PostStatePFZ}{median}{contribdist}{lookup_key}{post_state_local}
\DTLgetvalueforkey{\PostStatePSF}{p75}{contribdist}{lookup_key}{post_state_local}
\DTLgetvalueforkey{\PostValCounts}{counts}{contribdist}{lookup_key}{post_valuation}
\DTLgetvalueforkey{\PostValAvg}{average}{contribdist}{lookup_key}{post_valuation}
\DTLgetvalueforkey{\PostValPTF}{p25}{contribdist}{lookup_key}{post_valuation}
\DTLgetvalueforkey{\PostValPFZ}{median}{contribdist}{lookup_key}{post_valuation}
\DTLgetvalueforkey{\PostValPSF}{p75}{contribdist}{lookup_key}{post_valuation}
\DTLgetvalueforkey{\PostBookTaxCounts}{counts}{contribdist}{lookup_key}{post_book_v_tax}
\DTLgetvalueforkey{\PostBookTaxAvg}{average}{contribdist}{lookup_key}{post_book_v_tax}
\DTLgetvalueforkey{\PostBookTaxPTF}{p25}{contribdist}{lookup_key}{post_book_v_tax}
\DTLgetvalueforkey{\PostBookTaxPFZ}{median}{contribdist}{lookup_key}{post_book_v_tax}
\DTLgetvalueforkey{\PostBookTaxPSF}{p75}{contribdist}{lookup_key}{post_book_v_tax}
\DTLgetvalueforkey{\PostOtherCounts}{counts}{contribdist}{lookup_key}{post_other}
\DTLgetvalueforkey{\PostOtherAvg}{average}{contribdist}{lookup_key}{post_other}
\DTLgetvalueforkey{\PostOtherPTF}{p25}{contribdist}{lookup_key}{post_other}
\DTLgetvalueforkey{\PostOtherPFZ}{median}{contribdist}{lookup_key}{post_other}
\DTLgetvalueforkey{\PostOtherPSF}{p75}{contribdist}{lookup_key}{post_other}

\DTLloaddb{figoneyear}{\Inputs/decomp_tikz/fig_drivers_tikz_data.csv}
\DTLgetvalueforkey{\FigOneTimeInit}{tcja_onetime}{figoneyear}{year}{2004}
\DTLgetvalueforkey{\FigIntlInit}{international}{figoneyear}{year}{2004}

\DTLloaddb{figoneavg}{\Inputs/tables/contribution_gap_distribution_lookup.csv}
\DTLgetvalueforkey{\FigPreTotalAvg}{average}{figoneavg}{lookup_key}{pre_total_gap}
\DTLgetvalueforkey{\FigPreIntlAvg}{average}{figoneavg}{lookup_key}{pre_international}
\DTLgetvalueforkey{\FigPostTotalAvg}{average}{figoneavg}{lookup_key}{post_total_gap}
\DTLgetvalueforkey{\FigPostIntlAvg}{average}{figoneavg}{lookup_key}{post_international}

\DTLloaddb{figonechg}{\Inputs/tables/contribution_gap_changes.csv}
\DTLgetvalueforkey{\FigIntlChange}{change}{figonechg}{row_key}{international}
\DTLgetvalueforkey{\FigMaxNonIntlChange}{abs_change}{figonechg}{row_key}{max_non_international}

\DTLloaddb{figtwo}{\Inputs/decomp_ind_new_tikz/summary.csv}
\DTLgetvalueforkey{\FigTwoTechPreGap}{gap_pre_pp}{figtwo}{industry}{Tech}
\DTLgetvalueforkey{\FigTwoTechPostGap}{gap_post_pp}{figtwo}{industry}{Tech}
\DTLgetvalueforkey{\FigTwoTechIntlDelta}{international_delta_pp}{figtwo}{industry}{Tech}
\DTLgetvalueforkey{\FigTwoPharmaIntlDelta}{international_delta_pp}{figtwo}{industry}{Pharma/Biotech}
\DTLgetvalueforkey{\FigTwoManufIntlDelta}{international_delta_pp}{figtwo}{industry}{Manufacturing}
% Oil & Gas cannot be matched on the industry column (the ampersand does not
% survive as a Datatool key), so this one row is fetched by its plot position.
\DTLgetvalueforkey{\FigTwoOilIntlDelta}{international_delta_pp}{figtwo}{y}{3}

\DTLloaddb{industrydist}{\Inputs/tables/industry_gap_distribution_lookup.csv}
\DTLgetvalueforkey{\PreTechCounts}{counts}{industrydist}{lookup_key}{pre_tech}
\DTLgetvalueforkey{\PreTechAvg}{average}{industrydist}{lookup_key}{pre_tech}
\DTLgetvalueforkey{\PreTechPTF}{p25}{industrydist}{lookup_key}{pre_tech}
\DTLgetvalueforkey{\PreTechPFZ}{median}{industrydist}{lookup_key}{pre_tech}
\DTLgetvalueforkey{\PreTechPSF}{p75}{industrydist}{lookup_key}{pre_tech}
\DTLgetvalueforkey{\PrePharmaCounts}{counts}{industrydist}{lookup_key}{pre_pharma_biotech}
\DTLgetvalueforkey{\PrePharmaAvg}{average}{industrydist}{lookup_key}{pre_pharma_biotech}
\DTLgetvalueforkey{\PrePharmaPTF}{p25}{industrydist}{lookup_key}{pre_pharma_biotech}
\DTLgetvalueforkey{\PrePharmaPFZ}{median}{industrydist}{lookup_key}{pre_pharma_biotech}
\DTLgetvalueforkey{\PrePharmaPSF}{p75}{industrydist}{lookup_key}{pre_pharma_biotech}
\DTLgetvalueforkey{\PreManufCounts}{counts}{industrydist}{lookup_key}{pre_manufacturing}
\DTLgetvalueforkey{\PreManufAvg}{average}{industrydist}{lookup_key}{pre_manufacturing}
\DTLgetvalueforkey{\PreManufPTF}{p25}{industrydist}{lookup_key}{pre_manufacturing}
\DTLgetvalueforkey{\PreManufPFZ}{median}{industrydist}{lookup_key}{pre_manufacturing}
\DTLgetvalueforkey{\PreManufPSF}{p75}{industrydist}{lookup_key}{pre_manufacturing}
\DTLgetvalueforkey{\PreOilCounts}{counts}{industrydist}{lookup_key}{pre_oil_and_gas}
\DTLgetvalueforkey{\PreOilAvg}{average}{industrydist}{lookup_key}{pre_oil_and_gas}
\DTLgetvalueforkey{\PreOilPTF}{p25}{industrydist}{lookup_key}{pre_oil_and_gas}
\DTLgetvalueforkey{\PreOilPFZ}{median}{industrydist}{lookup_key}{pre_oil_and_gas}
\DTLgetvalueforkey{\PreOilPSF}{p75}{industrydist}{lookup_key}{pre_oil_and_gas}
\DTLgetvalueforkey{\PreRetailCounts}{counts}{industrydist}{lookup_key}{pre_retail}
\DTLgetvalueforkey{\PreRetailAvg}{average}{industrydist}{lookup_key}{pre_retail}
\DTLgetvalueforkey{\PreRetailPTF}{p25}{industrydist}{lookup_key}{pre_retail}
\DTLgetvalueforkey{\PreRetailPFZ}{median}{industrydist}{lookup_key}{pre_retail}
\DTLgetvalueforkey{\PreRetailPSF}{p75}{industrydist}{lookup_key}{pre_retail}
\DTLgetvalueforkey{\PreUtilCounts}{counts}{industrydist}{lookup_key}{pre_utilities}
\DTLgetvalueforkey{\PreUtilAvg}{average}{industrydist}{lookup_key}{pre_utilities}
\DTLgetvalueforkey{\PreUtilPTF}{p25}{industrydist}{lookup_key}{pre_utilities}
\DTLgetvalueforkey{\PreUtilPFZ}{median}{industrydist}{lookup_key}{pre_utilities}
\DTLgetvalueforkey{\PreUtilPSF}{p75}{industrydist}{lookup_key}{pre_utilities}
\DTLgetvalueforkey{\PreIndOtherCounts}{counts}{industrydist}{lookup_key}{pre_other}
\DTLgetvalueforkey{\PreIndOtherAvg}{average}{industrydist}{lookup_key}{pre_other}
\DTLgetvalueforkey{\PreIndOtherPTF}{p25}{industrydist}{lookup_key}{pre_other}
\DTLgetvalueforkey{\PreIndOtherPFZ}{median}{industrydist}{lookup_key}{pre_other}
\DTLgetvalueforkey{\PreIndOtherPSF}{p75}{industrydist}{lookup_key}{pre_other}
\DTLgetvalueforkey{\PostTechCounts}{counts}{industrydist}{lookup_key}{post_tech}
\DTLgetvalueforkey{\PostTechAvg}{average}{industrydist}{lookup_key}{post_tech}
\DTLgetvalueforkey{\PostTechPTF}{p25}{industrydist}{lookup_key}{post_tech}
\DTLgetvalueforkey{\PostTechPFZ}{median}{industrydist}{lookup_key}{post_tech}
\DTLgetvalueforkey{\PostTechPSF}{p75}{industrydist}{lookup_key}{post_tech}
\DTLgetvalueforkey{\PostPharmaCounts}{counts}{industrydist}{lookup_key}{post_pharma_biotech}
\DTLgetvalueforkey{\PostPharmaAvg}{average}{industrydist}{lookup_key}{post_pharma_biotech}
\DTLgetvalueforkey{\PostPharmaPTF}{p25}{industrydist}{lookup_key}{post_pharma_biotech}
\DTLgetvalueforkey{\PostPharmaPFZ}{median}{industrydist}{lookup_key}{post_pharma_biotech}
\DTLgetvalueforkey{\PostPharmaPSF}{p75}{industrydist}{lookup_key}{post_pharma_biotech}
\DTLgetvalueforkey{\PostManufCounts}{counts}{industrydist}{lookup_key}{post_manufacturing}
\DTLgetvalueforkey{\PostManufAvg}{average}{industrydist}{lookup_key}{post_manufacturing}
\DTLgetvalueforkey{\PostManufPTF}{p25}{industrydist}{lookup_key}{post_manufacturing}
\DTLgetvalueforkey{\PostManufPFZ}{median}{industrydist}{lookup_key}{post_manufacturing}
\DTLgetvalueforkey{\PostManufPSF}{p75}{industrydist}{lookup_key}{post_manufacturing}
\DTLgetvalueforkey{\PostOilCounts}{counts}{industrydist}{lookup_key}{post_oil_and_gas}
\DTLgetvalueforkey{\PostOilAvg}{average}{industrydist}{lookup_key}{post_oil_and_gas}
\DTLgetvalueforkey{\PostOilPTF}{p25}{industrydist}{lookup_key}{post_oil_and_gas}
\DTLgetvalueforkey{\PostOilPFZ}{median}{industrydist}{lookup_key}{post_oil_and_gas}
\DTLgetvalueforkey{\PostOilPSF}{p75}{industrydist}{lookup_key}{post_oil_and_gas}
\DTLgetvalueforkey{\PostRetailCounts}{counts}{industrydist}{lookup_key}{post_retail}
\DTLgetvalueforkey{\PostRetailAvg}{average}{industrydist}{lookup_key}{post_retail}
\DTLgetvalueforkey{\PostRetailPTF}{p25}{industrydist}{lookup_key}{post_retail}
\DTLgetvalueforkey{\PostRetailPFZ}{median}{industrydist}{lookup_key}{post_retail}
\DTLgetvalueforkey{\PostRetailPSF}{p75}{industrydist}{lookup_key}{post_retail}
\DTLgetvalueforkey{\PostUtilCounts}{counts}{industrydist}{lookup_key}{post_utilities}
\DTLgetvalueforkey{\PostUtilAvg}{average}{industrydist}{lookup_key}{post_utilities}
\DTLgetvalueforkey{\PostUtilPTF}{p25}{industrydist}{lookup_key}{post_utilities}
\DTLgetvalueforkey{\PostUtilPFZ}{median}{industrydist}{lookup_key}{post_utilities}
\DTLgetvalueforkey{\PostUtilPSF}{p75}{industrydist}{lookup_key}{post_utilities}
\DTLgetvalueforkey{\PostIndOtherCounts}{counts}{industrydist}{lookup_key}{post_other}
\DTLgetvalueforkey{\PostIndOtherAvg}{average}{industrydist}{lookup_key}{post_other}
\DTLgetvalueforkey{\PostIndOtherPTF}{p25}{industrydist}{lookup_key}{post_other}
\DTLgetvalueforkey{\PostIndOtherPFZ}{median}{industrydist}{lookup_key}{post_other}
\DTLgetvalueforkey{\PostIndOtherPSF}{p75}{industrydist}{lookup_key}{post_other}

\DTLloaddb{identification}{\Inputs/identification_tikz/intl_pre_by_industry.csv}
\DTLgetvalueforkey{\AllInSampleN}{n_firms}{identification}{row_key}{all_in_sample}
\DTLgetvalueforkey{\AllInSampleMean}{mean}{identification}{row_key}{all_in_sample}
\DTLgetvalueforkey{\AllInSampleSD}{sd}{identification}{row_key}{all_in_sample}
\DTLgetvalueforkey{\AllInSamplePTF}{p25}{identification}{row_key}{all_in_sample}
\DTLgetvalueforkey{\AllInSamplePFZ}{median}{identification}{row_key}{all_in_sample}
\DTLgetvalueforkey{\AllInSamplePSF}{p75}{identification}{row_key}{all_in_sample}
\DTLgetvalueforkey{\AllInSampleNonzero}{share_nonzero_pct}{identification}{row_key}{all_in_sample}
\DTLgetvalueforkey{\TechN}{n_firms}{identification}{row_key}{tech}
\DTLgetvalueforkey{\TechMean}{mean}{identification}{row_key}{tech}
\DTLgetvalueforkey{\TechSD}{sd}{identification}{row_key}{tech}
\DTLgetvalueforkey{\TechPTF}{p25}{identification}{row_key}{tech}
\DTLgetvalueforkey{\TechPFZ}{median}{identification}{row_key}{tech}
\DTLgetvalueforkey{\TechPSF}{p75}{identification}{row_key}{tech}
\DTLgetvalueforkey{\TechNonzero}{share_nonzero_pct}{identification}{row_key}{tech}
\DTLgetvalueforkey{\PharmaN}{n_firms}{identification}{row_key}{pharma_biotech}
\DTLgetvalueforkey{\PharmaMean}{mean}{identification}{row_key}{pharma_biotech}
\DTLgetvalueforkey{\PharmaSD}{sd}{identification}{row_key}{pharma_biotech}
\DTLgetvalueforkey{\PharmaPTF}{p25}{identification}{row_key}{pharma_biotech}
\DTLgetvalueforkey{\PharmaPFZ}{median}{identification}{row_key}{pharma_biotech}
\DTLgetvalueforkey{\PharmaPSF}{p75}{identification}{row_key}{pharma_biotech}
\DTLgetvalueforkey{\PharmaNonzero}{share_nonzero_pct}{identification}{row_key}{pharma_biotech}
\DTLgetvalueforkey{\ManufN}{n_firms}{identification}{row_key}{manufacturing}
\DTLgetvalueforkey{\ManufMean}{mean}{identification}{row_key}{manufacturing}
\DTLgetvalueforkey{\ManufSD}{sd}{identification}{row_key}{manufacturing}
\DTLgetvalueforkey{\ManufPTF}{p25}{identification}{row_key}{manufacturing}
\DTLgetvalueforkey{\ManufPFZ}{median}{identification}{row_key}{manufacturing}
\DTLgetvalueforkey{\ManufPSF}{p75}{identification}{row_key}{manufacturing}
\DTLgetvalueforkey{\ManufNonzero}{share_nonzero_pct}{identification}{row_key}{manufacturing}
\DTLgetvalueforkey{\OilN}{n_firms}{identification}{row_key}{oil_and_gas}
\DTLgetvalueforkey{\OilMean}{mean}{identification}{row_key}{oil_and_gas}
\DTLgetvalueforkey{\OilSD}{sd}{identification}{row_key}{oil_and_gas}
\DTLgetvalueforkey{\OilPTF}{p25}{identification}{row_key}{oil_and_gas}
\DTLgetvalueforkey{\OilPFZ}{median}{identification}{row_key}{oil_and_gas}
\DTLgetvalueforkey{\OilPSF}{p75}{identification}{row_key}{oil_and_gas}
\DTLgetvalueforkey{\OilNonzero}{share_nonzero_pct}{identification}{row_key}{oil_and_gas}
\DTLgetvalueforkey{\RetailN}{n_firms}{identification}{row_key}{retail}
\DTLgetvalueforkey{\RetailMean}{mean}{identification}{row_key}{retail}
\DTLgetvalueforkey{\RetailSD}{sd}{identification}{row_key}{retail}
\DTLgetvalueforkey{\RetailPTF}{p25}{identification}{row_key}{retail}
\DTLgetvalueforkey{\RetailPFZ}{median}{identification}{row_key}{retail}
\DTLgetvalueforkey{\RetailPSF}{p75}{identification}{row_key}{retail}
\DTLgetvalueforkey{\RetailNonzero}{share_nonzero_pct}{identification}{row_key}{retail}
\DTLgetvalueforkey{\UtilitiesN}{n_firms}{identification}{row_key}{utilities}
\DTLgetvalueforkey{\UtilitiesMean}{mean}{identification}{row_key}{utilities}
\DTLgetvalueforkey{\UtilitiesSD}{sd}{identification}{row_key}{utilities}
\DTLgetvalueforkey{\UtilitiesPTF}{p25}{identification}{row_key}{utilities}
\DTLgetvalueforkey{\UtilitiesPFZ}{median}{identification}{row_key}{utilities}
\DTLgetvalueforkey{\UtilitiesPSF}{p75}{identification}{row_key}{utilities}
\DTLgetvalueforkey{\UtilitiesNonzero}{share_nonzero_pct}{identification}{row_key} {utilities}
\DTLgetvalueforkey{\OtherN}{n_firms}{identification}{row_key}{other}
\DTLgetvalueforkey{\OtherMean}{mean}{identification}{row_key}{other}
\DTLgetvalueforkey{\OtherSD}{sd}{identification}{row_key}{other}
\DTLgetvalueforkey{\OtherPTF}{p25}{identification}{row_key}{other}
\DTLgetvalueforkey{\OtherPFZ}{median}{identification}{row_key}{other}
\DTLgetvalueforkey{\OtherPSF}{p75}{identification}{row_key}{other}
\DTLgetvalueforkey{\OtherNonzero}{share_nonzero_pct}{identification}{row_key}{other}

\DTLloaddb{pddid}{\Inputs/tables/pooled_did_results_table.csv}
\DTLgetvalueforkey{\IntlCoef}{coef_table}{pddid}{row_key}{international}
\DTLgetvalueforkey{\IntlSe}{se_table}{pddid}{row_key}{international}
\DTLgetvalueforkey{\IntlN}{n_firms}{pddid}{row_key}{international}
\DTLgetvalueforkey{\IntlObs}{n_obs}{pddid}{row_key}{international}
\DTLgetvalueforkey{\IntlStars}{stars}{pddid}{row_key}{international}
\DTLgetvalueforkey{\EtrCoef}{coef_table}{pddid}{row_key}{etr}
\DTLgetvalueforkey{\EtrSe}{se_table}{pddid}{row_key}{etr}
\DTLgetvalueforkey{\EtrN}{n_firms}{pddid}{row_key}{etr}
\DTLgetvalueforkey{\EtrObs}{n_obs}{pddid}{row_key}{etr}
\DTLgetvalueforkey{\EtrStars}{stars}{pddid}{row_key}{etr}
\DTLgetvalueforkey{\PiShrCoef}{coef_table}{pddid}{row_key}{foreign_profit_share}
\DTLgetvalueforkey{\PiShrSe}{se_table}{pddid}{row_key}{foreign_profit_share}
\DTLgetvalueforkey{\PiShrN}{n_firms}{pddid}{row_key}{foreign_profit_share}
\DTLgetvalueforkey{\PiShrObs}{n_obs}{pddid}{row_key}{foreign_profit_share}
\DTLgetvalueforkey{\IntlCoefAbs}{coef_abs}{pddid}{row_key}{international}
\DTLgetvalueforkey{\EtrCoefAbs}{coef_abs}{pddid}{row_key}{etr}
\DTLgetvalueforkey{\IntlCoefLongWindow}{coef_table}{pddid}{row_key}{international_long_window}
% Appendix reconciliations (A-1, A-2 of the Section 2 review).
\DTLloaddb{apprec}{\Inputs/tables/appendix_reconciliation.csv}
\DTLgetvalueforkey{\ArNoPacket}{value}{apprec}{row_key}{no_packet_canonical}
\DTLgetvalueforkey{\ArRescueOk}{value}{apprec}{row_key}{rescue_ok}
\DTLgetvalueforkey{\ArFailExtract}{value}{apprec}{row_key}{fail_at_extraction}
\DTLgetvalueforkey{\ArFailPanel}{value}{apprec}{row_key}{fail_at_panel}
\DTLgetvalueforkey{\ArPassThenFail}{value}{apprec}{row_key}{pass_then_fail}
\DTLgetvalueforkey{\ArFailThenPass}{value}{apprec}{row_key}{fail_then_pass}

% Real-activity outcomes in logs (Table \ref{tab:PooledDiDRA}).
\DTLloaddb{ralog}{\Inputs/tables/real_activity_log_table.csv}
\DTLgetvalueforkey{\RasCapx}{all_coef}{ralog}{row_key}{log_capx}
\DTLgetvalueforkey{\RasCapxSe}{all_se}{ralog}{row_key}{log_capx}
\DTLgetvalueforkey{\RasCapxStars}{all_stars}{ralog}{row_key}{log_capx}
\DTLgetvalueforkey{\RasCapxPft}{profitable_coef}{ralog}{row_key}{log_capx}
\DTLgetvalueforkey{\RasCapxPftSe}{profitable_se}{ralog}{row_key}{log_capx}
\DTLgetvalueforkey{\RasCapxPftStars}{profitable_stars}{ralog}{row_key}{log_capx}
\DTLgetvalueforkey{\RasCapxFirms}{all_firms}{ralog}{row_key}{log_capx}
\DTLgetvalueforkey{\RasCapxObs}{all_obs}{ralog}{row_key}{log_capx}
\DTLgetvalueforkey{\RasRd}{all_coef}{ralog}{row_key}{log_xrd}
\DTLgetvalueforkey{\RasRdSe}{all_se}{ralog}{row_key}{log_xrd}
\DTLgetvalueforkey{\RasRdStars}{all_stars}{ralog}{row_key}{log_xrd}
\DTLgetvalueforkey{\RasRdPft}{profitable_coef}{ralog}{row_key}{log_xrd}
\DTLgetvalueforkey{\RasRdPftSe}{profitable_se}{ralog}{row_key}{log_xrd}
\DTLgetvalueforkey{\RasRdPftStars}{profitable_stars}{ralog}{row_key}{log_xrd}
\DTLgetvalueforkey{\RasRdFirms}{all_firms}{ralog}{row_key}{log_xrd}
\DTLgetvalueforkey{\RasRdObs}{all_obs}{ralog}{row_key}{log_xrd}
\DTLgetvalueforkey{\RasSales}{all_coef}{ralog}{row_key}{log_sale}
\DTLgetvalueforkey{\RasSalesSe}{all_se}{ralog}{row_key}{log_sale}
\DTLgetvalueforkey{\RasSalesStars}{all_stars}{ralog}{row_key}{log_sale}
\DTLgetvalueforkey{\RasSalesPft}{profitable_coef}{ralog}{row_key}{log_sale}
\DTLgetvalueforkey{\RasSalesPftSe}{profitable_se}{ralog}{row_key}{log_sale}
\DTLgetvalueforkey{\RasSalesPftStars}{profitable_stars}{ralog}{row_key}{log_sale}
\DTLgetvalueforkey{\RasSalesFirms}{all_firms}{ralog}{row_key}{log_sale}
\DTLgetvalueforkey{\RasSalesObs}{all_obs}{ralog}{row_key}{log_sale}
\DTLgetvalueforkey{\RasEmp}{all_coef}{ralog}{row_key}{log_emp}
\DTLgetvalueforkey{\RasEmpSe}{all_se}{ralog}{row_key}{log_emp}
\DTLgetvalueforkey{\RasEmpStars}{all_stars}{ralog}{row_key}{log_emp}
\DTLgetvalueforkey{\RasEmpPft}{profitable_coef}{ralog}{row_key}{log_emp}
\DTLgetvalueforkey{\RasEmpPftSe}{profitable_se}{ralog}{row_key}{log_emp}
\DTLgetvalueforkey{\RasEmpPftStars}{profitable_stars}{ralog}{row_key}{log_emp}
\DTLgetvalueforkey{\RasEmpFirms}{all_firms}{ralog}{row_key}{log_emp}
\DTLgetvalueforkey{\RasEmpObs}{all_obs}{ralog}{row_key}{log_emp}
\DTLgetvalueforkey{\RasAt}{all_coef}{ralog}{row_key}{log_at}
\DTLgetvalueforkey{\RasAtSe}{all_se}{ralog}{row_key}{log_at}
\DTLgetvalueforkey{\RasAtStars}{all_stars}{ralog}{row_key}{log_at}
\DTLgetvalueforkey{\RasAtPft}{profitable_coef}{ralog}{row_key}{log_at}
\DTLgetvalueforkey{\RasAtPftSe}{profitable_se}{ralog}{row_key}{log_at}
\DTLgetvalueforkey{\RasAtPftStars}{profitable_stars}{ralog}{row_key}{log_at}
\DTLgetvalueforkey{\RasAtFirms}{all_firms}{ralog}{row_key}{log_at}
\DTLgetvalueforkey{\RasAtObs}{all_obs}{ralog}{row_key}{log_at}
\DTLgetvalueforkey{\RasPi}{all_coef}{ralog}{row_key}{log_pi}
\DTLgetvalueforkey{\RasPiSe}{all_se}{ralog}{row_key}{log_pi}
\DTLgetvalueforkey{\RasPiStars}{all_stars}{ralog}{row_key}{log_pi}
\DTLgetvalueforkey{\RasPiPft}{profitable_coef}{ralog}{row_key}{log_pi}
\DTLgetvalueforkey{\RasPiPftSe}{profitable_se}{ralog}{row_key}{log_pi}
\DTLgetvalueforkey{\RasPiPftStars}{profitable_stars}{ralog}{row_key}{log_pi}
\DTLgetvalueforkey{\RasPiFirms}{all_firms}{ralog}{row_key}{log_pi}
\DTLgetvalueforkey{\RasPiObs}{all_obs}{ralog}{row_key}{log_pi}
\DTLgetvalueforkey{\RasPiFo}{all_coef}{ralog}{row_key}{log_pifo}
\DTLgetvalueforkey{\RasPiFoSe}{all_se}{ralog}{row_key}{log_pifo}
\DTLgetvalueforkey{\RasPiFoStars}{all_stars}{ralog}{row_key}{log_pifo}
\DTLgetvalueforkey{\RasPiFoPft}{profitable_coef}{ralog}{row_key}{log_pifo}
\DTLgetvalueforkey{\RasPiFoPftSe}{profitable_se}{ralog}{row_key}{log_pifo}
\DTLgetvalueforkey{\RasPiFoPftStars}{profitable_stars}{ralog}{row_key}{log_pifo}
\DTLgetvalueforkey{\RasPiFoFirms}{all_firms}{ralog}{row_key}{log_pifo}
\DTLgetvalueforkey{\RasPiFoObs}{all_obs}{ralog}{row_key}{log_pifo}
\DTLgetvalueforkey{\RasPiDom}{all_coef}{ralog}{row_key}{log_pidom}
\DTLgetvalueforkey{\RasPiDomSe}{all_se}{ralog}{row_key}{log_pidom}
\DTLgetvalueforkey{\RasPiDomStars}{all_stars}{ralog}{row_key}{log_pidom}
\DTLgetvalueforkey{\RasPiDomPft}{profitable_coef}{ralog}{row_key}{log_pidom}
\DTLgetvalueforkey{\RasPiDomPftSe}{profitable_se}{ralog}{row_key}{log_pidom}
\DTLgetvalueforkey{\RasPiDomPftStars}{profitable_stars}{ralog}{row_key}{log_pidom}
\DTLgetvalueforkey{\RasPiDomFirms}{all_firms}{ralog}{row_key}{log_pidom}
\DTLgetvalueforkey{\RasPiDomObs}{all_obs}{ralog}{row_key}{log_pidom}
\DTLgetvalueforkey{\RasPiPftFirms}{profitable_firms}{ralog}{row_key}{log_pi}
\DTLgetvalueforkey{\RasPiPftObs}{profitable_obs}{ralog}{row_key}{log_pi}
\DTLgetvalueforkey{\RasPiFoPftFirms}{profitable_firms}{ralog}{row_key}{log_pifo}
\DTLgetvalueforkey{\RasPiFoPftObs}{profitable_obs}{ralog}{row_key}{log_pifo}
\DTLgetvalueforkey{\RasPiDomPftFirms}{profitable_firms}{ralog}{row_key}{log_pidom}
\DTLgetvalueforkey{\RasPiDomPftObs}{profitable_obs}{ralog}{row_key}{log_pidom}
% Balance-sheet margins behind the cash sentence in 3.2.3.
\DTLloaddb{bsheet}{\Inputs/tables/balance_sheet_table.csv}
\DTLgetvalueforkey{\BsCash}{coef}{bsheet}{row_key}{cash}
\DTLgetvalueforkey{\BsCashSe}{se}{bsheet}{row_key}{cash}
\DTLgetvalueforkey{\BsCashStars}{stars}{bsheet}{row_key}{cash}
\DTLgetvalueforkey{\BsTaxPaid}{coef}{bsheet}{row_key}{taxes_paid}
\DTLgetvalueforkey{\BsTaxPaidSe}{se}{bsheet}{row_key}{taxes_paid}
\DTLgetvalueforkey{\BsTaxPaidStars}{stars}{bsheet}{row_key}{taxes_paid}
\DTLgetvalueforkey{\BsDiv}{coef}{bsheet}{row_key}{dividends}
\DTLgetvalueforkey{\BsDivSe}{se}{bsheet}{row_key}{dividends}
\DTLgetvalueforkey{\BsDivStars}{stars}{bsheet}{row_key}{dividends}
\DTLgetvalueforkey{\BsBuyback}{coef}{bsheet}{row_key}{repurchases}
\DTLgetvalueforkey{\BsBuybackSe}{se}{bsheet}{row_key}{repurchases}
\DTLgetvalueforkey{\BsBuybackStars}{stars}{bsheet}{row_key}{repurchases}
\DTLgetvalueforkey{\BsDebt}{coef}{bsheet}{row_key}{debt}
\DTLgetvalueforkey{\BsDebtSe}{se}{bsheet}{row_key}{debt}
\DTLgetvalueforkey{\BsDebtStars}{stars}{bsheet}{row_key}{debt}
\DTLgetvalueforkey{\BsAcq}{coef}{bsheet}{row_key}{acquisitions}
\DTLgetvalueforkey{\BsAcqSe}{se}{bsheet}{row_key}{acquisitions}
\DTLgetvalueforkey{\BsAcqStars}{stars}{bsheet}{row_key}{acquisitions}
% Segment-based capital-location estimates quoted in 3.2.3.
\DTLloaddb{caploc}{\Inputs/tables/capital_location_table.csv}
\DTLgetvalueforkey{\CapLocPpe}{coef}{caploc}{row_key}{ppe_share}
\DTLgetvalueforkey{\CapLocPpeSe}{se}{caploc}{row_key}{ppe_share}
\DTLgetvalueforkey{\CapLocPpeFirms}{n_firms}{caploc}{row_key}{ppe_share}
\DTLgetvalueforkey{\CapLocPpeObs}{n_obs}{caploc}{row_key}{ppe_share}
\DTLgetvalueforkey{\CapLocCapx}{coef}{caploc}{row_key}{capx_share}
\DTLgetvalueforkey{\CapLocCapxSe}{se}{caploc}{row_key}{capx_share}
\DTLgetvalueforkey{\CapLocCapxFirms}{n_firms}{caploc}{row_key}{capx_share}
\DTLgetvalueforkey{\CapLocCapxObs}{n_obs}{caploc}{row_key}{capx_share}
% Windfall check (Section 3.2.4).
\DTLloaddb{wctrl}{\Inputs/tables/windfall_control_table.csv}
\DTLgetvalueforkey{\WcCorr}{value}{wctrl}{row_key}{corr_w_intl}
\DTLgetvalueforkey{\WcMeanExp}{value}{wctrl}{row_key}{mean_w_exposed}
\DTLgetvalueforkey{\WcMeanUnexp}{value}{wctrl}{row_key}{mean_w_unexposed}
\DTLgetvalueforkey{\WcShrBase}{value}{wctrl}{row_key}{foreign_share_base}
\DTLgetvalueforkey{\WcShrBaseSe}{value}{wctrl}{row_key}{foreign_share_base_se}
\DTLgetvalueforkey{\WcShrCtrl}{value}{wctrl}{row_key}{foreign_share_ctrl}
\DTLgetvalueforkey{\WcShrCtrlSe}{value}{wctrl}{row_key}{foreign_share_ctrl_se}
\DTLgetvalueforkey{\WcShrWCoef}{value}{wctrl}{row_key}{foreign_share_w_coef}
\DTLgetvalueforkey{\WcShrWSe}{value}{wctrl}{row_key}{foreign_share_w_se}
\DTLgetvalueforkey{\WcShrFirms}{value}{wctrl}{row_key}{foreign_share_firms}
% Case studies (Section 3.3).
\DTLloaddb{cases}{\Inputs/tables/case_study_table.csv}
\DTLloaddb{caselit}{\Inputs/tables/case_study_litigation.csv}
\DTLgetvalueforkey{\CaseAlphaPre}{intl_pre}{cases}{row_key}{alphabet}
\DTLgetvalueforkey{\CaseAlphaPost}{intl_post}{cases}{row_key}{alphabet}
\DTLgetvalueforkey{\CaseAlphaClose}{intl_closure}{cases}{row_key}{alphabet}
\DTLgetvalueforkey{\CaseAlphaShrPre}{share_pre_pct}{cases}{row_key}{alphabet}
\DTLgetvalueforkey{\CaseAlphaShrPost}{share_post_pct}{cases}{row_key}{alphabet}
\DTLgetvalueforkey{\CaseMsftPre}{intl_pre}{cases}{row_key}{microsoft}
\DTLgetvalueforkey{\CaseMsftPost}{intl_post}{cases}{row_key}{microsoft}
\DTLgetvalueforkey{\CaseMsftClose}{intl_closure}{cases}{row_key}{microsoft}
\DTLgetvalueforkey{\CaseMsftShrPre}{share_pre_pct}{cases}{row_key}{microsoft}
\DTLgetvalueforkey{\CaseMsftShrPost}{share_post_pct}{cases}{row_key}{microsoft}
\DTLgetvalueforkey{\CaseOnSemiPre}{intl_pre}{cases}{row_key}{on_semiconductor}
\DTLgetvalueforkey{\CaseOnSemiPost}{intl_post}{cases}{row_key}{on_semiconductor}
\DTLgetvalueforkey{\CaseOnSemiClose}{intl_closure}{cases}{row_key}{on_semiconductor}
\DTLgetvalueforkey{\CaseOnSemiShrPre}{share_pre_pct}{cases}{row_key}{on_semiconductor}
\DTLgetvalueforkey{\CaseOnSemiShrPost}{share_post_pct}{cases}{row_key}{on_semiconductor}
\DTLgetvalueforkey{\CaseQcomPre}{intl_pre}{cases}{row_key}{qualcomm}
\DTLgetvalueforkey{\CaseQcomPost}{intl_post}{cases}{row_key}{qualcomm}
\DTLgetvalueforkey{\CaseQcomClose}{intl_closure}{cases}{row_key}{qualcomm}
\DTLgetvalueforkey{\CaseQcomShrPre}{share_pre_pct}{cases}{row_key}{qualcomm}
\DTLgetvalueforkey{\CaseQcomShrPost}{share_post_pct}{cases}{row_key}{qualcomm}
\DTLgetvalueforkey{\CaseKoffPre}{intl_pre}{cases}{row_key}{coca_cola}
\DTLgetvalueforkey{\CaseKoffPost}{intl_post}{cases}{row_key}{coca_cola}
\DTLgetvalueforkey{\CaseKoffClose}{intl_closure}{cases}{row_key}{coca_cola}
\DTLgetvalueforkey{\CaseKoffShrPre}{share_pre_pct}{cases}{row_key}{coca_cola}
\DTLgetvalueforkey{\CaseKoffShrPost}{share_post_pct}{cases}{row_key}{coca_cola}
\DTLgetvalueforkey{\LitNLit}{value}{caselit}{row_key}{n_litigants_in_sample}
\DTLgetvalueforkey{\LitNFirms}{value}{caselit}{row_key}{n_sample_firms}
\DTLgetvalueforkey{\LitIntlFull}{value}{caselit}{row_key}{intl_full}
\DTLgetvalueforkey{\LitIntlFullSe}{value}{caselit}{row_key}{intl_full_se}
\DTLgetvalueforkey{\LitIntlDrop}{value}{caselit}{row_key}{intl_drop}
\DTLgetvalueforkey{\LitIntlDropSe}{value}{caselit}{row_key}{intl_drop_se}
\DTLgetvalueforkey{\LitEtrFull}{value}{caselit}{row_key}{etr_full}
\DTLgetvalueforkey{\LitEtrFullSe}{value}{caselit}{row_key}{etr_full_se}
\DTLgetvalueforkey{\LitEtrDrop}{value}{caselit}{row_key}{etr_drop}
\DTLgetvalueforkey{\LitEtrDropSe}{value}{caselit}{row_key}{etr_drop_se}
\DTLgetvalueforkey{\LitShrFull}{value}{caselit}{row_key}{share_full}
\DTLgetvalueforkey{\LitShrFullSe}{value}{caselit}{row_key}{share_full_se}
\DTLgetvalueforkey{\LitShrDrop}{value}{caselit}{row_key}{share_drop}
\DTLgetvalueforkey{\LitShrDropSe}{value}{caselit}{row_key}{share_drop_se}

% Profit-shifting heterogeneity (Table \ref{tab:Heterogeneity}).
\DTLloaddb{het}{\Inputs/tables/heterogeneity_gamma.csv}
\DTLloaddb{hethead}{\Inputs/tables/heterogeneity_headlines.csv}
\DTLgetvalueforkey{\HetRd}{gamma_pp}{het}{row_key}{pre_rd_at_z}
\DTLgetvalueforkey{\HetRdSe}{gamma_se_pp}{het}{row_key}{pre_rd_at_z}
\DTLgetvalueforkey{\HetRdStars}{stars}{het}{row_key}{pre_rd_at_z}
\DTLgetvalueforkey{\HetRdObs}{n_obs}{het}{row_key}{pre_rd_at_z}
\DTLgetvalueforkey{\HetCash}{gamma_pp}{het}{row_key}{pre_cash_at_z}
\DTLgetvalueforkey{\HetCashSe}{gamma_se_pp}{het}{row_key}{pre_cash_at_z}
\DTLgetvalueforkey{\HetCashStars}{stars}{het}{row_key}{pre_cash_at_z}
\DTLgetvalueforkey{\HetCashObs}{n_obs}{het}{row_key}{pre_cash_at_z}
\DTLgetvalueforkey{\HetIntan}{gamma_pp}{het}{row_key}{pre_intan_ex_gw_at_z}
\DTLgetvalueforkey{\HetIntanSe}{gamma_se_pp}{het}{row_key}{pre_intan_ex_gw_at_z}
\DTLgetvalueforkey{\HetIntanStars}{stars}{het}{row_key}{pre_intan_ex_gw_at_z}
\DTLgetvalueforkey{\HetIntanObs}{n_obs}{het}{row_key}{pre_intan_ex_gw_at_z}
\DTLgetvalueforkey{\HetFgnRev}{gamma_pp}{het}{row_key}{pre_foreign_rev_share_z}
\DTLgetvalueforkey{\HetFgnRevSe}{gamma_se_pp}{het}{row_key}{pre_foreign_rev_share_z}
\DTLgetvalueforkey{\HetFgnRevStars}{stars}{het}{row_key}{pre_foreign_rev_share_z}
\DTLgetvalueforkey{\HetFgnRevObs}{n_obs}{het}{row_key}{pre_foreign_rev_share_z}
\DTLgetvalueforkey{\HetHaven}{gamma_pp}{het}{row_key}{pre_has_haven_or_conduit}
\DTLgetvalueforkey{\HetHavenSe}{gamma_se_pp}{het}{row_key}{pre_has_haven_or_conduit}
\DTLgetvalueforkey{\HetHavenStars}{stars}{het}{row_key}{pre_has_haven_or_conduit}
\DTLgetvalueforkey{\HetHavenObs}{n_obs}{het}{row_key}{pre_has_haven_or_conduit}
\DTLgetvalueforkey{\HetClassicHaven}{gamma_pp}{het}{row_key}{pre_has_classic_haven}
\DTLgetvalueforkey{\HetClassicHavenSe}{gamma_se_pp}{het}{row_key}{pre_has_classic_haven}
\DTLgetvalueforkey{\HetClassicHavenStars}{stars}{het}{row_key}{pre_has_classic_haven}
\DTLgetvalueforkey{\HetClassicHavenObs}{n_obs}{het}{row_key}{pre_has_classic_haven}
\DTLgetvalueforkey{\HetHavenCount}{gamma_pp}{het}{row_key}{pre_log1p_haven_subs_z}
\DTLgetvalueforkey{\HetHavenCountSe}{gamma_se_pp}{het}{row_key}{pre_log1p_haven_subs_z}
\DTLgetvalueforkey{\HetHavenCountStars}{stars}{het}{row_key}{pre_log1p_haven_subs_z}
\DTLgetvalueforkey{\HetHavenCountObs}{n_obs}{het}{row_key}{pre_log1p_haven_subs_z}
\DTLgetvalueforkey{\HetLogAt}{gamma_pp}{het}{row_key}{pre_log_at_z}
\DTLgetvalueforkey{\HetLogAtSe}{gamma_se_pp}{het}{row_key}{pre_log_at_z}
\DTLgetvalueforkey{\HetLogAtStars}{stars}{het}{row_key}{pre_log_at_z}
\DTLgetvalueforkey{\HetLogAtObs}{n_obs}{het}{row_key}{pre_log_at_z}
\DTLgetvalueforkey{\HetExpExposure}{value}{hethead}{row_key}{typicalExposedExposurePp}
\DTLgetvalueforkey{\HetExpShare}{value}{hethead}{row_key}{typicalExposedSharePct}
\DTLgetvalueforkey{\HetExpDrop}{value}{hethead}{row_key}{typicalExposedDropPp}
\DTLgetvalueforkey{\HetExpDropPct}{value}{hethead}{row_key}{typicalExposedDropPct}
\DTLgetvalueforkey{\HetPiDom}{value}{hethead}{row_key}{betaPiDom}
\DTLgetvalueforkey{\HetPiDomSe}{value}{hethead}{row_key}{betaPiDomSe}
\DTLgetvalueforkey{\HetPiFo}{value}{hethead}{row_key}{betaPiFo}
\DTLgetvalueforkey{\HetPiFoSe}{value}{hethead}{row_key}{betaPiFoSe}

% Per-channel pooled DiD (Table \ref{tab:ChannelDiD}) and the identity check.
\DTLloaddb{chandid}{\Inputs/tables/channel_did_table.csv}
\DTLgetvalueforkey{\ChIntlCoef}{coef}{chandid}{row_key}{international}
\DTLgetvalueforkey{\ChIntlSe}{se}{chandid}{row_key}{international}
\DTLgetvalueforkey{\ChIntlStars}{stars}{chandid}{row_key}{international}
\DTLgetvalueforkey{\ChStateLocalCoef}{coef}{chandid}{row_key}{state_local}
\DTLgetvalueforkey{\ChStateLocalSe}{se}{chandid}{row_key}{state_local}
\DTLgetvalueforkey{\ChStateLocalStars}{stars}{chandid}{row_key}{state_local}
\DTLgetvalueforkey{\ChFedCredCoef}{coef}{chandid}{row_key}{federal_credits}
\DTLgetvalueforkey{\ChFedCredSe}{se}{chandid}{row_key}{federal_credits}
\DTLgetvalueforkey{\ChFedCredStars}{stars}{chandid}{row_key}{federal_credits}
\DTLgetvalueforkey{\ChBookTaxCoef}{coef}{chandid}{row_key}{permanent_items}
\DTLgetvalueforkey{\ChBookTaxSe}{se}{chandid}{row_key}{permanent_items}
\DTLgetvalueforkey{\ChBookTaxStars}{stars}{chandid}{row_key}{permanent_items}
\DTLgetvalueforkey{\ChValnCoef}{coef}{chandid}{row_key}{reserves_losses}
\DTLgetvalueforkey{\ChValnSe}{se}{chandid}{row_key}{reserves_losses}
\DTLgetvalueforkey{\ChValnStars}{stars}{chandid}{row_key}{reserves_losses}
\DTLgetvalueforkey{\ChOneTimeCoef}{coef}{chandid}{row_key}{tcja_onetime}
\DTLgetvalueforkey{\ChOneTimeSe}{se}{chandid}{row_key}{tcja_onetime}
\DTLgetvalueforkey{\ChOneTimeStars}{stars}{chandid}{row_key}{tcja_onetime}
\DTLgetvalueforkey{\ChResidCoef}{coef}{chandid}{row_key}{other}
\DTLgetvalueforkey{\ChResidSe}{se}{chandid}{row_key}{other}
\DTLgetvalueforkey{\ChResidStars}{stars}{chandid}{row_key}{other}
\DTLgetvalueforkey{\ChSumCoef}{coef}{chandid}{row_key}{sum_channels}
\DTLgetvalueforkey{\ChNonIntlCoef}{coef}{chandid}{row_key}{sum_non_international}
\DTLgetvalueforkey{\ChOffsetPctCoef}{coef}{chandid}{row_key}{offset_share_pct}
\DTLgetvalueforkey{\ChObs}{n_obs}{chandid}{row_key}{international}
\DTLgetvalueforkey{\ChFirms}{n_firms}{chandid}{row_key}{international}

% Figure 3 estimates eq. (2), whose sample keeps the 2017-2018 transition
% years the pooled DiD drops, so its N is larger than Table 7's.
\DTLloaddb{esintl}{\Inputs/first_stage_tikz/first_stage_event_summary.csv}
\DTLloaddb{esetr}{\Inputs/first_etr_tikz/first_etr_event_summary.csv}
\DTLloaddb{espishr}{\Inputs/first_pishr_tikz/first_pishr_event_summary.csv}
\DTLgetvalueforkey{\IntlEventObs}{event_obs}{esintl}{pooled_stars}{***}
\DTLgetvalueforkey{\EtrEventObs}{event_obs}{esetr}{pooled_stars}{***}
\DTLgetvalueforkey{\PiShrEventObs}{event_obs}{espishr}{pooled_stars}{***}
\DTLgetvalueforkey{\PiShrStars}{stars}{pddid}{row_key}{foreign_profit_share}

%% Stage 3 category-assignment diagnostics
\DTLloaddb{catdiag}{\Inputs/tables/category_diagnostics.csv}
\DTLgetvalueforkey{\CatDistinctLabels}{value}{catdiag}{key}{catDistinctLabels}
\DTLgetvalueforkey{\CatNumRules}{value}{catdiag}{key}{catNumRules}
\DTLgetvalueforkey{\CatNumCategories}{value}{catdiag}{key}{catNumCategories}
\DTLgetvalueforkey{\CatNoIntlFyPct}{value}{catdiag}{key}{catNoIntlFyPct}
\DTLgetvalueforkey{\CatOtherEqualPct}{value}{catdiag}{key}{catOtherEqualPct}
\DTLgetvalueforkey{\CatOtherSalesPct}{value}{catdiag}{key}{catOtherSalesPct}
\DTLgetvalueforkey{\CatIntlSalesPct}{value}{catdiag}{key}{catIntlSalesPct}
\DTLgetvalueforkey{\CatOtherForeignRowPct}{value}{catdiag}{key}{catOtherForeignRowPct}
\DTLgetvalueforkey{\CatOtherForeignPpPct}{value}{catdiag}{key}{catOtherForeignPpPct}
\DTLgetvalueforkey{\EtrAgreeHalfPct}{value}{catdiag}{key}{etrAgreeHalfPct}
\DTLgetvalueforkey{\EtrAgreeMedianAbs}{value}{catdiag}{key}{etrAgreeMedianAbs}
\DTLgetvalueforkey{\EtrAgreeN}{value}{catdiag}{key}{etrAgreeN}
\DTLgetvalueforkey{\PanelPriorYearN}{value}{catdiag}{key}{panelPriorYearN}
\DTLgetvalueforkey{\PanelPriorYearPct}{value}{catdiag}{key}{panelPriorYearPct}

%% Section 4 preamble
\DTLloaddb{theorycalib}{\Inputs/tables/theory_calibration.csv}
\DTLgetvalueforkey{\BetaIntlSame}{value}{theorycalib}{key}{betaIntlSame}
\DTLgetvalueforkey{\BetaIntlSameSe}{value}{theorycalib}{key}{betaIntlSameSe}
\DTLgetvalueforkey{\BetaIntlSameG}{value}{theorycalib}{key}{betaIntlSameG}
\DTLgetvalueforkey{\BetaIntlSameN}{value}{theorycalib}{key}{betaIntlSameN}
\DTLgetvalueforkey{\BetaIntlFullAbs}{value}{theorycalib}{key}{betaIntlFullAbs}
\DTLgetvalueforkey{\GTauHPostStatZero}{value}{theorycalib}{key}{gTauHPostStatZero}
\DTLgetvalueforkey{\GTauHPostStatFive}{value}{theorycalib}{key}{gTauHPostStatFive}
\DTLgetvalueforkey{\GTauHPostStatTen}{value}{theorycalib}{key}{gTauHPostStatTen}
\DTLgetvalueforkey{\GTauFNoGiltiZero}{value}{theorycalib}{key}{gTauFNoGiltiZero}
\DTLgetvalueforkey{\GTauFNoGiltiFive}{value}{theorycalib}{key}{gTauFNoGiltiFive}
\DTLgetvalueforkey{\GTauFNoGiltiTen}{value}{theorycalib}{key}{gTauFNoGiltiTen}
\DTLgetvalueforkey{\GTauFStatZero}{value}{theorycalib}{key}{gTauFStatZero}
\DTLgetvalueforkey{\GTauFStatFive}{value}{theorycalib}{key}{gTauFStatFive}
\DTLgetvalueforkey{\GTauFStatTen}{value}{theorycalib}{key}{gTauFStatTen}
\DTLgetvalueforkey{\GPsiNoGiltiZero}{value}{theorycalib}{key}{gPsiNoGiltiZero}
\DTLgetvalueforkey{\GPsiNoGiltiFive}{value}{theorycalib}{key}{gPsiNoGiltiFive}
\DTLgetvalueforkey{\GPsiNoGiltiTen}{value}{theorycalib}{key}{gPsiNoGiltiTen}
\DTLgetvalueforkey{\GPsiStatZero}{value}{theorycalib}{key}{gPsiStatZero}
\DTLgetvalueforkey{\GPsiStatFive}{value}{theorycalib}{key}{gPsiStatFive}
\DTLgetvalueforkey{\GPsiStatTen}{value}{theorycalib}{key}{gPsiStatTen}
\DTLgetvalueforkey{\GShiftedPreNoGiltiZero}{value}{theorycalib}{key}{gShiftedPreNoGiltiZero}
\DTLgetvalueforkey{\GShiftedPreNoGiltiFive}{value}{theorycalib}{key}{gShiftedPreNoGiltiFive}
\DTLgetvalueforkey{\GShiftedPreNoGiltiTen}{value}{theorycalib}{key}{gShiftedPreNoGiltiTen}
\DTLgetvalueforkey{\GShiftedPreStatZero}{value}{theorycalib}{key}{gShiftedPreStatZero}
\DTLgetvalueforkey{\GShiftedPreStatFive}{value}{theorycalib}{key}{gShiftedPreStatFive}
\DTLgetvalueforkey{\GShiftedPreStatTen}{value}{theorycalib}{key}{gShiftedPreStatTen}
\DTLgetvalueforkey{\GShiftedPostNoGiltiZero}{value}{theorycalib}{key}{gShiftedPostNoGiltiZero}
\DTLgetvalueforkey{\GShiftedPostNoGiltiFive}{value}{theorycalib}{key}{gShiftedPostNoGiltiFive}
\DTLgetvalueforkey{\GShiftedPostNoGiltiTen}{value}{theorycalib}{key}{gShiftedPostNoGiltiTen}
\DTLgetvalueforkey{\GShiftedPostStatZero}{value}{theorycalib}{key}{gShiftedPostStatZero}
\DTLgetvalueforkey{\GShiftedPostStatFive}{value}{theorycalib}{key}{gShiftedPostStatFive}
\DTLgetvalueforkey{\GShiftedPostStatTen}{value}{theorycalib}{key}{gShiftedPostStatTen}
\DTLgetvalueforkey{\BetaIntlNegSame}{value}{theorycalib}{key}{betaIntlNegSame}
\DTLgetvalueforkey{\BetaShareNeg}{value}{theorycalib}{key}{betaShareNeg}
\DTLgetvalueforkey{\NegDropFirmsSamePct}{value}{theorycalib}{key}{negDropFirmsSamePct}
\DTLgetvalueforkey{\PsiCiLoStatPct}{value}{theorycalib}{key}{psiCiLoStatPct}
\DTLgetvalueforkey{\PsiCiHiStatPct}{value}{theorycalib}{key}{psiCiHiStatPct}
\DTLgetvalueforkey{\PsiCiLoNoGiltiPct}{value}{theorycalib}{key}{psiCiLoNoGiltiPct}
\DTLgetvalueforkey{\PsiCiHiNoGiltiPct}{value}{theorycalib}{key}{psiCiHiNoGiltiPct}
\DTLgetvalueforkey{\TauFCiLoStatPct}{value}{theorycalib}{key}{tauFCiLoStatPct}
\DTLgetvalueforkey{\TauFCiHiStatPct}{value}{theorycalib}{key}{tauFCiHiStatPct}
\DTLgetvalueforkey{\TauFCiLoNoGiltiPct}{value}{theorycalib}{key}{tauFCiLoNoGiltiPct}
\DTLgetvalueforkey{\TauFCiHiNoGiltiPct}{value}{theorycalib}{key}{tauFCiHiNoGiltiPct}
\DTLgetvalueforkey{\PsiInvFdiiStatPct}{value}{theorycalib}{key}{psiInvFdiiStatPct}
\DTLgetvalueforkey{\PsiInvFdiiNoGiltiPct}{value}{theorycalib}{key}{psiInvFdiiNoGiltiPct}
\DTLgetvalueforkey{\TauFImpliedFdiiStatPct}{value}{theorycalib}{key}{tauFImpliedFdiiStatPct}
\DTLgetvalueforkey{\RentRebookedStatPct}{value}{theorycalib}{key}{rentRebookedStatPct}
\DTLgetvalueforkey{\RentRebookedNoGiltiPct}{value}{theorycalib}{key}{rentRebookedNoGiltiPct}

\DTLloaddb{indpass}{\Inputs/tables/industry_passthrough_examples.csv}
\DTLgetvalueforkey{\TechIntlPre}{international_pre_pp}{indpass}{row_key}{tech}
\DTLgetvalueforkey{\TechIntlPost}{international_post_pp}{indpass}{row_key}{tech}
\DTLgetvalueforkey{\TechIntlClosure}{international_closure_pp}{indpass}{row_key}{tech}
\DTLgetvalueforkey{\TechGapClosure}{total_gap_closure_pp}{indpass}{row_key}{tech}
\DTLgetvalueforkey{\TechImpliedEtr}{implied_etr_increase_pp}{indpass}{row_key}{tech}
\DTLgetvalueforkey{\TechOtherChannels}{other_channels_abs_pp}{indpass}{row_key}{tech}
\DTLgetvalueforkey{\TechTotalGapClosure}{total_gap_closure_pp}{indpass}{row_key}{tech}

\DTLloaddb{profloc}{\Inputs/tables/profit_location_headlines.csv}
\DTLgetvalueforkey{\IncomeLocBetaForeignShare}{value}{profloc}{row_key}{incomeLocBetaForeignShare}
\DTLgetvalueforkey{\IncomeLocBetaForeignSharePp}{value}{profloc}{row_key}{incomeLocBetaForeignSharePp}
\DTLgetvalueforkey{\IncomeLocBetaForeignShareSe}{value}{profloc}{row_key}{incomeLocBetaForeignShareSe}
\DTLgetvalueforkey{\IncomeLocBetaForeignShareSePp}{value}{profloc}{row_key}{incomeLocBetaForeignShareSePp}
\DTLgetvalueforkey{\IncomeLocBetaEtrSubsample}{value}{profloc}{row_key}{incomeLocBetaEtrSubsample}
\DTLgetvalueforkey{\IncomeLocClosureRatio}{value}{profloc}{row_key}{incomeLocClosureRatio}
\DTLgetvalueforkey{\ForeignShareDeclinePp}{value}{profloc}{row_key}{foreignShareDeclinePp}
\DTLgetvalueforkey{\GeoRevenueShareFirms}{value}{profloc}{row_key}{geoRevenueShareFirms}
\DTLgetvalueforkey{\GeoRevenueShareFirmYears}{value}{profloc}{row_key}{geoRevenueShareFirmYears}
\DTLgetvalueforkey{\GeoRevenueBetaForeignSharePp}{value}{profloc}{row_key}{geoRevenueBetaForeignSharePp}
\DTLgetvalueforkey{\GeoRevenueBetaForeignShareSePp}{value}{profloc}{row_key}{geoRevenueBetaForeignShareSePp}


\DTLloaddb{windfallni}{\Inputs/tables/windfall_ni_headline.csv}
\DTLgetvalueforkey{\ThetaShock}{coef}{windfallni}{row_key}{ni_shock_2017}
\DTLgetvalueforkey{\ThetaShockLo}{lower95}{windfallni}{row_key}{ni_shock_2017}
\DTLgetvalueforkey{\ThetaShockHi}{upper95}{windfallni}{row_key}{ni_shock_2017}
\DTLgetvalueforkey{\ThetaPost}{coef}{windfallni}{row_key}{ni_post_2018}
\DTLgetvalueforkey{\ThetaPostSe}{se}{windfallni}{row_key}{ni_post_2018}


\title{Foreign Tax Advantages after the U.S.~Tax Reform of 2017}
\author{Anmol Bhandari\thanks{University of Minnesota} 
  \and Ellen McGrattan\thanks{University of Minnesota and Federal Reserve Bank of Minneapolis}}

\begin{document}
\maketitle
\pagenumbering{gobble}
\thispagestyle{empty} 

\vskip 2.4truein
\begin{abstract}
\singlespacing
\noindent
We construct a firm-year panel of income-tax rate reconciliations from
annual 10-K filings of public companies to study the effects of the Tax Cuts and Jobs 
Act (TCJA) on U.S.~multinational profit shifting. These reconciliations 
decompose the gap between company-level effective tax rates and the federal 
statutory corporate tax rate, allowing us to isolate the contribution of
foreign tax advantages before and after the reform. We find that TCJA 
significantly reduced the tax rate gap and increased effective tax rates
for companies most exposed to pre-TCJA international tax benefits.
This resulted in an onshoring of profits with foreign profit shares falling
\ShareNum{\IncomeLocClosureRatio}\ 
percentage points per percentage point of tax increase.
Among more exposed companies, we find no differential change in 
investment, employment, or the geographic distribution of sales and installed capital.
A neoclassical model of a multinational that chooses where to book its
income rationalizes this pattern: the reform raised the average tax on
income booked abroad relative to income booked at home, which governs
where profits are reported, while the marginal tax on capital, which
governs investment, changed by an amount common to all firms.
The tax reform also changed company balance sheets through 
a one-time remeasurement of deferred taxes.
Companies with larger remeasurement gains did not invest more,
and controlling for these gains leaves our main results unchanged.\footnote{%
We thank Vishruth Parikh for valuable research assistance.  This research 
has been supported by National Science Foundation grant SES-2214248 and the 
Heller-Hurwicz Economics Institute.}
\end{abstract}

\newpage
%\setcounter{page}{1}
\pagenumbering{arabic}

\section{Introduction}

One of the central policy goals of the Tax Cuts and Jobs Act (TCJA) of 2017 was to
reduce incentives for profit shifting by U.S.~multinationals and to 
encourage the onshoring of profits in the United States. Toward this goal, 
the federal corporate tax rate was lowered from 35
to 21 percent, special deductions for foreign-derived income were created,
tax floors for foreign intangible income and deductible payments to foreign
affiliates were introduced,
and
indefinite tax deferrals abroad were disallowed.
This paper provides a post-TCJA policy evaluation using a novel dataset
from Securities and Exchange Commission (SEC) filings of company annual reports.
We combine automated rule-based parsing with large language models (LLMs)
to extract the full income tax
reconciliation table that is reported on SEC 10-K forms. The table reconciles
a company's effective tax rate and the federal statutory corporate income tax rate by listing
all adjustments for items such as foreign and local tax differentials, tax credits,
one-time changes in law, and valuation adjustments.\footnote{%
The Compustat North America Fundamentals Annual database (hereafter \citet{CompustatAnnual})
reports the effective tax rate,
but does not provide details on reconciliation adjustments.}
This reconciliation in turn is used along with a difference-in-differences 
strategy to estimate the causal impacts of TCJA and the
extent to which the tax reform achieved the stated objective of onshoring
profits. 

The empirical strategy exploits cross-sectional variation in pre-TCJA exposure to foreign 
tax benefits and compares company post-TCJA outcomes, most notably, effective tax rates, 
foreign profit shares, and real productive activities.
To measure foreign exposure, we use adjustment items in the income tax reconciliation
table related to the taxation of foreign corporate profits. In the pre-TCJA period,
companies reported foreign and domestic corporate income tax rate differentials.
In the post-TCJA period, companies also included adjustments related to the international 
provisions of the tax reform act such as the tax floor on foreign intangible income.
We aggregate all reconciliation items reflecting differences in the
taxation of foreign and domestic income
and use this international contribution to the overall tax rate
gap as our main explanatory variable in the difference-in-differences regression.
We also include firm and industry-year fixed effects and perform preliminary 
validity checks to verify that there is meaningful within-industry variation
in our measure of international exposure and parallel pre-TCJA trends.  

The main data input for the empirical analysis is a reproducible firm-year panel of 
income tax reconciliation tables 
over 1998--2024, linked to the \citet{CompustatAnnual} database. 
For every 10-K filing of a Compustat firm available from the SEC Electronic Data Gathering, Analysis, and
Retrieval System database (hereafter \citet{EDGAR}), an automated parser locates 
the reconciliation table and a second stage parser---that is rule-based if the table is regular,
and a large language model if it is not---converts it to 
a common row format. Every extracted table must satisfy the accounting identity 
that the statutory rate plus the reported adjustments equals the reported 
effective rate. The panel covers \ShareNum{\FullSampleSales}\ percent of 
Compustat sales for our main sample period 2013--2024.

Once we have the firm-year panel, we group reconciliation adjustments into 
economically meaningful categories.  There are two main categories related to TCJA. 
The first category, which was described earlier, 
includes all adjustments related to differences in 
taxation of profits booked
in the United States and profits booked abroad. 
The second category includes all adjustments related to current-year changes in law such
as those indicating a change under TCJA in 2017. 
We document a significant closure in the average sales-weighted gap between 
the effective and statutory tax rates after the 2017 reform.
We find that the largest changes are attributable to the international and one-time adjustments,
with very little change in the remaining reconciliation categories---namely, 
federal credits, state and local taxes, and general accounting adjustments. 
We find similar patterns across industries, with the largest changes
in the international tax gap occurring in technology, 
pharmaceuticals, manufacturing, and oil and gas. 

Using our difference-in-differences framework, we estimate that one percentage 
point more pre-TCJA exposure is associated with a 
\AbsRatioNum{\IntlCoef}\ percentage points 
closure in the international contribution to the tax rate gap.
Furthermore, we find there was a significant pass-through to higher 
effective tax rates:
one percentage point more pre-TCJA exposure is associated with a
\AbsRatioNum{\EtrCoef}\ percentage points increase in the rate.
To provide context, 
consider two companies in the same industry whose international contributions
to the tax rate gap differed by 
\RoundNum{\WorkedExampleGapPp}\ 
percentage points before the reform---say,
$-\RoundNum{\WorkedExampleFromPp}$ and $-\RoundNum{\WorkedExampleToPp}$.
After TCJA, the more exposed company's international contribution closed by 
\ShareNum{\WorkedExampleIntlPp}\ 
percentage points more than the less exposed
company's, and its effective tax rate rose by 
\ShareNum{\WorkedExampleEtrPp}\ 
percentage points more.
The difference between the two is what companies recovered through the
remaining reconciliation categories.

The same empirical strategy is used to estimate the extent of post-TCJA
onshoring and changes in real activities.
We find that TCJA significantly 
lowered reported foreign profit shares of the internationally
exposed companies, but that more exposed companies did not
differentially change their productive activities---neither the
location nor the scale. 
For every one percentage point of TCJA-induced effective tax rate increase,
the company's foreign profit share falls by 
\RatioNum{\IncomeLocClosureRatio}\ percentage points.
Because the reform's changes to the user cost
of capital were common to all companies in an industry, our design is
silent on whether investment rose in the aggregate.
We find that companies most responsive in relocating reported
profits are R\&D intensive with 
large cash balances relative to total assets, but we find no difference in 
income shifting for firms with and without subsidiaries in tax havens.
While the companies shifted profits back to the United
States, the geography of their sales and installed capital hardly changed.
Similarly, levels of capital and R\&D expenditures and growth in 
sales and employment show no significant response to the
reduction in foreign tax benefits.

Separately from the international tax benefits, the domestic corporate
tax cut of 14 percentage points under TCJA generated a one-time 
deferred-tax windfall that could relax financial constraints for
many companies and spur investment---potentially confounding our 
baseline empirical results.
We investigate whether this one-time gain
translated into higher investment by exploiting differences in 
net deferred tax liabilities across company balance sheets. 
We run a difference-in-differences regression with the main
right-hand-side variable being the windfall gain that is equal 
to the net deferred tax balance---deferred tax assets less 
liabilities---in 2016 divided by the total assets in 
2015 and multiplied by the tax change. We find
no detectable effect on capital expenditures or R\&D expenses
in the aggregate or across different subgroups of firms. 
As a robustness check on our baseline results, we also include
the windfall variable as a control and find the main results
are unchanged.

To interpret our empirical findings, we propose a minimal neoclassical model
of a multinational that chooses R\&D, tangible investment at home and abroad,
and where to book its income. In the model,
investment depends on marginal user costs of capital, and
booking depends on the wedge between the domestic and foreign tax
rates on booked income. TCJA changed user costs by amounts common
within an industry, which our fixed effects absorb, and collapsed the tax
wedge---most for the firms that had booked the most income abroad,
which is what the international exposure measure captures
in our empirical design. We use our difference-in-differences
regression coefficients and key relations in the model to estimate
the share of the foreign tax advantage---as measured by the
international contribution in the income tax reconciliation 
table---attributable to pure profit-shifting to tax havens
with no substantive productive operations. We find shares
ranging from \RoundNum{\GPsiStatTen} to \RoundNum{\GPsiNoGiltiZero}
percent, depending on the effectiveness of the \emph{global
intangible low-taxed income} (GILTI) provisions. The
implied foreign tax rates on income that was not shifted are in the
range of \RoundNum{\GTauFStatZero} to
\RoundNum{\GTauFNoGiltiZero} percent.  


\subsection{Related Literature}

The income tax reconciliation table has been studied before, but collected by hand 
and used to answer different questions. \citet{DyrengHanlonMaydewThornock2017} 
document a quarter-century decline in the effective tax rates of U.S.~companies. 
\citet{DrakeHamiltonLusch2020} hand-collected reconciliation items 
over the period 1996--2015 and use it to study the role of valuation allowances in
explaining differences in effective tax rates over time and across companies.
\citet{SchwabStombergXia2022} use the reconciliation table to separately 
study factors related and unrelated to tax avoidance over the period 2008--2016.
We automate the extraction 
over every 10-K filer on EDGAR from 1998 to 2024 and link the result to 
Compustat, which is what makes it possible to follow the international 
contribution to the tax gap across the 2017 reform for the full 
population of public companies rather than a hand-collected subset. 
\citet{DonohoeMcGillOutslay2012} catalogue what these SEC disclosures 
cannot reveal about a multinational's foreign operations. The
classification diagnostics in Appendix~\ref{app:categories} 
are our accounting of those limits.

The main focus of our paper is multinational profit shifting. 
Since \citet{HinesRice1994}, who regressed affiliates' reported 
profits on tax-rate differences across countries, the responsiveness 
of booked income to tax differentials has been estimated many ways. 
\citet{Dharmapala2014} surveys this literature. \citet{DowdLandefeldMoore2017} 
use U.S.~tax returns to show that profit-shifting responses are highly
nonlinear, which is why we distinguish rates in tax havens from the average
rates on foreign operations when interpreting the data through the lens of
neoclassical theory (in Section~\ref{sec:Theory}).
At the aggregate level, \citet{TorslovWierZucman2023} estimate that 
roughly a third of global multinational profits are shifted to tax havens, 
with U.S.~companies shifting more than half of their profits and 
multinationals from other nations shifting about a quarter.
Their estimates are the natural comparison for the pre-TCJA profit shifted 
shares we recover later in Section~\ref{sec:Theory}.
\citet{GuvenenMataloniRassierRuhl2022} show that this profit shifting 
distorts measured U.S.~output, productivity, and the labor share. 

Closest to our analysis are the studies by \citet{Grubert2003,Grubert2012}, who used
company-level Internal Revenue Service (IRS) microdata from Forms 1120 and 5471.
\citet{Grubert2003} studied R\&D intensive controlled foreign corporations 
and estimated that half of observed differences in profitability between high 
and low-tax countries was due to shifting income derived from R\&D operations.
\citet{Grubert2012} estimates that a one percentage point increase in foreign effective 
tax rates reduces the foreign share of multinational income by roughly 0.7 percentage 
points.  His empirical strategy relies on variation across firms in the
change in their effective foreign tax rates between 1996 and 2004,
whereas we exploit TCJA as a 
large policy shock and compare firms with different pre-reform exposure to 
foreign tax advantages. This empirical design provides a more transparent
policy-induced source 
of identifying variation for estimating the effect of tax incentives on the 
location of reported profits.  In terms of the estimate of interest,
we find a larger response of reported foreign profits to tax incentives than in 
Grubert's pre-TCJA sample, most likely because TCJA generated a larger 
and more salient shock to tax incentives than the cross-country variation exploited
in the earlier work. With our dataset that includes information on capital expenditures,
research and development, and growth in sales and employment, we also
extend the analysis of \citet{Grubert2012} who did not have such measures
of real activity.

The literature evaluating TCJA is already large (see, for example, 
\citet{Dharmapala2024} and \citet{GravelleMarples2025} for surveys).  
Two findings in that literature frame ours. First, on aggregate profit shifting, 
\citet{GarciaBernardoJanskyZucman2026} combine survey, tax, and financial-statement 
data and find that the share of U.S.~multinationals' profits booked abroad fell by 
only one to five percentage points after the reform and that the haven share 
of foreign profits was roughly stable.
\citet{Dharmapala2024} concludes that shifting fell by no more than the rate 
cut alone would predict. Our design is complementary rather than contradictory. 
We do not estimate the aggregate impacts. Instead, we estimate how much more the 
most exposed companies onshored than their less exposed peers within the 
same industry. We find that response to be large and concentrated 
in R\&D-intensive firms---the heterogeneity the aggregate studies observe 
but cannot attribute. Second, on investment, \citet{CSZZ2024} and 
\citet{KennedyDobridgeLandefeldMortenson2022} identify responses to the 
reform's changes in the user cost of capital. Those changes were largely 
common to firms within an industry and are absorbed by fixed effects in our
regressions. 

The theory used to interpret our empirical findings distinguishes between two tax
margins. Firms choose where to book income based on average tax-rate differences
across jurisdictions and how much to invest based on marginal user costs.
This distinction follows \citet{GrubertSlemrod1998} and mirrors the standard
distinction between average and marginal effective tax rates in
\citet{DevereuxGriffith1998, DevereuxGriffith2003}. The theory combines the
user-cost framework of \citet{HallJorgenson1967}, the standard treatment of
deferral in \citet{Hartman1985}, and the concealment-cost approach to profit
shifting in \citet{HinesRice1994}.


\section{Data\label{sec:Data}}

The central contribution of our data construction is a new firm-year 
panel of income tax reconciliation tables extracted from SEC Form 
10-K filings. In this section, we first describe the reconciliation 
tables and then the procedures used to extract and link them to the
\citet{CompustatAnnual} database.\footnote{%
Details on the samples and methods used to extract data 
are provided in Appendix A and online code repository.}

\subsection{Reconciliation Table Description}

SEC filings of Form 10-K must be prepared in accordance with the U.S.~Generally
Accepted Accounting Principles (GAAP), which are officially organized and 
referenced within the Accounting Standards Codification (ASC). The specific
code of interest here is ASC 740, which governs how companies must account
for income taxes under GAAP. 
Companies are required to reconcile their effective corporate tax rate 
starting from the U.S.~federal statutory rate 
by disclosing reconciling adjustments 
such as 
foreign tax differentials,
federal credits,
state and local taxes,
valuation allowances,
and other accounting remeasurements in accordance with ASC 740.

In Table \ref{tab:Alphabet}, we provide an example of an income tax reconciliation
table, which was sourced from the 2018 annual report of Alphabet, Inc.
As is standard for larger filers, the table has three columns of data for the
reporting year and two prior years.  The first row of the table reports the federal statutory rate.
Over the period 2016--2018, Alphabet reported a rate of 35 percent 
for 2016 and 2017 and then 21 percent for 2018.\footnote{%
Not every filer reports 35 or 21 percent. Firms with non-calendar fiscal years
report a blended rate for the year straddling the enactment date, and
foreign-domiciled filers report their home-country rate. In all cases, we
measure the adjustments relative to the statutory rate reported in the
table itself.}
An important part of the parsing step is to find this row and the associated
columns with the filing years.
We also check if the company uses tax rates or tax amounts and, if the
latter, we convert the amounts to rates using the pretax income from
continuing operations reported in the filing itself.
As we can see from Table \ref{tab:Alphabet}, Alphabet reports rates.
The next set of rows reports the tax rate adjustments used
to compute the effective tax rate.  
These adjustment rows can be different company by company, but the first and last 
row of the table are always the same. 
Once we have extracted a table, we verify that the reported adjustments 
sum to the reported effective tax rate. Because reconciliation 
tables satisfy a simple accounting identity---the statutory rate plus 
all reported adjustments must equal the reported effective 
tax rate---this check validates the extracted data before they enter our
analysis. 

% 
% TABLE 1: Example Income Tax Reconciliation: Alphabet, 2018 Form 10-K}}
%

\begin{table}[!tp]
\centering
\begin{threeparttable}
\caption{\textsc{Example Income Tax Reconciliation: Alphabet, 2018 Form 10-K}}
\label{tab:Alphabet}

   \begin{tabular}{clcrrrc}
   \hline
   \hline
   \noalign{\vskip 5pt}
& \ \ \  &\phantom{helloooo} &   \textsc{2016}\phantom{)} &  \textsc{2017}\phantom{)} & \textsc{2018}\phantom{)} & \ \ \ \tabularnewline
   \noalign{\vskip 4pt}
   \cline{4-6}
   \noalign{\vskip 4pt}
&   \textsc{U.S.~federal statutory tax rate}      & &   35.0\phantom{)} &  35.0\phantom{)} & 21.0\phantom{)}&\tabularnewline
   \noalign{\vskip 5pt}
&   \ \ \ Foreign income taxed at different rates & &  (11.0)           & (14.2)           & (4.9) &\tabularnewline
   \noalign{\vskip 3pt}
&   \ \ \ TCJA one-time transition tax            & &    0.0\phantom{)} &  37.6\phantom{)} & (0.1) &\tabularnewline
   \noalign{\vskip 3pt}
&   \ \ \ TCJA deferred tax effects               & &    0.0\phantom{)} &  (1.4)           & (1.2) &\tabularnewline
   \noalign{\vskip 3pt}
&   \ \ \ Research credit                         & &   (2.0)           &  (1.8)           & (2.4) &\tabularnewline
   \noalign{\vskip 3pt}
&   \ \ \ Stock-based compensation expense        & &   (3.4)           &  (4.5)           & (2.2) &\tabularnewline
   \noalign{\vskip 3pt}
&   \ \ \ European Commission fine                & &    0.0\phantom{)} &   3.5\phantom{)} &  3.1\phantom{)} &\tabularnewline
   \noalign{\vskip 3pt}
&   \ \ \ Deferred tax asset valuation allowance  & &    0.1\phantom{)} &   0.9\phantom{)} & (2.0) &\tabularnewline
   \noalign{\vskip 3pt}
&   \ \ \ Other adjustments                       & &    0.6\phantom{)} &  (1.7)           &  0.7\phantom{)} &\tabularnewline
   \noalign{\vskip 3pt}
   \cline{4-6}
   \noalign{\vskip 3pt}
&   \textsc{Effective tax rate}                 & &  19.3\phantom{)}  &  53.4\phantom{)} & 12.0\phantom{)} &\tabularnewline
   \noalign{\vskip 4pt}
   \hline
   \end{tabular}
\bigskip
\begin{tablenotes}[flushleft]
\item[] \footnotesize{\textit{Notes:} Entries are percentage-point contributions to the effective tax rate. Parentheses denote negative contributions. The table reports Alphabet's reconciliation from the U.S.~federal statutory corporate tax rate to its effective tax rate for fiscal years 2016--2018.}
\end{tablenotes}
\end{threeparttable}

\bigskip
\bigskip
\end{table}
 

Alphabet's reconciliation shows large adjustments for the domestic--foreign
rate differential in every year and, in 2017--2018, for the one-time TCJA
items. We think of 2016---the first column of Table \ref{tab:Alphabet}---as 
a pre-TCJA tax year with the
two TCJA adjustments equal to zero. In that year, lower foreign
corporate tax rates implied an 11 percentage-point lower effective
tax rate when compared to the 35 percent
U.S.~federal statutory rate. The research credit row shows tax credits on R\&D
yielding a further $-$2 percentage-point adjustment. Another negative adjustment
is made for stock-based compensation expenses if the company's stock
price rises between the grant date and the vesting date. 
In the case of Alphabet, these expenses lowered the effective rate another
3.4 percentage points. 
The valuation-allowance row near the bottom of the table records 
changes in how much of its future tax
benefits the company expects to be able to use. When the company writes down
those benefits, its reported tax rate rises, as in 2016 and 2017. When it
restores some of them, its reported tax rate falls, as in 2018.
The final adjustment called ``other''
includes small adjustments not elsewhere classified.

When TCJA passed in 2017, new items appeared on the income
tax reconciliation as shown in the second column of Table \ref{tab:Alphabet},
namely, the one-time transition tax and the deferred tax
effects.  The transition tax---sometimes listed as IRC Section 965---is a tax 
on foreign earnings that had not yet been repatriated under the previous world-wide
tax system. Under TCJA, the United States moved toward a territorial system, in which
dividends from foreign subsidiaries are exempt from U.S.\ tax through a
100 percent dividends-received deduction. The one-time transition allowed companies to apply discounted tax rates:
15.5 percent on accumulated deferred foreign earnings held in cash-like assets and
8 percent on the remainder.  Despite the discounted rates,
the one-time transition implied a positive adjustment of 37.6 percentage
points for Alphabet.

Deferred tax effects reflect the remeasurement of existing deferred 
tax assets and liabilities following changes in tax law. Because 
TCJA reduced the statutory corporate tax rate from 35 to 21 percent, 
both were written down. The sign of the adjustment therefore depends
on a firm's net position: companies with net deferred tax liabilities
saw future obligations revalued downward and recorded a one-time
reduction in tax expense, while companies with net deferred tax
assets saw future deductions become less valuable and recorded an
increase. Both are common in our sample, and it is this variation
that the deferred-tax windfall analysis of Section \ref{sec:Windfall}
exploits. For Alphabet, the adjustments were $-$1.4 and $-$1.2
percentage points in 2017 and 2018, respectively.
In addition to the new TCJA items, Alphabet also paid non-deductible fines
to the European Commission, which generated a positive adjustment 
to the effective tax rate.

In Table \ref{tab:Taxonomy}, we provide a general taxonomy to classify
the different categories of adjustments to be used in our empirical
analysis.\footnote{%
In December 2023, the Financial Accounting Standards Board (FASB) issued
ASU 2023-09, which expands the required disclosure detail in the income tax
reconciliation tables. It requires the reconciliation to be disaggregated into
prescribed categories---including foreign tax effects by jurisdiction---and
reconciling items above a quantitative threshold of 5 percent of pretax income
times the statutory rate to be reported separately. The standard is effective
for public companies for annual periods beginning after December 15, 2024, so
apart from any early adopters it takes effect after our sample period ends and
does not affect the disclosures we extract. It should make reconciliation
tables more comparable going forward.}
The first category
on the list---which we will later call \emph{international}---is 
the primary focus
of our study. Here, we include foreign rate differentials, foreign
tax credits and withholding taxes, jurisdiction-specific tax holidays,
Subpart F inclusions, and the three international provisions of TCJA: GILTI, 
\emph{foreign-derived intangible income} (FDII), and 
\emph{base erosion and anti-abuse tax} (BEAT).
 GILTI and FDII both require companies to compute
a deemed 10 percent return on tangible assets and classify income in
excess of that return as residual income for further tax calculations.
GILTI subjects foreign residual income to a minimum tax, while FDII
provides a deduction for qualifying domestic residual income derived
from exporting goods and services. BEAT is a minimum tax on deductible
payments made to foreign affiliates. Subpart F, which predates TCJA, taxes some
categories of foreign income to the U.S.\ parent as it is earned
rather than when it is repatriated.
In the second category, which we 
call \emph{one time}, we include
adjustments arising from the enactment of a change in tax law,
including the provisional amounts firms record on enactment and the
revisions they make as those amounts are finalized. In the case of
TCJA, these are the Section 965 transition tax and the deferred tax
remeasurements discussed earlier in the case of Alphabet, which appear
mainly in 2017 and 2018 rather than in the year of enactment
alone.\footnote{%
SEC Staff Accounting Bulletin 118 granted a one-year measurement period
to complete the accounting, and fiscal-year filers recorded the
transition tax in fiscal 2018. Earlier reforms enter the same category:
the repatriation charges taken under the American Jobs Creation Act of
2004, which will be visible in Figure \ref{fig:Contributions}, are one-time
consequences of a change in tax law rather than an ongoing rate
differential.}

The remaining categories in Table \ref{tab:Taxonomy} are not directly
related to the international provisions of TCJA, though some are
affected by its other provisions. The reform repealed the Section 199
deduction, which sits in \emph{federal credits}, and lowered the
federal benefit of deducting state taxes, which mechanically raises the
\emph{state \& local} contribution at a 21 percent rate. Both are
useful for our descriptive analysis of the pre- and post-reform
period.
The category in the third row is \emph{Federal credits} and includes
federal tax credits such as the research credit and, before its repeal
by TCJA, the 9 percent Section 199 deduction for qualified production
activities inside the United States. Next is \emph{state \& local}, which is primarily
state corporate income taxes.
The \emph{valuation} category includes adjustments arising from the
measurement of deferred tax assets and tax reserves and from revisions
to prior-year estimates. These include valuation allowances, uncertain
tax positions, adjustments related to net operating losses, and
true-ups that reconcile the tax provision booked on financial
statements to the actual liability calculated on the filed tax return.
Purely temporary differences between tax and financial records do not
appear in the reconciliation at all: they change when tax is paid
and therefore leave the total tax expense---and hence the
effective rate---unchanged.
The \emph{book v.~tax} category includes adjustments related to
permanent differences between tax and financial records. 
These include non-deductible expenses, stock-based compensation,
goodwill, tax-exempt income, expenses related to mergers and
acquisitions, and percentage depletion in excess of cost depletion. 
Reconciliation items that are not elsewhere categorized (n.e.c.)~are 
included in \emph{other}. Using Table \ref{tab:Alphabet} as an example, the
foreign-rate row is \emph{international}, the two TCJA rows are
\emph{one time}, the research credit is \emph{federal credits}, stock
compensation and the European Commission fine are \emph{book v.~tax},
the valuation allowance is \emph{valuation}, and ``other adjustments''
is \emph{other}.\footnote{%
Categorizing the free-text row labels required a set of guidelines, 
which are detailed in Appendix \ref{app:categories}.}
Throughout what follows, a firm-year's contribution in a given category
is the sum of the percentage-point values of all its rows assigned to
that category, winsorized at plus or minus 15 percentage points.


% 
% TABLE 2: Taxonomy of Income Tax Reconciliation Adjustments
%

\begin{table}[!tp]
\centering
\begin{threeparttable}
\caption{\textsc{Taxonomy of Income Tax Reconciliation Adjustments}}
\label{tab:Taxonomy}

   \begin{tabular}{lc>{\raggedright\arraybackslash}p{0.62\textwidth}}
   \hline
   \hline
   \noalign{\vskip 5pt}
   \textsc{Category}     & &  \textsc{Included Tax Adjustments} \tabularnewline
   \noalign{\vskip 5pt}
   \hline
   \noalign{\vskip 6pt}
   \textsc{International}         & &  Foreign rate differentials; foreign tax credits and withholding; tax holidays; Subpart F; GILTI; FDII; BEAT \tabularnewline
   \noalign{\vskip 3pt}
   \textsc{One time}              & &  TCJA rate remeasurement; Section 965 transition tax \tabularnewline
   \noalign{\vskip 3pt}
   \textsc{Federal credits}       & &  R\&D credit; Section 199 domestic production deduction\tabularnewline
   \noalign{\vskip 3pt}
   \textsc{State \& local}       & &  State corporate income tax net of federal benefit\tabularnewline
   \noalign{\vskip 3pt}
   \textsc{Valuation}             & &  Valuation allowances; uncertain tax positions; NOL; true-ups\tabularnewline
   \noalign{\vskip 3pt}
   \textsc{Book v.~tax}          & &  Non-deductibles; stock comp; goodwill; tax-exempt income; M\&A; percentage depletion\tabularnewline
   \noalign{\vskip 3pt}
   \textsc{Other}                 & &  Unclassified residual; other adjustments n.e.c.\tabularnewline
   \noalign{\vskip 6pt}
   \hline
   \end{tabular}


\bigskip
\begin{tablenotes}[flushleft]
\item[] \footnotesize{\textit{Notes:} 
The categories are constructed from reconciliation line items reported in SEC Form 10-K filings.
GILTI denotes global intangible low-taxed income;
FDII denotes foreign-derived intangible income;
NOL denotes net operating losses;
M\&A denotes mergers and acquisitions; and
stock comp denotes stock-based compensation. 
True-ups are revisions to prior tax estimates resulting from updated information, audit resolutions, or amended filings.}
\end{tablenotes}
\end{threeparttable}

\bigskip
\bigskip
\end{table}





\subsection{Reconciliation Table Extraction}

We next describe the procedure used to extract income tax reconciliation data from company annual 
10-K reports. We start from the annual filings on the SEC's
\citet{EDGAR} system for firms in \citet{CompustatAnnual} with a
CIK--GVKEY link---forms 10-K, 10-K405, 10-KSB, 10-KSB40, 10-KT, and
their amendments---supplemented by a small number of manually audited
links for firms whose link was missing. The universe is therefore
Compustat-linked filings rather than all of EDGAR. In this section, we describe how
we first parse the 10-K filings to find tables and text most likely to contain the reconciliation
data and then parse the results from the first stage for a second stage in which
we construct a panel of standardized tables. The panel data are 
the inputs to our empirical analysis.

We begin with \CountNum{\ReconUniqueFilings}\ unique 10-K filings from \citet{EDGAR}.
Of these unique filings, \CountNum{\PhaseOneStatusOK} are available in 
HTML format suitable for our parsing procedures.
Each of these filings is parsed into text blocks and candidate 
tables, which are scored according to how
closely their language and structure resemble an income tax reconciliation table. 
For example, text blocks or tables  containing phrases such as ``effective tax rate,''
``statutory tax rate,'' or ``rate reconciliation'' would have high scores.
We use the scores to rank candidates, not to exclude filings: every
packet goes to the second stage, which turns it into standardized rows.
We parse with a rule-based procedure if the structure of the table
allows it and otherwise with a large language model. We then
apply an arithmetic check to every observation. This recovers
reconciliation data from \CountNum{\ReconPhaseTwoFilings}\ filings and
\CountNum{\PhaseTwoTotalRows}\ filing-year observations. Appendix A
describes the parsing steps and reports the counts at each of them.

Because these tables typically report three years of data, the same
firm-year often appears in as many as three filings. 
When a firm-year appears in more than one filing, we prefer the observation 
from the earliest filing year, then one that passes the arithmetic check, 
then an original 10-K over an amendment, and finally a higher-confidence extraction.
The repetition serves a dual purpose. It gives an independent check on the
arithmetic of each table, and it recovers firm-years whose own filing
could not be parsed. We find that \ShareNum{\PanelPriorYearPct}\ percent of the
panel comes from the prior-year column of a later filing rather than
from the filing for that year. Once we have eliminated duplicates, we link filings to
\citet{CompustatAnnual}, keeping observations for 1998--2024.
From this process, we have panel of
\CountNum{\AnalysisPanelFirmYears}\ firm-year reconciliations. 
As a final check, we compare the extracted effective
tax rate with the one implied by \citet{CompustatAnnual}: across
linked firm-years the two agree to within half a
percentage point in \ShareNum{\EtrAgreeHalfPct}\ percent of cases, with
a median absolute difference of \RatioNum{\EtrAgreeMedianAbs}\
percentage points. 

Over the period 1998--2024, the year-by-year coverage rate has
grown dramatically in line with higher usage of the HTML formatting.
In 1998, our coverage is
\ShareNum{\ReconCovInit}\ percent but 
\ShareNum{\ReconCovGR}\ percent by 2008 (see Appendix A for the full series).
In Table \ref{tab:Coverage},
we report the coverage rates when compared to all firm-years in \citet{CompustatAnnual}
for the period 2013--2024 of our empirical analysis.
Rows in the table are subsamples of the population of companies and
columns report coverage. The first column gives the number of
Compustat firm-years in the group, which for the top-100, top-250, and
top-500 groups is fixed by construction.  Our sample
contains \CountNum{\TopHundredCovered}\ of them, a coverage rate of
\ShareNum{\TopHundredEq}\ percent for that group. Across the full
population of \CountNum{\FullSampleTotal}\ firm-years we cover
\CountNum{\FullSampleCovered}, or \ShareNum{\FullSampleEq}\ percent. On an equal-weighted basis, 
the highest coverage
is achieved by firms with the most sales, which implies that the sales-weighted coverage
is even higher. The coverage rate for the full sample---when sales-weighted---is
\ShareNum{\FullSampleSales}\ percent and 
\ShareNum{\TopHundredSales}\ percent for companies in the top 100 of sales.

Averages are sales-weighted throughout, and industries follow the
SIC-based classification reported in Appendix Table \ref{tab:SICcodes}. Beyond
the reconciliation tables themselves, we use \citet{CompustatAnnual} for
sales, pretax income, capital expenditures, research and development,
employment, deferred taxes, and geographic segment sales.


%
% TABLE 3. Coverage of Extracted Tax Reconciliations, 2013--2024
%

\begin{table}[!tp]
\centering
\begin{threeparttable}
\caption{\textsc{Coverage of Extracted Tax Reconciliations, 2013--2024}}
\label{tab:Coverage}

   \begin{tabular}{lrrrr}
   \hline
   \hline
   \noalign{\vskip 5pt}
                       &\multicolumn{1}{c}{\textsc{Compustat}}
                           &\multicolumn{1}{c}{\textsc{Covered}} 
                           &\multicolumn{1}{c}{\textsc{Equal-weighted}}
                           &\multicolumn{1}{c}{\textsc{Sales-weighted}}\tabularnewline
   \textsc{Group }         &\multicolumn{1}{c}{\ \textsc{Firm-years}\ \ }
                           &\multicolumn{1}{c}{\ \ \textsc{Firm-years}\ \ }  
                           &\multicolumn{1}{c}{\textsc{Share (\%)}}  
                           &\multicolumn{1}{c}{\textsc{Share (\%)}}\tabularnewline
   \noalign{\vskip 5pt}
   \hline
   \noalign{\vskip 6pt}
   \textsc{Top 100}        & \CountNum{\TopHundredTotal}\phantom{hello}
                           & \CountNum{\TopHundredCovered}\phantom{hello}  
                           & \ShareNum{\TopHundredEq}\phantom{helloooo}
                           & \ShareNum{\TopHundredSales}\phantom{helloooo}\tabularnewline
   \noalign{\vskip 3pt}
   \textsc{Top 250}        & \CountNum{\TopTwoFiftyTotal}\phantom{hello}
                           & \CountNum{\TopTwoFiftyCovered}\phantom{hello}   
                           & \ShareNum{\TopTwoFiftyEq}\phantom{helloooo}
                           & \ShareNum{\TopTwoFiftySales}\phantom{helloooo}\tabularnewline
   \noalign{\vskip 3pt}
   \textsc{Top 500}        & \CountNum{\TopFiveHundredTotal}\phantom{hello}
                           & \CountNum{\TopFiveHundredCovered}\phantom{hello}   
                           & \ShareNum{\TopFiveHundredEq}\phantom{helloooo}
                           & \ShareNum{\TopFiveHundredSales}\phantom{helloooo}\tabularnewline
   \noalign{\vskip 3pt}
   \textsc{Full Sample}    & \CountNum{\FullSampleTotal}\phantom{hello}
                           & \CountNum{\FullSampleCovered}\phantom{hello}   
                           & \ShareNum{\FullSampleEq}\phantom{helloooo}
                           & \ShareNum{\FullSampleSales}\phantom{helloooo}\tabularnewline
   \noalign{\vskip 6pt}
   \hline
   \end{tabular}


\bigskip
\begin{tablenotes}[flushleft]
\item[] \footnotesize{\textit{Notes:} The Compustat population is
\citet{CompustatAnnual} firm-years for fiscal years 2013--2024 that are
incorporated in the United States, excluding depository institutions,
investment offices, and real estate investment trusts. The
top 100, 250, and 500 are those with the highest sales within the year,
so their Compustat firm-year counts are fixed by construction. The
equal-weighted share is the average across years of the covered share
of firm-years. The sales-weighted share is the sales of covered
firm-years as a fraction of the sales of all firm-years in the group.}

\end{tablenotes}
\end{threeparttable}

\bigskip
\bigskip
\end{table}




\section{Results\label{sec:Results}}

In this section, we use the income tax reconciliation panel dataset described
in Section \ref{sec:Data} to investigate the impact of TCJA on effective tax rates,
profit shifting, real activity, and company balance sheets. 
We start with descriptive evidence that motivates our
empirical strategies and case studies.

\subsection{Descriptive Evidence\label{sec:DescriptiveEvidence}}

In this section, we report summary statistics for our income tax reconciliation
database---at the aggregate and disaggregate levels and by industry.
We document that the largest
post-TCJA changes in the gap between effective and statutory corporate tax
 rates are attributed to a reduction in foreign tax advantages.
The figures and tables that follow use every covered firm-year with
positive sales, whether or not the firm was profitable:
\CountNum{\PreTotalCounts}\ firm-years in 2013--2016 and
\CountNum{\PostTotalCounts}\ in 2019--2024, with 2017 and 2018 omitted
as transition years around enactment.

\subsubsection{Aggregate}

In Figure \ref{fig:Contributions}, we plot contributions to the 
gap between the effective corporate tax rates and the U.S.~federal  statutory rate
over the period 2004 to 2024, when coverage of the companies
for which we have parsed reconciliation tables
is at least two-thirds of the merged \citet{CompustatAnnual} sample.
Each colored segment of the figure shows the sales-weighted 
average contribution to the rate gap---in percentage points---of 
that income tax reconciliation category for the given year.
These categories match the list in our taxonomy of Table
\ref{tab:Taxonomy} and are constructed as described in Section
\ref{sec:Data}. 
Across the sample, we see that some categories---like state and local corporate
taxes---increase effective rates, and some categories---like federal tax 
credits---decrease effective rates.
The dotted line marks the TCJA enactment and
hollow circles show the sales-weighted average total gap in the year.

We use Figure \ref{fig:Contributions} to provide a stark visual
description of the aggregate impact of TCJA on the effective tax rates paid
by publicly traded firms reporting to the SEC.
The most notable change is the collapse of the international contribution 
(in red).
Prior to TCJA, the international contribution provided the largest tax relief
on average. The post-TCJA collapse is
especially noticeable against the other categories, none of which
changes by more than \ShareNum{\FigMaxNonIntlChange}\ percentage points
in absolute value between the two periods, against
\ShareNum{\FigIntlChange}\ for international.
Averaging over 2013--2016, 
we find a total gap of 
\ShareNum{\FigPreTotalAvg}\ percentage points, with the international contribution equal to 
\ShareNum{\FigPreIntlAvg}.
Averaging over 2019--2024,
the total gap falls to
\ShareNum{\FigPostTotalAvg}\ percentage points with the international contribution down to
\ShareNum{\FigPostIntlAvg}.


%
% FIGURE 1: Contributions to the Gap Between Statutory and Effective Tax Rates
%
\begin{figure}[!tp]
\centering
\caption{\textsc{Contributions to the Gap Between Statutory and Effective Tax Rates}}
\label{fig:Contributions}
\begin{tikzpicture}
\pgfplotstableread[col sep=comma]{\Inputs/decomp_tikz/fig_drivers_tikz_data.csv}\figdriversdata

\begin{axis}[
    width=16cm,
    height=9cm,
    ymin=-16.8,
    ymax=6.5,
    xmin=2003.65,
    xmax=2024.65,
    axis lines=left,
    ybar stacked,
    stack negative=separate,
    bar width=8.5pt,
    enlarge x limits=0.01,
    xtick={1998,2000,2002,2004,2006,2008,2010,2012,2014,2016,2018,2020,2022,2024},
    xticklabel style={
       font=\scriptsize,
       /pgf/number format/fixed,
       /pgf/number format/precision=0,
       /pgf/number format/1000 sep={}},
    yticklabel style={font=\scriptsize},
    ylabel={\textsc{Contribution to Tax Rate Gap (pp)}},
    tick align=outside,
    grid=none,
    legend style={
        at={(0.5,-0.10)},
        anchor=north,
        draw=none,
        font=\scriptsize,
        /tikz/every even column/.append style={column sep=0.45cm}
    },
    legend columns=4,
]

\draw[dashed, line width=0.9pt, color=black!55] (axis cs:2017.5,-18) -- (axis cs:2017.5,6);
\node[font=\scriptsize\itshape, text=black!60] at (axis cs:2018.2,5.0) {TCJA};

% 1. International (Vibrant Crimson Red)
\addplot+[ybar, draw=white, line width=0.2pt, fill=IntRed,
 legend image code/.code={\draw[fill=IntRed, draw=IntRed](0cm,0.0cm) rectangle (0.20cm,0.15cm);}]
table[x=year, y=international] {\figdriversdata}; \addlegendentry{\textsc{International}} 

% 2. One time (Bright Burnt Orange)
\addplot+[ybar, draw=white, line width=0.2pt, fill=OneOrange,
 legend image code/.code={\draw[fill=OneOrange, draw=OneOrange](0cm,0.0cm) rectangle (0.20cm,0.15cm);}]
table[x=year, y=tcja_onetime] {\figdriversdata}; \addlegendentry{\textsc{One time}} 

% 3. Federal credits (Vibrant Amber Gold)
\addplot+[ybar, draw=white, line width=0.2pt, fill=FedGold,
 legend image code/.code={\draw[fill=FedGold, draw=FedGold](0cm,0.0cm) rectangle (0.20cm,0.15cm);}]
table[x=year, y=federal_credits] {\figdriversdata}; \addlegendentry{\textsc{Federal credits}} 

% 4. State & local (Clean Kelly Green)
\addplot+[ybar, draw=white, line width=0.2pt, fill=StateGreen,
 legend image code/.code={\draw[fill=StateGreen, draw=StateGreen](0cm,0.0cm) rectangle (0.20cm,0.15cm);}]
table[x=year, y=state_local] {\figdriversdata}; \addlegendentry{\textsc{State \& local}}

% 5. Valuation (Deep Peacock Teal)
\addplot+[ybar, draw=white, line width=0.2pt, fill=ValTeal,
 legend image code/.code={\draw[fill=ValTeal, draw=ValTeal](0cm,0.0cm) rectangle (0.20cm,0.15cm);}]
table[x=year, y=reserves_losses] {\figdriversdata}; \addlegendentry{\textsc{Valuation}} 

% 6. Book v. Tax (Bright Sky Blue - Shifted from royal blue to create distinction)
\addplot+[ybar, draw=white, line width=0.2pt, fill=BookBlue,
 legend image code/.code={\draw[fill=BookBlue, draw=BookBlue](0cm,0.0cm) rectangle (0.20cm,0.15cm);}]
table[x=year, y=permanent_items] {\figdriversdata}; \addlegendentry{\textsc{Book v.~tax}} 

% 7. Other (Orchid Violet - Clear separation at the bottom of the rainbow)
\addplot+[ybar, draw=white, line width=0.2pt, fill=OthViolet,
 legend image code/.code={\draw[fill=OthViolet, draw=OthViolet](0cm,0.0cm) rectangle (0.20cm,0.15cm);}]
table[x=year, y=other] {\figdriversdata}; \addlegendentry{\textsc{Other}}

\addplot+[
    only marks,
    color=black,
    stack plots=false,
    mark=*,
    mark size=1.8pt,
    mark options={
        solid,
        fill=white,
        draw=black,
        line width=0.8pt},
    legend image code/.code={\draw[mark=*, 
           mark options={fill=white,line width=0.8pt}] plot coordinates {(0.27em,0.3ex)};}
]
table[x=year, y=total_gap_rev_wtd] {\figdriversdata};
\addlegendentry{\textsc{Total Gap}}

\end{axis}
\end{tikzpicture}

\bigskip

\begin{minipage}{0.95\textwidth}
\footnotesize{\textit{Notes:} The figure reports the sales-weighted average contribution of each reconciliation category to the gap between the effective corporate tax rates of companies and the U.S.\ federal statutory rate.  Negative values reduce the effective tax rate relative to the statutory rate. The dashed vertical line marks the enactment of the Tax Cuts and Jobs Act (TCJA) of 2017.}
\end{minipage}

\bigskip
\bigskip
\end{figure}



\subsubsection{Distributional}

In Table \ref{tab:Distribution}, we show the full distributional results of the
tax rate gaps.
The first column shows the firm-year counts, which differ across
categories because a firm-year enters a category only if the firm
reports an entry.
The second column shows the sales-weighted averages for 
2013--2016 in the pre-TCJA panel and 2019--2024 in the post-TCJA panel.
As in Figure \ref{fig:Contributions}, the largest average change
occurs in the international category.
In the last columns, we report selected sales-weighted percentiles of the
distribution.
Here, we see how much variation there is and how the
entire distribution shifts toward zero after TCJA. 
For example, the international contribution 
has a pre-TCJA median gap of 
\ShareNum{\PreIntlPFZ}\ percentage points and an interquartile range of
[\ShareNum{\PreIntlPTF},\ShareNum{\PreIntlPSF}].
Because of its importance, the wide range in this category translates 
into a wide range for the overall tax rate gap.
After TCJA, the range for the international contribution collapses to 
[\ShareNum{\PostIntlPTF},\ShareNum{\PostIntlPSF}] 
with the median close to zero.
A perusal of the remaining categories also shows relative stability
across these periods.  For example, the distribution of state and local tax 
contributions changes only slightly.

%
% TABLE 4: Distribution of the Tax Rate Gap by Contribution
%

\begin{table}[!tp]
\centering
\begin{threeparttable}
\caption{\textsc{Distribution of the Tax Rate Gap by Contribution}}
\label{tab:Distribution}

   \begin{tabular}{lcrrrrr}
   \hline
   \hline
   \noalign{\vskip 5pt}
           &\phantom{hell} &
                           &
                           &\multicolumn{3}{c}{\textsc{Percentiles}} \tabularnewline
   \noalign{\vskip 2pt}
   \cline{5-7}
   \noalign{\vskip 5pt}
   \textsc{Contribution} & &\multicolumn{1}{c}{\textsc{Counts\ \ }}
                           &\multicolumn{1}{c}{\textsc{Mean}}
                           &\multicolumn{1}{c}{\textsc{25$^{\rm th}$}}
                           &\multicolumn{1}{c}{\textsc{50$^{\rm th}$}}
                           &\multicolumn{1}{c}{\textsc{75$^{\rm th}$}}\tabularnewline
   \noalign{\vskip 5pt}
   \hline
   \noalign{\vskip 6pt}
           &               &\multicolumn{5}{c}{\textsc{Pre-TCJA}} \tabularnewline
   \noalign{\vskip 2pt}
   \cline{3-7}
   \noalign{\vskip 6pt}
   \textsc{International}& & \CountNum{\PreIntlCounts}\phantom{he}  
                           & \ShareNum{\PreIntlAvg}\phantom{hel}
                           & \ShareNum{\PreIntlPTF}\phantom{hel}
                           & \ShareNum{\PreIntlPFZ}\phantom{hel}
                           & \ShareNum{\PreIntlPSF}\phantom{hel}\tabularnewline
   \noalign{\vskip 3pt}
   \textsc{One time}     & & \CountNum{\PreOneTimeCounts}\phantom{he}  
                           & \ShareNum{\PreOneTimeAvg}\phantom{hel}
                           & \ShareNum{\PreOneTimePTF}\phantom{hel}
                           & \ShareNum{\PreOneTimePFZ}\phantom{hel}
                           & \ShareNum{\PreOneTimePSF}\phantom{hel}\tabularnewline
   \noalign{\vskip 3pt}
   \textsc{Federal Credits}& & \CountNum{\PreFedCredCounts}\phantom{he}  
                           & \ShareNum{\PreFedCredAvg}\phantom{hel}
                           & \ShareNum{\PreFedCredPTF}\phantom{hel}
                           & \ShareNum{\PreFedCredPFZ}\phantom{hel}
                           & \ShareNum{\PreFedCredPSF}\phantom{hel}\tabularnewline
   \noalign{\vskip 3pt}
   \textsc{State \& local}& & \CountNum{\PreStateCounts}\phantom{he}  
                           & \ShareNum{\PreStateAvg}\phantom{hel}
                           & \ShareNum{\PreStatePTF}\phantom{hel}
                           & \ShareNum{\PreStatePFZ}\phantom{hel}
                           & \ShareNum{\PreStatePSF}\phantom{hel}\tabularnewline
   \noalign{\vskip 3pt}
   \textsc{Valuation}    & & \CountNum{\PreValCounts}\phantom{he}  
                           & \ShareNum{\PreValAvg}\phantom{hel}
                           & \ShareNum{\PreValPTF}\phantom{hel}
                           & \ShareNum{\PreValPFZ}\phantom{hel}
                           & \ShareNum{\PreValPSF}\phantom{hel}\tabularnewline
   \noalign{\vskip 3pt}
   \textsc{Book v.~tax}  & & \CountNum{\PreBookTaxCounts}\phantom{he}  
                           & \ShareNum{\PreBookTaxAvg}\phantom{hel}
                           & \ShareNum{\PreBookTaxPTF}\phantom{hel}
                           & \ShareNum{\PreBookTaxPFZ}\phantom{hel}
                           & \ShareNum{\PreBookTaxPSF}\phantom{hel}\tabularnewline
   \noalign{\vskip 3pt}
   \textsc{Other}        & & \CountNum{\PreOtherCounts}\phantom{he}  
                           & \ShareNum{\PreOtherAvg}\phantom{hel}
                           & \ShareNum{\PreOtherPTF}\phantom{hel}
                           & \ShareNum{\PreOtherPFZ}\phantom{hel}
                           & \ShareNum{\PreOtherPSF}\phantom{hel}\tabularnewline
   \noalign{\vskip 3pt}
   \textsc{Total Gap}    & & \CountNum{\PreTotalCounts}\phantom{he}  
                           & \ShareNum{\PreTotalAvg}\phantom{hel}
                           & \ShareNum{\PreTotalPTF}\phantom{hel}
                           & \ShareNum{\PreTotalPFZ}\phantom{hel}
                           & \ShareNum{\PreTotalPSF}\phantom{hel}\tabularnewline
   \noalign{\vskip 8pt}
           &               &\multicolumn{5}{c}{\textsc{Post-TCJA}} \tabularnewline
   \noalign{\vskip 2pt}
   \cline{3-7}
   \noalign{\vskip 6pt}
   \textsc{International}& & \CountNum{\PostIntlCounts}\phantom{he}  
                           & \ShareNum{\PostIntlAvg}\phantom{hel}
                           & \ShareNum{\PostIntlPTF}\phantom{hel}
                           & \ShareNum{\PostIntlPFZ}\phantom{hel}
                           & \ShareNum{\PostIntlPSF}\phantom{hel}\tabularnewline
   \noalign{\vskip 3pt}
   \textsc{One time}     & & \CountNum{\PostOneTimeCounts}\phantom{he}  
                           & \ShareNum{\PostOneTimeAvg}\phantom{hel}
                           & \ShareNum{\PostOneTimePTF}\phantom{hel}
                           & \ShareNum{\PostOneTimePFZ}\phantom{hel}
                           & \ShareNum{\PostOneTimePSF}\phantom{hel}\tabularnewline
   \noalign{\vskip 3pt}
   \textsc{Federal Credits}& & \CountNum{\PostFedCredCounts}\phantom{he}  
                           & \ShareNum{\PostFedCredAvg}\phantom{hel}
                           & \ShareNum{\PostFedCredPTF}\phantom{hel}
                           & \ShareNum{\PostFedCredPFZ}\phantom{hel}
                           & \ShareNum{\PostFedCredPSF}\phantom{hel}\tabularnewline
   \noalign{\vskip 3pt}
   \textsc{State \& local}& & \CountNum{\PostStateCounts}\phantom{he}  
                           & \ShareNum{\PostStateAvg}\phantom{hel}
                           & \ShareNum{\PostStatePTF}\phantom{hel}
                           & \ShareNum{\PostStatePFZ}\phantom{hel}
                           & \ShareNum{\PostStatePSF}\phantom{hel}\tabularnewline
   \noalign{\vskip 3pt}
   \textsc{Valuation}    & & \CountNum{\PostValCounts}\phantom{he}  
                           & \ShareNum{\PostValAvg}\phantom{hel}
                           & \ShareNum{\PostValPTF}\phantom{hel}
                           & \ShareNum{\PostValPFZ}\phantom{hel}
                           & \ShareNum{\PostValPSF}\phantom{hel}\tabularnewline
   \noalign{\vskip 3pt}
   \textsc{Book v.~tax}  & & \CountNum{\PostBookTaxCounts}\phantom{he}  
                           & \ShareNum{\PostBookTaxAvg}\phantom{hel}
                           & \ShareNum{\PostBookTaxPTF}\phantom{hel}
                           & \ShareNum{\PostBookTaxPFZ}\phantom{hel}
                           & \ShareNum{\PostBookTaxPSF}\phantom{hel}\tabularnewline
   \noalign{\vskip 3pt}
   \textsc{Other}        & & \CountNum{\PostOtherCounts}\phantom{he}  
                           & \ShareNum{\PostOtherAvg}\phantom{hel}
                           & \ShareNum{\PostOtherPTF}\phantom{hel}
                           & \ShareNum{\PostOtherPFZ}\phantom{hel}
                           & \ShareNum{\PostOtherPSF}\phantom{hel}\tabularnewline
   \noalign{\vskip 3pt}
   \textsc{Total Gap}    & & \CountNum{\PostTotalCounts}\phantom{he}  
                           & \ShareNum{\PostTotalAvg}\phantom{hel}
                           & \ShareNum{\PostTotalPTF}\phantom{hel}
                           & \ShareNum{\PostTotalPFZ}\phantom{hel}
                           & \ShareNum{\PostTotalPSF}\phantom{hel}\tabularnewline
   \noalign{\vskip 6pt}
   \hline
   \end{tabular}

\bigskip
\begin{tablenotes}[flushleft]
\item[] \footnotesize{\textit{Notes:} 
Entries report the distribution of the contribution of each reconciliation category to the gap between the
effective tax rate and the statutory tax rate, in percentage points. The \emph{Total gap} 
row is computed on the full
sample using sales as weights. All other rows are computed using sales as weights 
for the subset of firm-year observations in which that contribution category is present in the reconciliation. 
The pre-TCJA period is 2013--2016 and the post-TCJA period is 2019–-2024.
Transition years 2017-–2018 are omitted. Negative values indicate items that reduce the
effective tax rate relative to the statutory rate.}
\end{tablenotes}
\end{threeparttable}

\bigskip
\bigskip
\end{table}





\subsubsection{By Industry}


Next, we investigate the differences in income tax reconciliation contributions before
and after TCJA for selected industries.\footnote{% 
See the Appendix A for details on industry categorizations by SIC code.}
As before, we compute the difference between the effective and statutory rates.
In Figure \ref{fig:Industry}, we report sales-weighted industry
averages of the tax gap in percentage points (pp) in two panels. 
Panel A shows the average total gap for 2013--2016 and 2019--2024 by industry. Panel B
shows the change in the average reconciliation contributions, that is,
the post-TCJA averages less the pre-TCJA averages.

%
% FIGURE 2: Industry Differences in Tax Rate Gap Before and After TCJA
%
\pgfplotstableread[col sep=comma]{\Inputs/decomp_ind_new_tikz/summary.csv}\indgapdata
\providecommand{\IndGapPanelWidth}{0.46\linewidth}
\providecommand{\IndGapPanelHeight}{6.4cm}
\newcommand{\IndGapSwatch}[1]{%
  \tikz[baseline=-0.55ex]\filldraw[fill=#1,draw=#1] (0,0) rectangle (0.18cm,0.12cm); }

\begin{figure}[!tp]
\centering
\caption{\textsc{Industry Differences in Tax Rate Gap Before and After TCJA}}
\label{fig:Industry}

\begin{tikzpicture}
\begin{groupplot}[
  group style={
    group size=2 by 1,
    horizontal sep=0.5cm,
  },
  height=\IndGapPanelHeight,
  width=\IndGapPanelWidth,
  y dir=reverse,
  ymin=-0.5,
  ymax=6.5,
  ytick={0,1,2,3,4,5,6},
  yticklabels={\textsc{Technology},\textsc{Pharmaceuticals},\textsc{Manufacturing},\textsc{Oil \& Gas},
               \textsc{Retail},\textsc{Utilities},\textsc{Other}},
  yticklabel style={font=\scriptsize, anchor=east, xshift=-1pt},
  tick align=outside,
  axis x line=bottom,
  axis y line=left,
  axis line style={-},
  grid=none,
]

\nextgroupplot[
  xmin=-20,
  xmax=8,
  xtick={-20,-15,-10,-5,0,5},
  title={\textsc{A. Average Total Gap}},
  title style={font=\fontsize{9}{9}\selectfont},
  xticklabel style={font=\scriptsize},
]
\addplot[
  only marks,
  mark=*,
  mark size=2.8pt,
  draw=Indpre,
  fill=Indpre,
  point meta=explicit symbolic,
  nodes near coords,
  every node near coord/.append style={font=\tiny, anchor=south west, xshift=-8pt, yshift=-1pt},
] table[x=gap_pre_pp, y=y, meta=gap_pre_lbl] {\indgapdata};

\addplot[
  only marks,
  mark=square*,
  mark size=2.3pt,
  draw=Indpost,
  fill=Indpost,
  point meta=explicit symbolic,
  nodes near coords,
  every node near coord/.append style={font=\tiny, anchor=south west, xshift=-8pt, yshift=-12pt},
] table[x=gap_post_pp, y=y, meta=gap_post_lbl] {\indgapdata};


\pgfplotsextra{
   \foreach \row/\pos in {0/0,1/1,2/2,3/3,4/4,5/5,6/6} {
    \pgfplotstablegetelem{\row}{gap_pre_pp}\of\indgapdata
    \pgfmathsetmacro{\xA}{\pgfplotsretval}

    \pgfplotstablegetelem{\row}{gap_post_pp}\of\indgapdata
    \pgfmathsetmacro{\xB}{\pgfplotsretval}

    \draw[->, thick, black!60]
      (axis cs:\xA,\pos) -- (axis cs:\xB,\pos);
  }
}

\node[
  anchor=south west,
  font=\scriptsize,
  align=left
] at (rel axis cs:0.03,0.03) {%
{\color{Indpre}\scalebox{1.6}{$\bullet$}}\,2013--16\\[-1pt]
{\color{Indpost}\scalebox{1.0}{$\blacksquare$}}\,2019--24
};


\draw[black, line width=0.8pt] (axis cs:0,-0.5) -- (axis cs:0,6.5);

\nextgroupplot[
  xbar stacked,
  stack negative=separate,
  /pgf/bar width=10pt,
  xmin=-8,
  xmax=15,
  xtick={-5,0,5,10,15},
  title={\textsc{B. Change in Contributions}},
  title style={font=\fontsize{9}{9}\selectfont},
  ytick=\empty,
  xticklabel style={font=\scriptsize},
]
\addplot+[draw=white, line width=0.2pt, fill=IntRed] table[x=international_delta_pp, y=y] {\indgapdata};
\addplot+[draw=white, line width=0.2pt, fill=OneOrange] table[x=tcja_onetime_delta_pp, y=y] {\indgapdata};
\addplot+[draw=white, line width=0.2pt, fill=FedGold] table[x=federal_credits_delta_pp, y=y] {\indgapdata};
\addplot+[draw=white, line width=0.2pt, fill=StateGreen] table[x=state_local_delta_pp, y=y] {\indgapdata};
\addplot+[draw=white, line width=0.2pt, fill=ValTeal] table[x=reserves_losses_delta_pp, y=y] {\indgapdata};
\addplot+[draw=white, line width=0.2pt, fill=BookBlue] table[x=permanent_items_delta_pp, y=y] {\indgapdata};
\addplot+[draw=white, line width=0.2pt, fill=OthViolet] table[x=other_delta_pp, y=y] {\indgapdata};

\draw[black, line width=0.8pt] (axis cs:0,-0.5) -- (axis cs:0,6.5);

\end{groupplot}

\path (group c1r1.south west) -- coordinate[pos=0.5] (indgaplegendanchor) (group c2r1.south east);

\node[
  anchor=north,
  font=\scriptsize,
  align=center,
  inner sep=0pt,
] at ([yshift=-0.65cm]indgaplegendanchor) {%
  \setlength{\tabcolsep}{8pt}%
  \renewcommand{\arraystretch}{0.95}%

  \begin{tabular}{cccc}
  \IndGapSwatch{IntRed} \textsc{International} &
  \IndGapSwatch{OneOrange} \textsc{One time} &
  \IndGapSwatch{FedGold} \textsc{Federal credits} &
  \IndGapSwatch{StateGreen} \textsc{State \& local} 
  \end{tabular}

  \\[-0.1em]

  \begin{tabular}{ccc}
  \IndGapSwatch{ValTeal} \textsc{Valuation} &
  \IndGapSwatch{BookBlue} \textsc{Book v.~tax} &
  \IndGapSwatch{OthViolet} \textsc{Other}
  \end{tabular}%
};

\end{tikzpicture}

\bigskip

\begin{minipage}{0.95\textwidth}
\footnotesize
\textit{Notes:} 
The figure reports sales-weighted industry averages of the gap between the
effective tax rate and the U.S.\ federal statutory corporate tax rate, and changes in the
average contribution of each reconciliation category to that gap. 
The industries are categorized using SIC codes reported in the appendix.
Panel A shows the average
total gap for 2013--2016 and 2019--2024. Panel B shows the change in average reconciliation
contributions between the two periods (2019--2024 minus 2013--2016). Contributions are
measured in percentage points. Negative values reduce the effective tax rate relative to
the statutory rate. The sample is the full sample, including loss firms. Years 2017 and
2018 are omitted as transition years around TCJA enactment. Industries follow the SIC-based
classification used elsewhere in the paper.
\end{minipage}

\bigskip
\bigskip
\end{figure}




Consider the total gap changes in Panel A of Figure \ref{fig:Industry}.
Here, we report sales-weighted averages for 2013--2016 (circles)
and 2019--2024 (squares).  As an example, the top row shows the changes
for technology companies that had an average tax rate gap of 
\ShareNum{\FigTwoTechPreGap} percentage points in the pre-TCJA period and
\ShareNum{\FigTwoTechPostGap} percentage points in the post-TCJA period.
As the figure shows, the tax rate gaps close for all industries except utilities.
Panel B shows that the main source of the closure---with utilities again
the exception---is the international category of reconciliations.
This closure contributes \ShareNum{\FigTwoTechIntlDelta},
\ShareNum{\FigTwoPharmaIntlDelta}, \ShareNum{\FigTwoManufIntlDelta}, and
\ShareNum{\FigTwoOilIntlDelta}\ percentage points for technology,
pharmaceuticals, manufacturing, and oil and gas, respectively.

In Table \ref{tab:DistributionByInd}, we reconsider the distributional
impacts of TCJA, in this case by industry.
We report the counts in the first column with the pre-TCJA data
in the top panel and the post-TCJA in the bottom.
The sales-weighted averages from Figure~\ref{fig:Industry}
are repeated here in the second column for convenience.
The final columns show the sales-weighted 25$^{\rm th}$, 50$^{\rm th}$,
and 75$^{\rm th}$ percentiles of the distribution for the
pre- and post-TCJA periods.
Comparing these periods, we find a significant shift to 
higher rates across industries and a narrowing of the interquartile
ranges for all but utilities. 

The main takeaways from the descriptive evidence are twofold.
First, the largest change in the tax gap decomposition comes from
the international reconciliation category.  Second, aggregate 
effective tax rates hide significant heterogeneity across companies and industries.


%
% TABLE 5: Distribution of the Tax Rate Gap by Industry
%

\begin{table}[!tp]
\centering
\begin{threeparttable}
\caption{\textsc{Distribution of the Tax Rate Gap by Industry}}
\label{tab:DistributionByInd}

   \begin{tabular}{lcrrrrr}
   \hline
   \hline
   \noalign{\vskip 5pt}
           &\phantom{h} &
                           &
                           &\multicolumn{3}{c}{\textsc{Percentiles}} \tabularnewline
   \noalign{\vskip 2pt}
   \cline{5-7}
   \noalign{\vskip 5pt}
   \textsc{Industry}   &   &\multicolumn{1}{c}{\textsc{Counts}}
                           &\multicolumn{1}{c}{\textsc{Mean}}
                           &\multicolumn{1}{c}{\textsc{25$^{\rm th}$}}
                           &\multicolumn{1}{c}{\textsc{50$^{\rm th}$}}
                           &\multicolumn{1}{c}{\textsc{75$^{\rm th}$}}\tabularnewline
   \noalign{\vskip 5pt}
   \hline
   \noalign{\vskip 6pt}
           &               &\multicolumn{5}{c}{\textsc{Pre-TCJA}} \tabularnewline
   \noalign{\vskip 2pt}
   \cline{3-7}
   \noalign{\vskip 6pt}
   \textsc{Technology}     & & \CountNum{\PreTechCounts}\phantom{he}
                           & \ShareNum{\PreTechAvg}\phantom{hel}
                           & \ShareNum{\PreTechPTF}\phantom{hell}
                           & \ShareNum{\PreTechPFZ}\phantom{hell}
                           & \ShareNum{\PreTechPSF}\phantom{hell}\tabularnewline
   \noalign{\vskip 3pt}
   \textsc{Pharmaceuticals}& & \CountNum{\PrePharmaCounts}\phantom{he}
                           & \ShareNum{\PrePharmaAvg}\phantom{hel}
                           & \ShareNum{\PrePharmaPTF}\phantom{hell}
                           & \ShareNum{\PrePharmaPFZ}\phantom{hell}
                           & \ShareNum{\PrePharmaPSF}\phantom{hell}\tabularnewline
   \noalign{\vskip 3pt}
   \textsc{Manufacturing}  & & \CountNum{\PreManufCounts}\phantom{he}
                           & \ShareNum{\PreManufAvg}\phantom{hel}
                           & \ShareNum{\PreManufPTF}\phantom{hell}
                           & \ShareNum{\PreManufPFZ}\phantom{hell}
                           & \ShareNum{\PreManufPSF}\phantom{hell}\tabularnewline
   \noalign{\vskip 3pt}
   \textsc{Oil \& gas}  & & \CountNum{\PreOilCounts}\phantom{he}
                           & \ShareNum{\PreOilAvg}\phantom{hel}
                           & \ShareNum{\PreOilPTF}\phantom{hell}
                           & \ShareNum{\PreOilPFZ}\phantom{hell}
                           & \ShareNum{\PreOilPSF}\phantom{hell}\tabularnewline
   \noalign{\vskip 3pt}
   \textsc{Retail}       & & \CountNum{\PreRetailCounts}\phantom{he}
                           & \ShareNum{\PreRetailAvg}\phantom{hel}
                           & \ShareNum{\PreRetailPTF}\phantom{hell}
                           & \ShareNum{\PreRetailPFZ}\phantom{hell}
                           & \ShareNum{\PreRetailPSF}\phantom{hell}\tabularnewline
   \noalign{\vskip 3pt}
   \textsc{Utilities}    & & \CountNum{\PreUtilCounts}\phantom{he}
                           & \ShareNum{\PreUtilAvg}\phantom{hel}
                           & \ShareNum{\PreUtilPTF}\phantom{hell}
                           & \ShareNum{\PreUtilPFZ}\phantom{hell}
                           & \ShareNum{\PreUtilPSF}\phantom{hell}\tabularnewline
   \noalign{\vskip 3pt}
   \textsc{Other}        & & \CountNum{\PreIndOtherCounts}\phantom{he}
                           & \ShareNum{\PreIndOtherAvg}\phantom{hel}
                           & \ShareNum{\PreIndOtherPTF}\phantom{hell}
                           & \ShareNum{\PreIndOtherPFZ}\phantom{hell}
                           & \ShareNum{\PreIndOtherPSF}\phantom{hell}\tabularnewline
   \noalign{\vskip 8pt}
           &               &\multicolumn{5}{c}{\textsc{Post-TCJA}} \tabularnewline
   \noalign{\vskip 2pt}
   \cline{3-7}
   \noalign{\vskip 6pt}
   \textsc{Technology}         & & \CountNum{\PostTechCounts}\phantom{he}
                           & \ShareNum{\PostTechAvg}\phantom{hel}
                           & \ShareNum{\PostTechPTF}\phantom{hell}
                           & \ShareNum{\PostTechPFZ}\phantom{hell}
                           & \ShareNum{\PostTechPSF}\phantom{hell}\tabularnewline
   \noalign{\vskip 3pt}
   \textsc{Pharmaceuticals}& & \CountNum{\PostPharmaCounts}\phantom{he}
                           & \ShareNum{\PostPharmaAvg}\phantom{hel}
                           & \ShareNum{\PostPharmaPTF}\phantom{hell}
                           & \ShareNum{\PostPharmaPFZ}\phantom{hell}
                           & \ShareNum{\PostPharmaPSF}\phantom{hell}\tabularnewline
   \noalign{\vskip 3pt}
   \textsc{Manufacturing}  & & \CountNum{\PostManufCounts}\phantom{he}
                           & \ShareNum{\PostManufAvg}\phantom{hel}
                           & \ShareNum{\PostManufPTF}\phantom{hell}
                           & \ShareNum{\PostManufPFZ}\phantom{hell}
                           & \ShareNum{\PostManufPSF}\phantom{hell}\tabularnewline
   \noalign{\vskip 3pt}
   \textsc{Oil \& gas}  & & \CountNum{\PostOilCounts}\phantom{he}
                           & \ShareNum{\PostOilAvg}\phantom{hel}
                           & \ShareNum{\PostOilPTF}\phantom{hell}
                           & \ShareNum{\PostOilPFZ}\phantom{hell}
                           & \ShareNum{\PostOilPSF}\phantom{hell}\tabularnewline
   \noalign{\vskip 3pt}
   \textsc{Retail}       & & \CountNum{\PostRetailCounts}\phantom{he}
                           & \ShareNum{\PostRetailAvg}\phantom{hel}
                           & \ShareNum{\PostRetailPTF}\phantom{hell}
                           & \ShareNum{\PostRetailPFZ}\phantom{hell}
                           & \ShareNum{\PostRetailPSF}\phantom{hell}\tabularnewline
   \noalign{\vskip 3pt}
   \textsc{Utilities}    & & \CountNum{\PostUtilCounts}\phantom{he}
                           & \ShareNum{\PostUtilAvg}\phantom{hel}
                           & \ShareNum{\PostUtilPTF}\phantom{hell}
                           & \ShareNum{\PostUtilPFZ}\phantom{hell}
                           & \ShareNum{\PostUtilPSF}\phantom{hell}\tabularnewline
   \noalign{\vskip 3pt}
   \textsc{Other}        & & \CountNum{\PostIndOtherCounts}\phantom{he}
                           & \ShareNum{\PostIndOtherAvg}\phantom{hel}
                           & \ShareNum{\PostIndOtherPTF}\phantom{hell}
                           & \ShareNum{\PostIndOtherPFZ}\phantom{hell}
                           & \ShareNum{\PostIndOtherPSF}\phantom{hell}\tabularnewline
   \noalign{\vskip 6pt}
   \hline
   \end{tabular}

\bigskip
\begin{tablenotes}[flushleft]
\item[] \footnotesize{\textit{Notes:} 
Entries report the distribution of the gap between the effective tax rate and the statutory tax rate, 
in percentage points, by industry. All statistics are computed on the full sample using
sales as weights. \emph{Counts} are the number of firm-year observations in each industry-period cell. The 
\emph{pre-TCJA}
period is 2013--2016 and the \emph{post-TCJA} period is 2019–-2024. 
Transition years 2017–2018 are omitted. Negative values
indicate effective tax rates below the statutory rate. 
Industries follow the SIC-based classification used elsewhere in the paper.}

\end{tablenotes}
\end{threeparttable}

\bigskip
\bigskip
\end{table}


\subsection{Causal Estimates of TCJA}

In this section, we exploit pre-TCJA heterogeneity in 
company income tax reconciliation and balance sheet items to estimate
causal effects of the post-TCJA outcomes. Of particular interest
is the policy response of profit shifting and real productive activities.

\subsubsection{Effective Tax Rate Pass-Through\label{sec:EffectiveTaxRate}} 


We first test the pass-through of TCJA policies to higher
effective tax rates. In Section \ref{sec:DescriptiveEvidence},
we provided descriptive evidence that TCJA compressed the
tax rate gap by reducing foreign tax advantages. Here, we
estimate the pass-through to effective tax rates. 

Our empirical strategy follows the standard difference-in-differences framework 
described in \citet{AngristPischke2009}. 
Our treatment intensity is each firm's average pre-TCJA international 
contribution $\bar x^{\text{pre}}_{\text{Intl},i}$ to the tax rate gap
in percentage points for company $i$ over period 2013--2016.
Firms with larger negative values 
had greater exposure to foreign tax advantages before TCJA. 
We compare these firms with less exposed firms after the reform.
More specifically, we use the following specification
\begin{equation}
  y_{it} =
  \alpha_i+\lambda_{j(i)t}
  +\beta
  \bigl(\bar x^{\text{pre}}_{\text{Intl},i}\times 1[t\in \text{post}]\bigr)
  +\varepsilon_{it}
\label{eq:PooledDiD}
\end{equation}
where $y_{it}$ is the outcome variable for $i$ in $t$, 
$\alpha_i$ is a firm fixed effect,
$\lambda_{j(i)t}$ are industry by year fixed effects, and 
``post'' is 2019--2022.
Transition years 2017 and 2018 are excluded, so the pooled
specification compares 2019--2022 to 2013--2016. Standard
errors are clustered at the firm level. 
The sample we will use when estimating $\beta$ in equation (\ref{eq:PooledDiD})
includes all companies observed at least twice before and twice after
the reform and, for the tax and profit-location outcomes, with positive
pretax income.\footnote{%
The regressions also end earlier than the descriptive tables of Section
\ref{sec:DescriptiveEvidence}, comparing 2019--2022 with 2013--2016
rather than running through 2024. We use the shorter window because it
is the more conservative one: extending the post-period to 2024 moves
the international coefficient from \CoefNum{\IntlCoef}\ to
\CoefNum{\IntlCoefLongWindow}.}

We run the regression in equation (\ref{eq:PooledDiD})
using different outcome variables ($y_{it}$). 
We first study tax outcomes before turning 
to profit shifting and real activity.
We start with the
international contribution for company $i$
in year $t$, which we denote by $x_{\text{Intl},it}$.
Since $\bar x^{\text{pre}}_{\text{Intl},i}$ is negative for 
companies doing business in low-tax foreign jurisdictions, we expect 
$\beta$ to be negative if the tax rate gap closes more for those 
with larger pre-TCJA foreign tax advantages.
We then explore the effective tax rate 
for firm $i$ in $t$, denoted here by ETR$_{it}$. Because the industry-by-year effects absorb the change common to an
industry, a negative coefficient $\beta$ in this case says that companies
with larger pre-TCJA foreign tax advantages had larger increases in their
effective tax rates.

Before reporting results for the difference-in-differences
regressions, we provide some checks on our empirical strategy.
Identification comes from comparisons of companies facing 
different pre-TCJA foreign tax advantages within the same industry.
Thus, we have included industry-year fixed effects.
In Table \ref{tab:INTLexposure}, we 
report on the distribution of $\bar x^{\text{pre}}_{\text{Intl},i}$ 
by industry for the sample of companies used to estimate $\beta$ in (\ref{eq:PooledDiD}).
We find that for most of the sample---especially companies in technology,
pharmaceuticals, and manufacturing---there is large within-industry variation in 
international exposure. For example, 
in the case of technology companies,
the standard deviation of 
$\bar x^{\text{pre}}_{\text{Intl},i}$ is 
\ShareNum{\TechSD}\ percentage points
and values at the 25$^{\rm th}$ and 75$^{\rm th}$ 
percentiles are
\ShareNum{\TechPTF}\  and
\ShareNum{\TechPSF}, respectively.
Because a firm-year with no international line is assigned a
contribution of zero, we also report the percentage of firms with a
nonzero value---the share reporting an international line at all.
For technology companies, we find
\ShareNum{\TechNonzero}\ percent do.
If we include all firms, the share of nonzeros is 
\ShareNum{\AllInSampleNonzero}\ percent, with all 
industries but utilities having a significant share of
companies with international exposure. 

A second check on the empirical strategy
is to run the following 
first-stage regression:
\begin{equation}
  y_{it} =
  \alpha_i+\lambda_{j(i)t}
  +\sum_{s\neq 2016} \beta_s
  \bigl(\bar x^{\text{pre}}_{\text{Intl},i}\times 1[t=s]\bigr)
  +\varepsilon_{it},
\label{eq:FSregression}
\end{equation}
with $\beta_{\rm 2016}$ normalized to zero.
We want to make sure that the estimates $\hat\beta$ are
close to zero before TCJA and significantly different from
zero after TCJA.
In Panels A and B of Figure \ref{fig:EventStudy},
we report results using the international
contribution and the effective tax rate as outcomes, 
respectively.\footnote{Panel C is discussed in the next section.}
In both panels, we find coefficients around zero prior to 2016
and significantly negative afterward. A negative $\hat\beta$ in Panel A
means that companies with a larger pre-TCJA international contribution 
had a larger move toward zero than their less-exposed peers.
A negative $\hat\beta$ in Panel B means the effective tax
rate faced by the more-exposed companies rises more after TCJA 
than for the less-exposed companies.
In 2017, there is a large decline in Panel B
because of the one-time transition tax assessed 
on foreign earnings that have not yet been repatriated 
(Section 965).
After 2018, we see that both sets of coefficients 
in Panels A and B stabilize to levels that provide information on the differential 
impacts for companies that were different  in their pre-TCJA exposure
to foreign tax advantages.






%
% TABLE 6: Distribution of Pre-TCJA International Exposure by Industry
%

\begin{table}[t!]
\centering
\begin{threeparttable}
\caption{\textsc{Distribution of Pre-TCJA International Exposure by Industry}}
\label{tab:INTLexposure}

   \begin{tabular}{lrrrrrrr}
   \hline
   \hline
   \noalign{\vskip 5pt}
                           &
                           &
                           &
                           &\multicolumn{3}{c}{\textsc{Percentiles}} \tabularnewline
   \noalign{\vskip 2pt}
   \cline{5-7}
   \noalign{\vskip -4pt}
                           & & & \textsc{Std.\,} & & & & \textsc{\%\ \ \ \ \ \ }\tabularnewline
   \noalign{\vskip -2pt}
   \textsc{Industry}       &\multicolumn{1}{c}{\textsc{Counts}}
                           &\multicolumn{1}{c}{\textsc{Mean}}
                           &\multicolumn{1}{c}{\textsc{Dev.}}
                           &\multicolumn{1}{c}{\textsc{25$^{\rm th}$}}
                           &\multicolumn{1}{c}{\textsc{50$^{\rm th}$}}
                           &\multicolumn{1}{c}{\textsc{75$^{\rm th}$}}
                           &\multicolumn{1}{c}{\textsc{Nonzero}}\tabularnewline
   \noalign{\vskip 5pt}
   \hline
   \noalign{\vskip 6pt}
   \textsc{Technology}     & \CountNum{\TechN}\phantom{he}
                           & \ShareNum{\TechMean}\phantom{hel}
                           & \ShareNum{\TechSD}\phantom{hel}
                           & \ShareNum{\TechPTF}\phantom{hell}
                           & \ShareNum{\TechPFZ}\phantom{hell}
                           & \ShareNum{\TechPSF}\phantom{he}
                           & \ShareNum{\TechNonzero}\phantom{hell}\tabularnewline
   \noalign{\vskip 3pt}
   \textsc{Pharmaceuticals}& \CountNum{\PharmaN}\phantom{he}
                           & \ShareNum{\PharmaMean}\phantom{hel}
                           & \ShareNum{\PharmaSD}\phantom{hel}
                           & \ShareNum{\PharmaPTF}\phantom{hell}
                           & \ShareNum{\PharmaPFZ}\phantom{hell}
                           & \ShareNum{\PharmaPSF}\phantom{he}
                           & \ShareNum{\PharmaNonzero}\phantom{hell}\tabularnewline
   \noalign{\vskip 3pt}
   \textsc{Manufacturing}  & \CountNum{\ManufN}\phantom{he}
                           & \ShareNum{\ManufMean}\phantom{hel}
                           & \ShareNum{\ManufSD}\phantom{hel}
                           & \ShareNum{\ManufPTF}\phantom{hell}
                           & \ShareNum{\ManufPFZ}\phantom{hell}
                           & \ShareNum{\ManufPSF}\phantom{he}
                           & \ShareNum{\ManufNonzero}\phantom{hell}\tabularnewline
   \noalign{\vskip 3pt}
   \textsc{Oil \& gas}     & \CountNum{\OilN}\phantom{he}
                           & \ShareNum{\OilMean}\phantom{hel}
                           & \ShareNum{\OilSD}\phantom{hel}
                           & \ShareNum{\OilPTF}\phantom{hell}
                           & \ShareNum{\OilPFZ}\phantom{hell}
                           & \ShareNum{\OilPSF}\phantom{he}
                           & \ShareNum{\OilNonzero}\phantom{hell}\tabularnewline
   \noalign{\vskip 3pt}
   \textsc{Retail}         & \CountNum{\RetailN}\phantom{he}
                           & \ShareNum{\RetailMean}\phantom{hel}
                           & \ShareNum{\RetailSD}\phantom{hel}
                           & \ShareNum{\RetailPTF}\phantom{hell}
                           & \ShareNum{\RetailPFZ}\phantom{hell}
                           & \ShareNum{\RetailPSF}\phantom{he}
                           & \ShareNum{\RetailNonzero}\phantom{hell}\tabularnewline
   \noalign{\vskip 3pt}
   \textsc{Utilities}      & \CountNum{\UtilitiesN}\phantom{he}
                           & \ShareNum{\UtilitiesMean}\phantom{hel}
                           & \ShareNum{\UtilitiesSD}\phantom{hel}
                           & \ShareNum{\UtilitiesPTF}\phantom{hell}
                           & \ShareNum{\UtilitiesPFZ}\phantom{hell}
                           & \ShareNum{\UtilitiesPSF}\phantom{he}
                           & \ShareNum{\UtilitiesNonzero}\phantom{hell}\tabularnewline
   \noalign{\vskip 3pt}
   \textsc{Other}          & \CountNum{\OtherN}\phantom{he}
                           & \ShareNum{\OtherMean}\phantom{hel}
                           & \ShareNum{\OtherSD}\phantom{hel}
                           & \ShareNum{\OtherPTF}\phantom{hell}
                           & \ShareNum{\OtherPFZ}\phantom{hell}
                           & \ShareNum{\OtherPSF}\phantom{he}
                           & \ShareNum{\OtherNonzero}\phantom{hell}\tabularnewline
   \noalign{\vskip 3pt}
   \cline{2-8}
   \noalign{\vskip 3pt}
   \textsc{All in sample}  & \CountNum{\AllInSampleN}\phantom{he}
                           & \ShareNum{\AllInSampleMean}\phantom{hel}
                           & \ShareNum{\AllInSampleSD}\phantom{hel}
                           & \ShareNum{\AllInSamplePTF}\phantom{hell}
                           & \ShareNum{\AllInSamplePFZ}\phantom{hell}
                           & \ShareNum{\AllInSamplePSF}\phantom{he}
                           & \ShareNum{\AllInSampleNonzero}\phantom{hell}\tabularnewline
   \noalign{\vskip 6pt}
   \hline
   \end{tabular}

\bigskip
\begin{tablenotes}[flushleft]
\item[] \footnotesize{\textit{Notes:} Entries report the distribution of firms' pre-TCJA international
  reconciliation exposure, $\bar x^{\text{pre}}_{\text{Intl},i}$, by industry. For each firm,
  $\bar x^{\text{pre}}_{\text{Intl},i}$ is the average international contribution over
  2013--2016. The firms in the sample must have positive pretax income, and have at 
  least two pre-TCJA and two post-TCJA observations. The final column reports the share of
  firms with nonzero pre-TCJA international exposure. Industries follow the SIC-based
  classification used elsewhere in the paper.}
\end{tablenotes}
\end{threeparttable}

\bigskip
\bigskip
\end{table}










%
% FIGURE 3: Event Study Responses to Pre-TCJA International Exposure
%

\pgfplotstableread[col sep=comma]{\Inputs/first_stage_tikz/first_stage_pre.csv}\fstpre
\pgfplotstableread[col sep=comma]{\Inputs/first_stage_tikz/first_stage_shock.csv}\fstshock
\pgfplotstableread[col sep=comma]{\Inputs/first_stage_tikz/first_stage_post.csv}\fstpost
\pgfplotstableread[col sep=comma]{\Inputs/first_stage_tikz/first_stage_reference.csv}\fstref
\pgfplotstableread[col sep=comma]{\Inputs/first_etr_tikz/first_etr_pre.csv}\etrpre
\pgfplotstableread[col sep=comma]{\Inputs/first_etr_tikz/first_etr_shock.csv}\etrshock
\pgfplotstableread[col sep=comma]{\Inputs/first_etr_tikz/first_etr_post.csv}\etrpost
\pgfplotstableread[col sep=comma]{\Inputs/first_etr_tikz/first_etr_reference.csv}\etrref
\pgfplotstableread[col sep=comma]{\Inputs/first_pishr_tikz/first_pishr_pre.csv}\pspre
\pgfplotstableread[col sep=comma]{\Inputs/first_pishr_tikz/first_pishr_shock.csv}\psshock
\pgfplotstableread[col sep=comma]{\Inputs/first_pishr_tikz/first_pishr_post.csv}\pspost
\pgfplotstableread[col sep=comma]{\Inputs/first_pishr_tikz/first_pishr_reference.csv}\psref

\definecolor{ESPre}{RGB}{120,120,120}
\definecolor{ESShock}{RGB}{217,95,2}
\definecolor{ESPost}{RGB}{31,119,180}

\pgfplotsset{
  espanel/.style={
    width=0.31\textwidth,
    height=0.26\textwidth,
    xmin=2012.7, xmax=2022.3,
    xtick={2013,2015,2017,2019,2021},
    axis x line=bottom,
    axis y line=left,
    tick align=outside,
    grid=none,
    title style={font=\small},
    label style={font=\scriptsize},
    ticklabel style={font=\scriptsize},
  },
  espre/.style={
    only marks,
    color=ESPre,
    mark=o,
    mark size=2pt,
    mark options={draw=ESPre, fill=white, line width=0.6pt},
    error bars/y dir=both,
    error bars/y explicit,
    error bars/error bar style={draw=ESPre, line width=0.7pt},
  },
  espost/.style={
    only marks,
    color=ESPost,
    mark=*,
    mark size=2.2pt,
    mark options={draw=ESPost, fill=ESPost, line width=0.6pt},
    error bars/y dir=both,
    error bars/y explicit,
    error bars/error bar style={draw=ESPost, line width=0.7pt},
  },
  esshock/.style={
    only marks,
    color=ESShock,
    mark=square*,
    mark size=2.4pt,
    mark options={draw=ESShock, fill=ESShock, line width=0.6pt},
    error bars/y dir=both,
    error bars/y explicit,
    error bars/error bar style={draw=ESShock, line width=0.7pt},
  }
}

\begin{figure}[!tp]
\centering
\caption{\textsc{Event Study Responses to Pre-TCJA International Exposure}}
\label{fig:EventStudy}

  \begin{tikzpicture}
  \begin{groupplot}[
    group style={
      group size=1 by 3,
      vertical sep=1.8cm,
    },
    width=0.85\textwidth,
    height=0.32\textwidth,
    xmin=2012.7, xmax=2022.3,
    xtick={2013,2016,2019,2022},
    xticklabel={\pgfmathprintnumber[fixed,precision=0,1000 sep={}]{\tick}},
    axis x line=bottom,
    axis y line=left,
    tick align=outside,
    grid=none,
    title style={font=\small},
    label style={font=\scriptsize},
    ticklabel style={font=\scriptsize},
  ]

  \nextgroupplot[
    title={\textsc{A. International Contribution}},
    ymin=-1.05, ymax=0.25,
    ylabel={\textsc{Coefficient}},
    ylabel style={
        at={(axis description cs:-0.12,.5)},
        anchor=center},
  ]
  \addplot+[espre]
    table[x=year,y=beta,y error=err95] {\fstpre};
  \addplot+[esshock]
    table[x=year,y=beta,y error=err95] {\fstshock};

  \addplot+[espost]
    table[x=year,y=beta,y error=err95] {\fstpost};

  \addplot+[
    only marks, mark=diamond*, mark size=2.2pt, draw=black, fill=white,
  ] table[x=year,y=beta] {\fstref};

  \addplot[black, thin] coordinates {(2012.7,0) (2022.3,0)};
  \addplot[black!60, dashed] coordinates {(2017.5,-1.05) (2017.5,0.25)};



  \nextgroupplot[
    title={\textsc{B. Effective Tax Rate}},
    ymin=-2.6, ymax=0.6,
    ylabel={\textsc{Coefficient}},
    ylabel style={
        at={(axis description cs:-0.12,.5)},
        anchor=center},
  ]
  \addplot+[espre]
    table[x=year,y=beta,y error=err95] {\etrpre};
  \addplot+[esshock]
    table[x=year,y=beta,y error=err95] {\etrshock};
  \addplot+[espost]
    table[x=year,y=beta,y error=err95] {\etrpost};
  \addplot+[only marks, mark=diamond*, mark size=2.2pt, draw=black, fill=white]
    table[x=year,y=beta] {\etrref};
  \addplot[black, thin] coordinates {(2012.7,0) (2022.3,0)};
  \addplot[black!60, dashed] coordinates {(2017.5,-2.6) (2017.5,0.6)};

  \nextgroupplot[
    title={\textsc{C. Foreign Profit Share}},
    ymin=-0.5, ymax=1.5,
    scaled y ticks=false,
    ylabel={\textsc{Coefficient}},
    ylabel style={
        at={(axis description cs:-0.12,.5)},
        anchor=center},
  ]
  \addplot+[espre]
    table[x=year,y expr={100*\thisrow{beta}},y error expr={100*\thisrow{err95}}] {\pspre};
  \addplot+[esshock]
    table[x=year,y expr={100*\thisrow{beta}},y error expr={100*\thisrow{err95}}] {\psshock};
  \addplot+[espost]
    table[x=year,y expr={100*\thisrow{beta}},y error expr={100*\thisrow{err95}}] {\pspost};
  \addplot+[only marks, mark=diamond*, mark size=2.2pt, draw=black, fill=white]
    table[x=year,y expr={100*\thisrow{beta}},y error expr={100*\thisrow{err95}}] {\psref};
  \addplot[black, thin] coordinates {(2012.7,0) (2022.3,0)};
  \addplot[black!60, dashed] coordinates {(2017.5,-0.5) (2017.5,1.5)};

  \end{groupplot}
  \end{tikzpicture}

  \bigskip
  \begin{minipage}{0.95\textwidth}
  \footnotesize
\textit{Notes:} Each panel reports year-by-year coefficients from event-study regressions of the
  indicated outcome on firms' pre-TCJA International exposure, $\bar x^{\text{pre}}
  _{\text{Intl},i}$, interacted with calendar-year indicators. The omitted reference year is 2016.
  Panel A uses the international reconciliation contribution, measured in percentage points; Panel B
  uses the effective tax rate, also in percentage points; and Panel C uses foreign profit share,
  reported in percentage points. All specifications include firm and industry-by-year fixed effects,
  and standard errors are clustered at the firm level. The treatment variable $\bar x^{\text{pre}}
  _{\text{Intl},i}$ is each firm's average international contribution over 2013--2016. More negative
  values correspond to larger pre-TCJA International tax benefits. Accordingly, negative
  coefficients in Panel A indicate greater post-TCJA closure of the international contribution, negative
  coefficients in Panel B indicate larger post-TCJA increases in effective tax rates, and positive
  coefficients in Panel C indicate larger post-TCJA declines in foreign profit share among more
  exposed firms. The dashed vertical line marks 2017, the year of TCJA enactment. The number of
  firm-year observations for Panels A, B, and C are \CountNum{\IntlEventObs},
  \CountNum{\EtrEventObs}, and \CountNum{\PiShrEventObs}, respectively; these
  exceed the counts in Table \ref{tab:PooledDiD} because the event study also
  uses the 2017 and 2018 transition years.
  \end{minipage}

  \bigskip
  \bigskip

\end{figure}



We turn now to the main results for the pooled difference-in-differences
regression in (\ref{eq:PooledDiD}), which are shown in Table \ref{tab:PooledDiD}.
Results for the international contribution and the effective tax rate
are shown in the first two rows. The coefficient in the first column is
$\hat\beta$ in the regression, with standard errors reported in parentheses.
The last two columns provide information on the counts of 
firms and total observations
for our samples.
Throughout, we write $\hat\beta_{\text{Intl}}$,
$\hat\beta_{\text{ETR}}$, and, below, $\hat\beta_{\text{Share}}$
for the estimates of $\beta$ in (\ref{eq:PooledDiD}) when the outcome
$y_{it}$ is the international contribution, the effective tax rate, and
the foreign profit share, respectively.
Both are ratios of percentage points to percentage points and so are
dimensionless. Our estimate of $\hat\beta_{\text{Intl}}$ is
\CoefNum{\IntlCoef}: for each additional percentage point of pre-TCJA
international exposure, the international contribution closes by
\CoefNum{\IntlCoefAbs}\ percentage points after TCJA. Our estimate of
$\hat\beta_{\text{ETR}}$ is \CoefNum{\EtrCoef}, so the same firm's
effective tax rate rises by \CoefNum{\EtrCoefAbs}\ percentage points.
The ratio of these coefficients $\hat\beta_{\rm ETR}/\hat \beta_{\rm Intl}$---which
is roughly \RoundNum{\PassthroughEtrPct} percent---is the pass-through
to higher effective rates.

To put these estimates in context, consider the impact of TCJA on effective tax
rates in the technology sector. 
The average international contribution in the income tax reconciliation before TCJA was 
\ShareNum{\TechIntlPre}\ percentage points. After TCJA, this fell to
\ShareNum{\TechIntlPost}, implying a closure of 
\ShareNum{\TechIntlClosure}.  We can see this change in the first row of Figure 
\ref{fig:Industry}, Panel B. 
Using our estimates from Table \ref{tab:PooledDiD}, we compute an
implied increase in the effective tax rate of
\ShareNum{\TechImpliedEtr}\ percentage points
(or, $\hat\beta_{\text{ETR}}/\hat\beta_{\text{Intl}} \times$ \ShareNum{\TechIntlClosure}).
This is smaller than the closure itself because part of it is returned
through the other reconciliation categories. Appendix Table \ref{tab:ChannelDiD}
estimates equation (\ref{eq:PooledDiD}) category by category. The
coefficients sum to \CoefNum{\ChSumCoef}, against
\CoefNum{\EtrCoef}\ for the effective tax rate itself, so the
decomposition is close to exact. Outside international the coefficients
sum to \CoefNum{\ChNonIntlCoef}, or
\RoundNum{\ChOffsetPctCoef}\ percent of the international closure, and the
two categories that carry it are valuation allowances and the residual
category. For technology that is roughly
\ShareNum{\TechOtherChannels}\ of the
\ShareNum{\TechIntlClosure}\ point closure. Companies did find some
other ways to offset the loss in foreign tax benefits, but they
recovered only about a quarter of it.


%
% TABLE 7: Pooled Difference-in-Differences Estimates of TCJA Pass-Through}}
%

\begin{table}[!tp]
\centering
\begin{threeparttable}
\caption{\textsc{Pooled Difference-in-Differences Estimates of TCJA Pass-Through}}
\label{tab:PooledDiD}

   \begin{tabular}{clcrrrc}
   \hline
   \hline
   \noalign{\vskip 5pt}
   \phantom{h} & \textsc{Outcome} &\phantom{hello}
                    &\multicolumn{1}{c}{\textsc{Coefficient}}
                    &\multicolumn{1}{c}{\textsc{Firms}}
                    &\multicolumn{1}{c}{\textsc{Firm-years}} &\phantom{h}\tabularnewline
   \noalign{\vskip 5pt}
   \hline
   \noalign{\vskip 6pt}
   &\textsc{International contribution}& & \CoefNum{\IntlCoef}\StarBox{\IntlStars}
                                        & \CountNum{\IntlN}%\phantom{he}
                                        & \CountNum{\IntlObs}\phantom{hell} &\tabularnewline
                                    & & & \SENum{\IntlSe}\StarBox{} & & & \tabularnewline
    \noalign{\vskip 3pt}
   &\textsc{Effective tax rate}       & & \CoefNum{\EtrCoef}\StarBox{\EtrStars}
                                        & \CountNum{\EtrN}%\phantom{he}
                                        & \CountNum{\EtrObs}\phantom{hell} &\tabularnewline
                                    & & & \SENum{\EtrSe}\StarBox{} & & & \tabularnewline
    \noalign{\vskip 3pt}
   &\textsc{Foreign profit share}     & & \CoefNum{\PiShrCoef}\StarBox{\PiShrStars}
                                        & \CountNum{\PiShrN}%\phantom{he}
                                        & \CountNum{\PiShrObs}\phantom{hell} &\tabularnewline
                                    &  & & \SENum{\PiShrSe}\StarBox{} & & &\tabularnewline
    \noalign{\vskip 3pt}
   &\textsc{Foreign sales share}   & & \CoefNum{\GeoRevenueBetaForeignSharePp}\StarBox{}
                                        & \CountNum{\GeoRevenueShareFirms}%\phantom{he}
                                        & \CountNum{\GeoRevenueShareFirmYears}\phantom{hell} &\tabularnewline
                                    &  & & \SENum{\GeoRevenueBetaForeignShareSePp}\StarBox{} & & &\tabularnewline
   \noalign{\vskip 6pt}
   \hline
   \end{tabular}

\bigskip
\begin{tablenotes}[flushleft]
\item[] \footnotesize{\textit{Notes:} Entries report pooled difference-in-differences estimates 
relating post-TCJA outcomes to firms’ pre-TCJA International exposure. %\bar x^{\text{pre}}_{\text{Intl},i}$.
  The international contribution, effective tax rate, and foreign profit share are all reported 
in percentage points. The pooled specification compares 2019–2022 to 2013–2016, omitting 2017–2018, 
and includes firm and industry-by-year fixed effects.  Standard errors clustered at the firm level 
are reported in parentheses. $^{*} p<0.10$, $^{**} p<0.05$, $^{***} p<0.01$.}
\end{tablenotes}
\end{threeparttable}

\bigskip
\bigskip
\end{table}





\subsubsection{Profit Shifting\label{sec:ProfitShifting}} 

These results prompt an obvious follow-up question:  was there any impact on the geography of 
profits and sales or, more generally, on real productive activity?
To answer this, we use the same empirical strategy as above, but now
apply it to different outcomes, namely, foreign profit share,
foreign sales share, capital investment, research and development,
employment growth, and  sales growth.


We start by investigating the location of booked profits, which was a central focus of the
TCJA reform.  We estimate the change in the foreign profit
share---that is, the ratio of foreign profits to total profits as a percent---in response
to TCJA-induced changes in effective tax rates.  As a first-stage check, we rerun
the regression in (\ref{eq:FSregression}) to make sure that the identification 
strategy is sound.  In Panel C of Figure \ref{fig:EventStudy},
we report the first-stage regression results with $y_{it}$ given
by the foreign profit share.  Here, again, we find that the 
coefficients prior to 2016 are not significantly different from zero,
while the coefficients after 2018 are. Given this, we proceed with the pooled
difference-in-differences regression in (\ref{eq:PooledDiD}).
In Table \ref{tab:PooledDiD}, we report our estimate in the third
row, namely $\hat\beta_{\rm Share}$ equal to \CoefNum{\PiShrCoef}.
Dividing this by $\hat\beta_{\rm Intl}$ re-estimated on the same
\CountNum{\PiShrN}\ firms gives \CoefNum{\ProfitPerPpClosure}, which implies that
a firm's foreign profit share fell by about that many percentage points
for each percentage point of closure in its international
contribution.

The last step is to construct an interpretable estimate
of the impact of profit-shifting in response to tax changes using the coefficient
for the regression with the effective tax rate as the outcome variable.
To be consistent, we need to use the same sample of companies that have nonmissing
values for foreign profit (that is, variable, pifo in \citet{CompustatAnnual}).
When we do, we find the coefficient in the effective tax rate regression is 
\CoefNum{\IncomeLocBetaEtrSubsample}.
This means that a 1 percentage point TCJA-induced effective tax rate increase is 
associated with a 
\RatioNum{\IncomeLocClosureRatio}\ percentage point decline in the foreign profit share. 
To put this into perspective, the average pre-TCJA exposure among firms with
a negative international contribution is
\ShareNum{\HetExpExposure}\ percentage points, and those firms held
\ShareNum{\HetExpShare}\ percent of their pretax income abroad. The
estimate implies a fall of \ShareNum{\HetExpDrop}\ percentage points in the foreign share
for such a firm, or about a tenth of the pre-TCJA share.

To put these estimates in context, we can compare our estimation strategy
and results to earlier work by \citet{Grubert2012}.
\citet{Grubert2012} used tax filings of  754 large multinational companies
between 1996 and 2004 and  studied the responsiveness of foreign profit location 
to tax incentives using variation in firms' average effective foreign tax rates. 
We instead exploit a large tax reform and compare firms with different
pre-reform exposure to foreign tax advantages. 
Our measure of tax incentives is derived from firms' tax reconciliation disclosures 
rather than tax return data,  and the variation we 
study primarily reflects changes in the tax advantage associated 
with foreign operations because the international contribution is so large.
The TCJA therefore provides a more direct source of 
identifying variation while retaining a close connection to the tax incentives 
examined by Grubert.  

In terms of headline results, \citet{Grubert2012} estimates the 
change in the foreign profit share with a change in the foreign effective tax rate
and finds it is $-$0.7  for his preferred empirical specification.
Our estimate of $-$\RatioNum{\IncomeLocClosureRatio}\ is the change in the
foreign profit share associated with 
a change in the total effective tax rate induced by the policy reform.
Since we have already established that most of the induced change occurs
through lower foreign tax advantages, the two ratios are conceptually close.
Our larger response is perhaps unsurprising given that our identification 
comes from a major tax reform that substantially reduced incentives for locating 
profits abroad, whereas Grubert relies on variation across firms in
effective foreign tax rate changes. Furthermore, our treatment intensity is based on firms' 
pre-reform exposure to foreign tax advantages---and thus our empirical
strategy exploits variation across firms with different opportunities to
engage in profit shifting.

In his study of the multinational  tax filings, \citet{Grubert2012} was also able to 
track the location of sales.\footnote{%
To construct data on foreign sales, we use the geographic segments
in \citet{CompustatSegments} and keep all firms in our sample with nonmissing data.}
Here, we investigate the impact of TCJA on location of
sales using the same empirical strategy as for profits. Specifically,
we run the first-stage regression in (\ref{eq:FSregression}) and
then the pooled difference-in-differences regression in (\ref{eq:PooledDiD}).
The main finding, like Grubert's famous title, is that profits---not sales---are
being globalized. 
In the first-stage regression, 
coefficients in all years are insignificantly different from zero.
In the pooled difference-in-differences regression, which is shown in the
fourth row of Table \ref{tab:PooledDiD}, the coefficient is also insignificantly
different from zero---the value 
is \CoefNum{\GeoRevenueBetaForeignSharePp} 
with a standard error of
\CoefNum{\GeoRevenueBetaForeignShareSePp}.
This stands in contrast to the foreign profit share, where the
corresponding coefficient is \CoefNum{\PiShrCoef}\ and precisely
estimated.

To investigate this finding further, we consider a variant
of our difference-in-differences specification as follows:
\begin{equation}
  y_{it} =
  \alpha_i+\lambda_{j(i)t}
  +\beta D_{it} 
  +\delta (Z_i \times 1[t\in \text{post}])
  +\gamma (D_{it} \times Z_i)
  +\varepsilon_{it}
\label{eq:DiDvariant}
\end{equation}
where $D_{it}= \bar x^{\text{pre}}_{\text{Intl},i}\times 1[t\in \text{post}]$
and $Z_i$ equal to a characteristic of company $i$
with a continuous value that is winsorized and standardized.
The idea is to improve our prediction of the profit-location responses
using pre-TCJA firm characteristics.  Obvious candidates based on earlier
work are R\&D intensities, measures of intangibles or intellectual property
(if available), use of tax haven subsidiaries, foreign scope (for example, 
counts of jurisdictions or subsidiaries), and balance sheet items 
(for example, cash or debt).

With the specification in equation (\ref{eq:DiDvariant}), we find
that the profit-location response is concentrated in companies that are
R\&D-intensive and liquid. Appendix Table \ref{tab:Heterogeneity} reports
$\hat\gamma$ one characteristic at a time. A one standard deviation
higher pre-TCJA R\&D-to-assets ratio raises the response by
\CoefNum{\HetRd}\ percentage points of the foreign profit share
with a standard error of \CoefNum{\HetRdSe}, and a one standard
deviation higher cash-to-assets ratio raises it by
\CoefNum{\HetCash}\ \SENum{\HetCashSe}. Our measures of the use of tax havens, on the
other hand, do not predict the response: the coefficient on having any
haven or conduit subsidiary is \CoefNum{\HetHaven}\
\SENum{\HetHavenSe}, and on the count of such subsidiaries is
\CoefNum{\HetHavenCount}\ \SENum{\HetHavenCountSe}. 


\subsubsection{Investment and Growth\label{sec:InvestmentGrowth}}

We next investigate if companies that lose their foreign tax advantages 
change the scale or location of their real activity once
reported income shifts back to the United States.
More specifically, we rerun our baseline first-stage and pooled
difference-in-differences regressions with $y_{it}$ given by one of 
seven outcomes. The real-activity outcomes are
capital expenditures, R\&D expenditures,
sales, and employment. The profit-location outcomes are pretax
income, foreign pretax income, and domestic pretax income.\footnote{%
The \citet{CompustatAnnual} variables used are capx, xrd, sale, emp, pi, pifo, and pidom.}

In the first-stage regression results (not shown here),
we find no significant change in capital expenditures, R\&D,
sales, or employment for more internationally exposed firms in any year.
In Table \ref{tab:PooledDiDRA} we report the coefficients and
firm and firm-year counts for the pooled difference-in-differences
regression. Here, again, we find no response in any of the four real-activity outcomes.
The lower panel shows where the response does occur. Total pretax
income is unchanged---the coefficient is
\CoefNum{\RasPiPft}\ with a standard error of
\CoefNum{\RasPiPftSe}---while the coefficients on foreign and
domestic pretax income are \CoefNum{\RasPiFoPft}\ and
\CoefNum{\RasPiDomPft}, so that foreign profits fall and domestic
profits rise for more exposed firms. More
exposed firms did not become smaller or invest less but instead booked the
same profits in different places.
If we check firm balance sheets, we find the contraction in cash
is tax driven with firms paying a transition tax on earnings not yet 
repatriated when they switched from the world-wide tax system
to a territorial system.  
If we rerun regression (\ref{eq:PooledDiD}) with balance sheet
items divided by lagged assets, we find 
the coefficient on cash is
\CoefNum{\BsCash}\ \SENum{\BsCashSe}\ and the coefficient on cash
taxes paid is \CoefNum{\BsTaxPaid}\ \SENum{\BsTaxPaidSe}. In other words, cash falls
and taxes paid rise for more exposed firms. Dividends increase modestly,
\CoefNum{\BsDiv}\ \SENum{\BsDivSe}. There is no response in share
repurchases, \CoefNum{\BsBuyback}\ \SENum{\BsBuybackSe}, in debt,
\CoefNum{\BsDebt}\ \SENum{\BsDebtSe}, or in acquisitions,
\CoefNum{\BsAcq}\ \SENum{\BsAcqSe}.

Because annual 10-K filings report consolidated worldwide operations, 
they may mask changes in the geographic location of real activity.
To investigate this, we use the \citet{CompustatSegments}
database to track domestic and foreign property, plant, and equipment.  
We construct foreign shares of capital expenditures  and capital stock 
and rerun the first-stage regression in equation (\ref{eq:FSregression})
and pooled difference-in-differences regression in equation (\ref{eq:PooledDiD})
(with Compustat variables capxs and ppents).\footnote{%
Geographic capital is reported by far fewer filers than the
consolidated series, so these regressions run on
\CountNum{\CapLocCapxFirms}\ firms for expenditures and
\CountNum{\CapLocPpeFirms}\ for stocks.}
In the case of capital
expenditures, we cannot rule out differential trends prior to TCJA. The
pre-TCJA event-study coefficients are negative and the post-TCJA ones,
though positive, are individually insignificant, so the pooled estimate
of \CoefNum{\CapLocCapx}\ \SENum{\CapLocCapxSe}\ percentage
points rests in part on that pre-trend.
In the case of capital stocks, the results are much clearer: there
is no evidence of any shift in the location of the installed capital
stock. The pooled coefficient is \CoefNum{\CapLocPpe}\
\SENum{\CapLocPpeSe}, and no year is significant.

One caveat applies to every regression in this subsection. The
industry-by-year effects absorb whatever TCJA did to all companies in
an industry, so these coefficients measure how much \emph{more} the
exposed companies invested, hired, or sold, not whether investment rose
after the reform. The reform's largest changes to the incentive to
invest---the cut in the statutory rate and full expensing---moved the
user cost of capital for every company, and our design is silent on this.
Our results therefore do not speak directly to the debate between
\citet{CSZZ2024} and \citet{GravelleMarples2025}, who use variation in
the user cost of capital.

Overall, the main takeaway from the difference-in-differences regression
using variation in international exposure is that the loss in foreign tax
advantages resulted in a significant profit relocation for more R\&D-intensive and
liquid companies, with little differential response in real activities.


%
% TABLE 8: Pooled Difference-in-Differences Estimates of Real Activity Impacts}}
%

\begin{table}[!tp]
\centering
\begin{threeparttable}
\caption{\textsc{Pooled Difference-in-Differences Estimates of \\ Real Activity and Profit Location}}
\label{tab:PooledDiDRA}

   \begin{tabular}{clcrrrc}
   \hline
   \hline
   \noalign{\vskip 5pt}
   \phantom{h} & \textsc{Outcome} &\phantom{hel}
                    &\multicolumn{1}{c}{\textsc{Coefficient}}
                    &\multicolumn{1}{c}{\textsc{Firms}}
                    &\multicolumn{1}{c}{\textsc{Firm-years}} &\phantom{h}\tabularnewline
   \noalign{\vskip 5pt}
   \hline
   \noalign{\vskip 6pt}
   & \textsc{A. Real activity} & & & & &\tabularnewline
   \noalign{\vskip 8pt}
   &\ \ \ \ \ {Capital expenditures}& & \CoefNum{\RasCapx}\StarBox{\RasCapxStars}
                                        & \CountNum{\RasCapxFirms}%
                                        & \CountNum{\RasCapxObs}\phantom{hell} &\tabularnewline
                                    & & & \SENum{\RasCapxSe}\StarBox{} & & & \tabularnewline
    \noalign{\vskip 3pt}
   &\ \ \ \ \ {R\&D expenditures}& & \CoefNum{\RasRd}\StarBox{\RasRdStars}
                                        & \CountNum{\RasRdFirms}%
                                        & \CountNum{\RasRdObs}\phantom{hell} &\tabularnewline
                                    & & & \SENum{\RasRdSe}\StarBox{} & & & \tabularnewline
    \noalign{\vskip 3pt}
   &\ \ \ \ \ {Sales}& & \CoefNum{\RasSales}\StarBox{\RasSalesStars}
                                        & \CountNum{\RasSalesFirms}%
                                        & \CountNum{\RasSalesObs}\phantom{hell} &\tabularnewline
                                    & & & \SENum{\RasSalesSe}\StarBox{} & & & \tabularnewline
    \noalign{\vskip 3pt}
   &\ \ \ \ \ {Employment}& & \CoefNum{\RasEmp}\StarBox{\RasEmpStars}
                                        & \CountNum{\RasEmpFirms}%
                                        & \CountNum{\RasEmpObs}\phantom{hell} &\tabularnewline
                                    & & & \SENum{\RasEmpSe}\StarBox{} & & & \tabularnewline
   \noalign{\vskip 10pt}
   & \textsc{B. Location of profits} & & & & &\tabularnewline
   \noalign{\vskip 8pt}
   &\ \ \ \ \ {Pretax income}& & \CoefNum{\RasPiPft}\StarBox{\RasPiPftStars}
                                        & \CountNum{\RasPiPftFirms}%
                                        & \CountNum{\RasPiPftObs}\phantom{hell} &\tabularnewline
                                    & & & \SENum{\RasPiPftSe}\StarBox{} & & & \tabularnewline
    \noalign{\vskip 3pt}
   &\ \ \ \ \ {Foreign pretax income}& & \CoefNum{\RasPiFoPft}\StarBox{\RasPiFoPftStars}
                                        & \CountNum{\RasPiFoPftFirms}%
                                        & \CountNum{\RasPiFoPftObs}\phantom{hell} &\tabularnewline
                                    & & & \SENum{\RasPiFoPftSe}\StarBox{} & & & \tabularnewline
    \noalign{\vskip 3pt}
   &\ \ \ \ \ {Domestic pretax income}& & \CoefNum{\RasPiDomPft}\StarBox{\RasPiDomPftStars}
                                        & \CountNum{\RasPiDomPftFirms}%
                                        & \CountNum{\RasPiDomPftObs}\phantom{hell} &\tabularnewline
                                    & & & \SENum{\RasPiDomPftSe}\StarBox{} & & & \tabularnewline
    \noalign{\vskip 3pt}
   \noalign{\vskip 6pt}
   \hline
   \end{tabular}

\bigskip
\begin{tablenotes}[flushleft]
\item[] \footnotesize{\textit{Notes:} Entries report pooled difference-in-differences estimates
  relating post-TCJA outcomes to firms' pre-TCJA International exposure,
  $\bar x^{\text{pre}}_{\text{Intl},i}$. The pooled specification compares
  2019--2022 to 2013--2016, omitting 2017--2018, and includes firm and
  industry-by-year fixed effects. Standard errors clustered at the firm level
  are reported in parentheses. All seven displayed dependent variables are logs
  of levels, so coefficients are in $100\times$ log points per percentage point
  of exposure. Because $\bar x^{\text{pre}}_{\text{Intl},i}$ is more negative
  for firms with larger pre-TCJA International tax benefits, a negative
  coefficient indicates a larger post-TCJA increase in the outcome among more
  exposed firms, while a positive coefficient indicates a relative decrease.
  Samples differ across rows because the outcomes have different coverage. The
  profit-location rows use the positive-profit sample for which the relevant
  logged profit outcome is defined. $^{*} p<0.10$, $^{**} p<0.05$,
  $^{***} p<0.01$.}
\end{tablenotes}
\end{threeparttable}

\bigskip
\bigskip
\end{table}








\subsubsection{Deferred-Tax Windfall\label{sec:Windfall}}

TCJA affected more than just the tax on foreign income.
After the reform, deferred tax balances were revalued, and that
revaluation raises a question about how to interpret results in the previous sections. Companies
carrying net deferred tax liabilities booked a one-time gain in equity,
which in theory relaxes a binding financial constraint and encourages
investment. If the companies more exposed to foreign tax advantages were also
the ones receiving the deferred-tax 
gain, then the absence of a differential real response could be two
effects cancelling rather than the absence of either. Our profit-sharing
results could reflect balance sheets rather than booking
incentives. 

To investigate this, we first construct a measure of the deferred-tax windfall, $W_i$:
\begin{equation*}
  W_i = \frac{21-35}{35}\left(\frac{\text{txndb}_{2016}}{\text{at}_{2015}}\right),
\end{equation*}
where ``txndb'' is the net deferred tax balance in 2016
and ``at'' is total assets in 2015, both from
\citet{CompustatAnnual}. 
The balance to asset ratio is multiplied by the change in corporate income
taxes so that it is appropriately scaled for TCJA.
$W_i$ is positive for companies with
a net deferred tax liability---these companies had
lower future tax bill liabilities and an effective increase in book equity.
Companies carrying a net deferred tax asset had the reverse, writing that
asset down.  Here, we exploit differences in $W_i$ using the following
empirical specification:
\begin{equation}
  y_{it} =
  \alpha_i+\lambda_{j(i)t}
  +\theta
  \bigl(W_i\times 1[t\in \text{post}]\bigr)
  +\varepsilon_{it}.
\label{eq:DiDwindfall}
\end{equation}

Before running the regression in equation (\ref{eq:DiDwindfall}), we check
that the one-time deferred-tax remeasurement shows up in
reported accounting earnings in 2017. We estimate the year-by-year analogue
of equation (\ref{eq:DiDwindfall}),
\begin{equation}
  y_{it} =
  \alpha_i+\lambda_{j(i)t}
  +\sum_{s\neq 2016}\theta_s
  \bigl(W_i\times 1[t=s]\bigr)
  +\varepsilon_{it},
\label{eq:FSwindfall}
\end{equation}
with $y_{it}$ equal to net income divided by lagged assets. The estimated
2017 coefficient is \CoefNum{\ThetaShock}, with a 95 percent confidence band
of [\CoefNum{\ThetaShockLo}, \CoefNum{\ThetaShockHi}], and the coefficients
return toward zero afterward. The pooled 2018--2022 coefficient is
\CoefNum{\ThetaPost}, with a standard error of \CoefNum{\ThetaPostSe}.
Thus the deferred-tax measure behaves as expected and produces a transitory
earnings shock in the TCJA year, not a persistent post-TCJA difference in
earnings.

Next, consider the impact of the windfall on investment.
If we use the ratios of capital expenditures or R\&D expenditures to lagged
assets as the outcome variable $y_{it}$,
we find that the deferred-tax windfall has no detectable effect on either.
The estimated coefficient $\hat\theta$ in the case of capital expenditures is equal to
\CoefNum{\WindfallThetaCapex}
with a standard error of 
\CoefNum{\WindfallThetaCapexSe}.
The estimated coefficient in the case of R\&D is equal to
\CoefNum{\WindfallThetaRd}
with a standard error of 
\CoefNum{\WindfallThetaRdSe}.
When we divide the population of firms by size, cash holdings,
leverage, and intangibles with three bins for each, none of the twelve subgroups produces a
positive investment response.


These results also make it unlikely that the deferred-tax windfall is masking
the effects we attribute to international exposure.
We check this possibility in two simple ways. First, we compute the correlation
of $W_i$ and our measure of pre-TCJA international exposure and find it
is \CoefNum{\WcCorr}. 
If anything, the more exposed companies received \emph{smaller}
windfalls, with an average $W_i$ of \CoefNum{\WcMeanExp} compared with
\CoefNum{\WcMeanUnexp} for companies with no exposure.
Second, we add $W_i$ interacted with the post-TCJA indicator to our pooled
difference-in-differences regression (\ref{eq:PooledDiD}) for the foreign profit share.
The international-exposure coefficient is essentially unchanged:
\CoefNum{\WcShrBase}\ \SENum{\WcShrBaseSe}\ without the windfall control and
\CoefNum{\WcShrCtrl}\ \SENum{\WcShrCtrlSe}\ with it, using the
\CountNum{\WcShrFirms}\ companies for which $W_i$ is available.



\subsection{Case Studies\label{sec:Cases}}

An advantage of working with company annual reports 
is that the numbers are typically accompanied by text.
If there is a large closure of the international contribution
to the tax gap, there are likely to be reasons provided
in the notes. Here, we provide examples of companies
with such notes for the TCJA period.

Alphabet reported
that it ``now license[s] intellectual property from the U.S.~that was previously
licensed from Bermuda.'' We see that its international contribution to
the tax reconciliation moved from
\ShareNum{\CaseAlphaPre}\ to \ShareNum{\CaseAlphaPost}\ percentage points and
its foreign profit share from \ShareNum{\CaseAlphaShrPre}\ to
\ShareNum{\CaseAlphaShrPost}\ percent. 
Microsoft reported an ``impact of
intangible property transfers'' worth $-5.9$ percentage points in fiscal 2019.
We see a partial move given they kept some operations in Ireland.
ON Semiconductor reported an ``impact of the domestication'' worth $-35.7$
percentage points after moving its Bermuda subsidiaries to Delaware. The
international contribution to the tax reconciliation closed by \ShareNum{\CaseOnSemiClose}\ points.
Qualcomm reported that one of its Singapore subsidiaries ``distributed certain
intellectual property to a U.S.~subsidiary'' and later claimed a \$419 million
FDII deduction. We see its foreign profit share fell from
\ShareNum{\CaseQcomShrPre}\ to \ShareNum{\CaseQcomShrPost}\ percent.
In these cases, the filings give a concrete reason for the change in the tax
reconciliation, and the reason fits the post-TCJA incentives we study.


Coca-Cola is the useful counterexample. Their international contribution
closes by \ShareNum{\CaseKoffClose}\ percentage points---among the
largest in the sample---but the 10-K filing attributes the change to a notice
received from the IRS in September 2015 related to tax years 2007 through
2009.  The IRS can require a
company to price cross-border transactions at arm's length and to
report a larger share of its profits in the United States under Section 482
of the Internal Revenue Code. 
Coca-Cola went through a five-year audit and a subsequent reallocation of
income.  The resulting
changes reported in the reconciliation tables look like the movements above, but the
cause was enforcement of existing law rather than the change in the law we
are trying to measure.

Transfer-pricing enforcement of this kind is therefore an alternative 
explanation for the decline in foreign tax advantages
and could be driving our main estimates. To investigate this
hypothsis, we collect published transfer-pricing opinions involving U.S.~multinationals 
and match the companies to \citet{CompustatAnnual}. This
identifies \CountNum{\LitNLit}\ companies in our estimation sample of
\CountNum{\LitNFirms}. Dropping these companies has little effect on the
estimates. 
The international coefficient $\hat\beta_{\text{Intl}}$  changes from
\CoefNum{\LitIntlFull}\ \SENum{\LitIntlFullSe}\ to
\CoefNum{\LitIntlDrop}\ \SENum{\LitIntlDropSe}, 
the effective tax rate coefficient
$\hat\beta_{\text{ETR}}$ changes from
\CoefNum{\LitEtrFull}\ \SENum{\LitEtrFullSe}\ to
\CoefNum{\LitEtrDrop}\ \SENum{\LitEtrDropSe}; and the foreign-profit-share coefficient
$\hat\beta_{\text{Share}}$ changes from
\CoefNum{\LitShrFull}\ \SENum{\LitShrFullSe}\ to
\CoefNum{\LitShrDrop}\ \SENum{\LitShrDropSe}. In other words,
two of the three estimates are farther from zero after dropping firms in litigation.
This is only a conservative check, because published opinions do not include
every company with a transfer-pricing dispute, but that limitation works against
finding no effect. If dropping the litigants we can identify does not change the
results, unidentified disputes are unlikely to explain them.


\section{A Neoclassical Interpretation\label{sec:Theory}}

In this section, we show that our main empirical result---that
reported multinational profits relocated  in response to TCJA,
while production did not---is consistent with a neoclassical theory
of multinational activity.
A minimal structure assumes that multinationals can pay to 
conceal some profits from the U.S.~government by shifting
income in amount $S$ to a low-tax affiliate through a transfer pricing scheme
(as in \citet{HinesRice1994}). We use the theory to estimate the
share of the international contribution $x_{\text{Intl}}$ 
attributable to shifting
profits to tax havens with no substantive productive operations.
Relatedly, we estimate the share of foreign profits that is
shifted income $S$, both before and after TCJA.

\subsection{Multinational Optimization}

Consider the problem of a multinational with domestic ($d$)
and foreign ($f$) subsidiaries that maximizes world wide
cash flow discounted at rate $r$ for its shareholders.
This cash flow is net of income taxes, capital adjustment
costs, and profit shifting concealment costs.

Let $X_{\scriptstyle{R\&D},t}$ be the 
investment in R\&D that is created in the United States and
used nonrivalrously at home and abroad, and let
$X_{\ell,t}$ be the investment in tangible capital in location
$\ell\in\{d,f\}$. For tangible capital, we assume that there
are costs of adjustment given by
$\Phi_\ell(X_{\ell,t},K_{\ell,t})$, where the stock accumulates as
$K_{\ell,t+1}=(1-\delta)K_{\ell,t}+X_{\ell,t}$.
Firms also pay a cost of concealing shifted profits $S$ equal
to $C_i(S)$.
Here, we assume that a shifting of profits is purely a rebooking
that lowers taxable profits in the United States without
any requirement of productive factors. In other words,
profits that should be included as taxable in the United States
are instead reported on the books of a subsidiary in a tax haven ($h$)
that does not produce anything. 

The multinational's problem then is
to choose $\{S_{it},X_{{\scriptstyle R\&D},it},X_{d,it},X_{f,it}\}$
to maximize the net present value of cash flows
to shareholders, $\sum_t(1+r)^{-t}\mathcal{N}_{it}$, after paying
taxes at rate $\tau_{\ell,t}$ on profits in $\ell\in\{d,f,h\}$ with
\begin{align}
  \mathcal{N}_{it} & = \sum_{\ell\in\{d,f\}} \left\{(1-\tau_{\ell,t}) \pi_{\ell,it}-(1-\tau_{\ell,t} z_{\ell})X_{\ell,it} 
                         - \Phi_\ell(X_{\ell,it},K_{\ell,it})\right\}\nonumber \\
                &\qquad\qquad\qquad -  (1-\tau_{d,t} z_{{\scriptstyle R\&D}}) X_{{\scriptstyle R\&D},it} 
                          +(\tau_{d,t}-\tau_{h,t})S_{it} - C_i(S_{it}),
\label{eq:cashflow}
\end{align}
where depreciation allowances by tax authorities are denoted by
$z_jX_{j,t}$, $j\in\{{\scriptstyle R\&D},d,f\}$.\footnote{%
Here, we simplify our analysis by using the same $\tau_f$ and the
same $\tau_h$ for every firm, despite the fact that firms in our
sample might face different effective rates abroad.}
The last two terms in (\ref{eq:cashflow}) are the tax savings paid
to shareholders for profit shifting and the costs of concealing this
activity.

Because shifting profits requires no 
additional productive resources, the firm's problem can be
separated into the real productive decision-making of the firm---that is,
how much to invest and where to locate the capital---and the profit 
shifting. This is perhaps clearer if we write $\mathcal{N}_{it}$ 
as the sum of two payments: the ``real'' block that includes only cash flow
from profits and the ``booking'' block that includes
only cash flow from profit shifting---which, in this case, is the last
two terms of equation (\ref{eq:cashflow}).
This is an important distinction when considering the impacts of TCJA.
Analyses of investment decisions following TCJA
compare outcomes of firms with different user costs of capital:
in a steady state without adjustment costs, the real block implies the
condition of \citet{HallJorgenson1967},
\begin{equation}
  \frac{\partial \pi_{\ell,it}}{\partial K_{\ell,it}}
   = (r+\delta)\,\frac{1-\tau_{\ell,t} z_\ell}{1-\tau_{\ell,t}},
   \qquad \ell\in\{d,f\},
\label{eq:usercost}
\end{equation}
and an analogous condition for R\&D, which equates
the marginal product of capital to the user cost of capital.
(See, for example, \citet{CSZZ2024} and \citet{GravelleMarples2025},
who build on this framework.)
In contrast, our analysis compares outcomes of firms with different exposure
to foreign tax advantages and thus is concerned with the booking
decision, which in our environment is the choice
$S^*_{it}$ that satisfies $C_i'(S^*_{it})=\tau_{d,t}-\tau_{h,t}$.
Neither the domestic nor foreign user cost in  (\ref{eq:usercost})
depends on the wedge $\tau_{d,t}-\tau_{h,t}$,
and $S^*_{it}$ depends on nothing else: this is the separation.
The distinction is the average-versus-marginal one of
\citet{DevereuxGriffith2003}: the booking choice responds to the
\emph{average} tax wedge on booked income, while the investment
choices respond to \emph{marginal} prices.

\subsection{Mapping Theory and Empirical Design}

We show next how the booking choice $S^*$ is related to the main variable of 
interest in our empirical design, namely the international
contribution in the reconciliation table. Assuming 
that this contribution were the main driver of the
gap between effective and statutory rates---which we will do here---we have
an exact relation:
\begin{align}
  x_{\text{Intl},it} & =  \text{ETR}_{it}-\tau_{d,t}
  \nonumber\\
  \noalign{\bigskip}
                     & = \frac{\tau_{d,t}(\pi_{d,it}-S_{it})+\tau_{h,t}S_{it}+\tau_{f,t}\pi_{f,it}}{\pi_{it}}-\tau_{d,t}
  \nonumber \\
  \noalign{\bigskip}
                     & = \frac{-(\tau_{d,t}-\tau_{h,t})S_{it}-(\tau_{d,t}-\tau_{f,t})\pi_{f,it}}{\pi_{it}}
\label{eq:xintit}
\end{align}
for firm $i$ at time $t$.
As we noted above, we want to estimate the extent of income shifting that is purely a rebooking
of $S$. To do this, we need to estimate the contribution of the first term in (\ref{eq:xintit})---which
is the tax wedge times the share of profits shifted---to the total
international contribution $x_{\text{Intl}}$ (both in absolute value),
evaluated at pre-TCJA quantities:
$\psi_i\equiv(\tau_d^{\text{pre}}-\tau_h^{\text{pre}})
(S_i^{\text{pre}}/\pi_i)/|x^{\text{pre}}_{\text{Intl},i}|$. We refer to
$\psi_i$ as the \emph{profit-shifting share} of the
international contribution to the tax reconciliation.

We show next that the profit-shifting share $\psi_i$ can be inferred from 
our empirical estimates if we have information on statutory tax rates and any
changes that occurred under TCJA.
Through the lens of the model, the coefficient $\beta$ in
(\ref{eq:PooledDiD}) is the change in an outcome per percentage point
of pre-TCJA exposure $x^{\text{pre}}_{\text{Intl},i}$. Let
$\beta_{\text{Real}}$ denote this coefficient when the outcome is a
measure of real activity---capital expenditures or sales growth,
say---and $\beta_{\text{Intl}}$ and $\beta_{\text{Share}}$ its model
counterparts for the international contribution and the foreign profit
share (in the first and third rows of Table \ref{tab:PooledDiD}). 
For clarity, the formulas below treat $\psi$ as common across firms.\footnote{%
If $\psi$ differs across firms, the regression coefficients should be read as
describing the firms that drive the identifying variation, not the average firm
in the population.}
The model implies $\beta_{\text{Real}}=0$ and
\begin{align}
  \beta_{\text{Intl}} & = -(1-\psi) \left[1-\frac{\tau_d^{\text{post}}-\tau_f^{\text{post}}}
                                                     {\tau_d^{\text{pre}}-\tau_f^{\text{pre}}}\right]
                              -\psi \left[1-\left(\frac{\tau_d^{\text{post}}-\tau_h^{\text{post}}}
                                                        {\tau_d^{\text{pre}}-\tau_h^{\text{pre}}}\right)^2\right]
 \label{eq:betaintl}\\
 \noalign{\bigskip}
  \beta_{\text{Share}} & = \left[\frac{\psi}{\tau_d^{\text{pre}}-\tau_h^{\text{pre}}}\right]
                                 \left[1-\frac{\tau_d^{\text{post}}-\tau_h^{\text{post}}}
                                            {\tau_d^{\text{pre}}-\tau_h^{\text{pre}}}\right].
 \label{eq:betashr}
\end{align}
We unpack each in turn.

The prediction $\beta_{\text{Real}}=0$ has two parts. First, by the separation,
the tax wedge $\tau_{d,t}-\tau_{h,t}$ relevant for the booking decision
enters no real first-order
condition. Second, TCJA's changes to the user costs in
(\ref{eq:usercost})---for example, the drop in $\tau_d$ from 35 to 21
percent---depend only on statutory parameters common to all firms in
an industry-year cell, so the industry-year fixed effects absorb them.
We maintain the assumption that firms' responses to these common
changes are orthogonal to their exposure. The prediction is borne out
in the data. Sales, capital expenditures, R\&D, and employment show no
significant differential response to exposure, and neither do the
foreign shares of sales, the installed capital stock, or total pretax income.
The balance sheets of more exposed firms do
contract, but that contraction is concentrated in cash rather than in
productive capital, reflecting repatriation and the transition tax.
Based on these results, we hold profits 
from foreign operations, $\pi_f$, and total profits, $\pi$, fixed when unpacking
the two coefficients in (\ref{eq:betaintl}) and (\ref{eq:betashr}).

The change in the international contribution to the tax reconciliation
combines two effects: a smaller
tax-rate advantage on foreign operations, and less income shifted to low-tax
affiliates. To derive the expression for the regression coefficient
$\beta_{\text{Intl}}$ in
(\ref{eq:betaintl}), note that with the quadratic concealment cost of
\citet{HinesRice1994}, the optimal shifting rule sets $S$ proportional
to the tax gap, so
$S_i^{\text{post}}/S_i^{\text{pre}}
=(\tau_d^{\text{post}}-\tau_h^{\text{post}})/
(\tau_d^{\text{pre}}-\tau_h^{\text{pre}})$. Using this 
expression and the definition of $\psi_i$, we can
derive the expression for $\beta_{\text{Intl}}$: 
\begin{align*}
  \Delta x_{\text{Intl},i} 
             & = -\frac{(\tau_d^{\text{post}}-\tau_h^{\text{post}})S_i^{\text{post}}
                      -(\tau_d^{\text{pre}}-\tau_h^{\text{pre}})S_i^{\text{pre}}}{\pi_i}\\
             &\qquad\qquad
                 -\frac{(\tau_d^{\text{post}}-\tau_f^{\text{post}}-
                    \tau_d^{\text{pre}}+\tau_f^{\text{pre}})\pi_{f,i}}{\pi_i}\\
  \noalign{\bigskip}
            & = \underbrace{-(1-\psi)\left[1-\frac{\tau_d^{\text{post}}-\tau_f^{\text{post}}}
                                                 {\tau_d^{\text{pre}}-\tau_f^{\text{pre}}}\right]
                              -\psi\left[1-\left(\frac{\tau_d^{\text{post}}-\tau_h^{\text{post}}}
                                                    {\tau_d^{\text{pre}}-\tau_h^{\text{pre}}}\right)^{2}\right]}
                _{\beta_{\text{Intl}}}\; x^{\text{pre}}_{\text{Intl},i} 
\end{align*}
with simple algebra.
This expression has a simple interpretation. The pre-TCJA international contribution
contains two kinds of foreign tax benefits. One comes from foreign
profits $\pi_f$ taxed below the U.S.~rate. The other comes from shifted income, $S$, booked in
the tax haven.  TCJA reduces both benefits, but for different reasons. For foreign
operations, the underlying profits $\pi_f$ are held fixed, so the tax benefit
shrinks only because the U.S.-foreign tax gap narrows. For shifted income, both
the tax saving per shifted dollar and the amount shifted fall with the
U.S.-haven tax gap. This is why the operating part enters with the simple gap
closure, while the shifted part enters with the square of the gap ratio.

The expression for $\beta_{\text{Share}}$ in (\ref{eq:betashr})
can be written in terms of $\psi$ and the U.S.-haven tax gap:
\begin{align*}
   \Delta \left(\frac{\pi_{f,i}+S_i}{\pi_i}\right) 
            &= 
               \frac{\pi_{f,i}+S_i^{post}}{\pi_i}
              -\frac{\pi_{f,i}+S_i^{pre}}{\pi_i}\\
  \noalign{\bigskip}
            & = \underbrace{\left[\frac{\psi}{\tau_d^{\text{pre}}-\tau_h^{\text{pre}}}\right]
                 \left[1-\frac{\tau_d^{\text{post}}-\tau_h^{\text{post}}}
                              {\tau_d^{\text{pre}}-\tau_h^{\text{pre}}}\right]}
                _{\beta_{\text{Share}}}\; x^{\text{pre}}_{\text{Intl},i}.
\end{align*}
Unlike $x_{\text{Intl},i}$, the foreign profit share is an
income share rather than a tax-rate category. Since real foreign operating profits,
$\pi_f$, are held fixed, the only part of the numerator that changes is shifted
income, $S$. The coefficient $\beta_{\text{Share}}$ therefore captures how much
shifted income falls relative to total profits.

We now use our estimated coefficients, together with statutory tax rates, to
infer two model-implied quantities: the profit-shifting share of the international
contribution to the tax reconciliation, $\psi$, and the tax
rate on real foreign operations, $\tau_f$, which we treat as unchanged by TCJA.\footnote{%
Following \citet{Hartman1985}, we treat the pre-TCJA residual U.S.~tax on
foreign income as deferred and not paid at the margin, so the relevant pre-TCJA
rates are the local rates $\tau_f$ and $\tau_h^{\text{pre}}$.}
These quantities then allow us to compute the share of foreign booked income accounted
for by shifted income, $S/(S+\pi_f)$.
The statutory rates for the United States pre- and post-TCJA
are $\tau_d^{\text{pre}}=35$ percent and $\tau_d^{\text{post}}=21$ percent, respectively.
Most estimates for pre-TCJA tax haven rates
are in the range of 0 to 10 percent, which we will investigate by considering
three cases, $\tau_h^{\text{pre}}\in \{0,5,10\}$ percent. 
For the post-TCJA tax haven rate, we consider two cases. In the first, firms have
enough foreign tax credits from other foreign income that low-tax haven income
faces no residual GILTI tax and $\tau_h^{\text{post}}=\tau_h^{\text{pre}}$. 
In the second, firms pay the statutory GILTI residual on tax haven income, 
which we model as:
\begin{equation}
   \tau_h^{\text{post}} = \tau_h^{\text{pre}}+\max\{0,\,0.105-0.8\tau_h^{\text{pre}}\}.
\label{eq:tauhpost}
\end{equation}
This provision adds a residual tax equal to the 10.5 percent GILTI rate less a
credit for 80 percent of the foreign tax already paid, or
$0.105-0.8\tau_h^{\text{pre}}$.\footnote{%
See \citet{GravelleGILTI2024} and \citet{Herzfeld2021} on the GILTI rate,
foreign tax credits, and the role of excess credits. We apply this adjustment to
tax haven income, but not to real foreign operations.}

In Table \ref{tab:InvertedPsi}, we show the main results, varying $\tau_h^{\text{pre}}$ (in columns)
and the GILTI regime (in panels). In order to use a consistent sample for estimating
coefficients $\beta_{\text{Intl}}$ and $\beta_{\text{Share}}$ we rerun
the pooled difference-in-differences 
(shown in the first row of Table \ref{tab:PooledDiD})
with the more restricted sample used
for the foreign profit share (shown in the third row). This yields
a value for $\hat\beta_{\text{Intl}}$
of \CoefNum{\BetaIntlSame}\ (s.e.\ \CoefNum{\BetaIntlSameSe}) on the
\RoundNum{\BetaIntlSameG}\ firms that report the split, against
\CoefNum{\IntlCoef}\ on the full panel.
We choose three values for $\tau_h^{\text{pre}}$ of 0, 5, and 
10 percent. In the no-GILTI case, we use the same value for $\tau_h^{\text{post}}$
and therefore do not report it. Instead we report the implied values
for the foreign tax on productive operations $\tau_f$, the profit-shifting
share $\psi$ of the international contribution
and the ratios of shifted profit $S$ to total foreign profits
$S+\pi_f$. As the table shows, the implied
foreign tax rate is well above the tax haven rate---in the range of
\RoundNum{\GTauFNoGiltiTen} to \RoundNum{\GTauFNoGiltiZero}
percent---but in line with estimates of \emph{effective} foreign
tax rates paid by multinationals outside of tax havens.\footnote{%
We use a firm-level bootstrap to re-estimate $\hat\beta_{\text{Intl}},\hat\beta_{\text{Share}}$ 
and construct 95 percent confidence intervals of 
[\RoundNum{\TauFCiLoNoGiltiPct}, \RoundNum{\TauFCiHiNoGiltiPct}] percent for $\tau_f$ 
and 
[\RoundNum{\PsiCiLoNoGiltiPct}, \RoundNum{\PsiCiHiNoGiltiPct}] percent for $\psi$ 
when $\tau^h_{\text{pre}}=5$ percent.}



\begin{table}[!tp]
\centering
\begin{threeparttable}
\caption{\textsc{Profit-Shifting Implications, Varying Pre-TCJA Tax Haven Rates }}
\label{tab:InvertedPsi}

   \begin{tabular}{lcrrr}
   \hline
   \hline
   \noalign{\vskip 5pt}
   & & \multicolumn{3}{c}{$\%\tau_h^{\text{pre}}$}\tabularnewline
    \noalign{\vskip 3pt}
   \cline{3-5}
   \noalign{\vskip 3pt}
   \textsc{Outcome} &\phantom{hello} & 0\ \  & 5\ \  & 10\ \  \tabularnewline
   \noalign{\vskip 3pt}
   \hline
   \noalign{\vskip 6pt}
   A. No GILTI & &   & &  \tabularnewline  
    \noalign{\vskip 5pt}
   \ \ \ \ \textsc{Tax rate on productive foreign profits} ($\tau_f$)  & & \ShareNum{\GTauFNoGiltiZero} & \ShareNum{\GTauFNoGiltiFive} & \ShareNum{\GTauFNoGiltiTen} \tabularnewline
    \noalign{\vskip 5pt}
   \ \ \ \ \textsc{Profit-shifting share of} $x^{\text{pre}}_{\text{Intl}}$ ($\psi$) & & \ShareNum{\GPsiNoGiltiZero} & \ShareNum{\GPsiNoGiltiFive} & \ShareNum{\GPsiNoGiltiTen} \tabularnewline
    \noalign{\vskip 5pt}
   \ \ \ \ \textsc{Pre-TCJA shifted share of profits} ($S^{\text{pre}}/(S^{\text{pre}}+\pi_f)$)  & & 
                      \ShareNum{\GShiftedPreNoGiltiZero}  & \ShareNum{\GShiftedPreNoGiltiFive} & \ShareNum{\GShiftedPreNoGiltiTen} \tabularnewline
    \noalign{\vskip 5pt}
   \ \ \ \ \textsc{Post-TCJA shifted share of profits} ($S^{\text{post}}/(S^{\text{post}}+\pi_f)$)  & & 
                      \ShareNum{\GShiftedPostNoGiltiZero} & \ShareNum{\GShiftedPostNoGiltiFive} & \ShareNum{\GShiftedPostNoGiltiTen} \tabularnewline
    \noalign{\vskip 10pt}
   B. Statutory GILTI & &   & &  \tabularnewline  
    \noalign{\vskip 5pt}
   \ \ \ \ \textsc{Post-TCJA tax haven tax rate} ($\tau_h^{\text{post}}$) & & \ShareNum{\GTauHPostStatZero} & \ShareNum{\GTauHPostStatFive} & \ShareNum{\GTauHPostStatTen} \tabularnewline
    \noalign{\vskip 5pt}
   \ \ \ \ \textsc{Tax rate on productive foreign profits} ($\tau_f$)  & & \ShareNum{\GTauFStatZero} & \ShareNum{\GTauFStatFive} & \ShareNum{\GTauFStatTen}  \tabularnewline
    \noalign{\vskip 5pt}
   \ \ \ \ \textsc{Profit-shifting share of} $x^{\text{pre}}_{\text{Intl}}$ ($\psi$) & & \ShareNum{\GPsiStatZero} & \ShareNum{\GPsiStatFive} & \ShareNum{\GPsiStatTen} \tabularnewline
    \noalign{\vskip 5pt}
   \ \ \ \ \textsc{Pre-TCJA shifted share of profits} ($S^{\text{pre}}/(S^{\text{pre}}+\pi_f)$)  & & 
                      \ShareNum{\GShiftedPreStatZero}  & \ShareNum{\GShiftedPreStatFive} & \ShareNum{\GShiftedPreStatTen} \tabularnewline
    \noalign{\vskip 5pt}
   \ \ \ \ \textsc{Post-TCJA shifted share of profits} ($S^{\text{post}}/(S^{\text{post}}+\pi_f)$)  & & 
                      \ShareNum{\GShiftedPostStatZero} & \ShareNum{\GShiftedPostStatFive} & \ShareNum{\GShiftedPostStatTen} \tabularnewline
   \noalign{\vskip 6pt}
   \hline
   \end{tabular}

\bigskip
\begin{tablenotes}[flushleft]
\item[] \footnotesize{\textit{Notes:} All entries are percentages.
The table reports estimates of key ratios starting from the
same-sample regression estimates
$\hat\beta_{\mathrm{Intl}}=\CoefNum{\BetaIntlSame}$ and
$\hat\beta_{\mathrm{Share}}=\CoefNum{\PiShrCoef}$ and a value for
$\tau_h^{\text{pre}}$. Panel A assumes no GILTI liability on shifted
profits, so that $\tau_h^{\text{post}}=\tau_h^{\text{pre}}$; Panel B
assumes the statutory GILTI residual, so that $\tau_h^{\text{post}}$
is given by (\ref{eq:tauhpost}). The inversion is closed form.
Writing $R_h$ and $R_o$ for the post-to-pre ratios of the tax gaps on
shifted and productive foreign profits, equation (\ref{eq:betashr})
gives $\psi=\hat\beta_{\mathrm{Share}}
(\tau_d^{\text{pre}}-\tau_h^{\text{pre}})/(1-R_h)$, equation
(\ref{eq:betaintl}) then gives
$R_o=1-[|\hat\beta_{\mathrm{Intl}}|-(1-R_h^2)\psi]/(1-\psi)$, and
$\tau_f$ solves $R_o=(0.21-\tau_f)/(0.35-\tau_f)$. The shifted
shares use
$S^{\text{pre}}/\pi=\psi|x^{\text{pre}}_{\text{Intl}}|/
(\tau_d^{\text{pre}}-\tau_h^{\text{pre}})$ and
$\pi_f/\pi=(1-\psi)|x^{\text{pre}}_{\text{Intl}}|/
(\tau_d^{\text{pre}}-\tau_f)$, so $S/(S+\pi_f)$ is free of the
exposure; post-TCJA, $S^{\text{post}}=R_hS^{\text{pre}}$.}
\end{tablenotes}
\end{threeparttable}

\bigskip
\bigskip
\end{table}



The key variable of interest $\psi$---which is listed next in Table \ref{tab:InvertedPsi}---provides
a measure of foreign tax advantages prior to the tax reform.
Across the range of estimates for $\tau_h^{\text{pre}}$,
we find the share to be economically large: \RoundNum{\GPsiNoGiltiTen} if the tax rate
is at 10 percent and \RoundNum{\GPsiNoGiltiZero} percent if it is zero.
Using the estimate for $\psi$, we can infer the ratio of shifted profits
to all foreign profits and compare this to \citet{TorslovWierZucman2023}
who find an estimate over one-half for the United States in 2015.  Our highest
estimate---assuming GILTI provisions are not binding and tax haven rates
are at 0 percent---is \RoundNum{\GShiftedPreNoGiltiZero} percent in the
period prior to TCJA. If we assume $\tau_h^{\text{pre}}$ is 10 percent,
then the ratio falls to \RoundNum{\GShiftedPreNoGiltiZero}.
The estimates differ from \citet{TorslovWierZucman2023}
because ours are incentive-based, 
using TCJA-induced changes in tax advantages, while theirs are residual-based, 
classifying unexplained tax haven profitability as shifted profits.
With our empirical design, we are also able to infer how
this ratio changed after the tax reform.
We find a significant decline, with the range of estimates given by
\RoundNum{\GShiftedPostNoGiltiTen} to
\RoundNum{\GShiftedPostNoGiltiZero} percent after-TCJA.


In the bottom panel of Table \ref{tab:InvertedPsi}, we show the cases with
GILTI provisions binding.  Applying the formula in
(\ref{eq:tauhpost}), we see that GILTI puts a relatively 
high tax rate floor on the tax havens. To account for the same foreign profit
share and closure of the international contribution as in the no-GILTI case, 
we need smaller implied values for the  foreign tax on profits ($\tau_f$) and 
the pre-TCJA profit-shifted share ($\psi$) of $\bar x^{\text{pre}}$.
It turns out that there is also much less variation in these estimates
for the different $\tau_h^{\text{pre}}$. The foreign tax rate $\tau_f$ is roughly constant
in the range of \RoundNum{\GTauFStatTen} to \RoundNum{\GTauFStatZero}
percent. The share $\psi$ is now in the range of
\RoundNum{\GPsiStatTen} to \RoundNum{\GPsiStatZero}.\footnote{%
When $\tau_h^{\text{pre}}=5$ percent, the 95 percent confidence intervals for $\tau_f$ is
[\RoundNum{\TauFCiLoStatPct}, \RoundNum{\TauFCiHiStatPct}] percent 
and for $\psi$ is
[\RoundNum{\PsiCiLoStatPct}, \RoundNum{\PsiCiHiStatPct}] percent.
We also consider an alternative in which rebooked profits are taxed at the 13.1 percent FDII 
rate rather than the 21 percent corporate rate. 
In this case, $\psi$ falls to \RoundNum{\PsiInvFdiiStatPct}--\RoundNum{\PsiInvFdiiNoGiltiPct} 
percent, while $\tau_f$ remains near \RoundNum{\TauFImpliedFdiiStatPct} percent.}
Not surprisingly, the GILTI provisions are relevant when the
tax haven rate is close to zero.
The final estimates in the table are the 
pre- and post-TCJA estimates for the shifted share of foreign 
profits, $S/(S+\pi_f)$, when GILTI provisions are binding.  As with $\tau_f$ and
$\psi$, these provisions imply less variation in the estimates
given GILTI puts an effective floor on the tax haven rates.
Here we find the pre-TCJA shifted share is 
in the range of \RoundNum{\GShiftedPreStatTen} to \RoundNum{\GShiftedPreStatZero},
which is roughly half of the residual-based estimate of \citet{TorslovWierZucman2023}.
Post-TCJA this share drops to a tight range of
\RoundNum{\GShiftedPostStatTen} to \RoundNum{\GShiftedPostStatZero},

In either GILTI case, the model gives a simple measure of the change in shifting.
Because shifted income $S$ is proportional to the tax gap, the fraction of
shifted income that returns to the United States,
$(S^{\text{pre}}-S^{\text{post}})/S^{\text{pre}}$, equals the decline in the
shifting incentive,
$1-(\tau_d^{\text{post}}-\tau_h^{\text{post}})/
(\tau_d^{\text{pre}}-\tau_h^{\text{pre}})$. This fraction ranges from 40 to
70 percent across the cells of Table \ref{tab:InvertedPsi}.



\section{Conclusion}

The TCJA substantially reduced the foreign tax advantages reported by U.S.~multinationals. 
Using firms' income tax reconciliations, we show that companies
with larger pre-TCJA international tax benefits experienced larger post-TCJA
increases in effective tax rates and larger declines in foreign profit shares.
At the same time, we find little evidence that these firms differentially changed
their sales, employment, investment, R\&D, or the geographic location of their
productive capital. The pattern points to a change in profit booking rather than
a broad relocation of real activity: the reform weakened the tax incentive to
book income abroad, while changes in the marginal tax on investment were largely
common to firms within the same industry.


\newpage


\bibliographystyle{jpe}
\bibliography{tcja_refs}


\newpage

\appendix

\counterwithin{table}{section}
\counterwithin{figure}{section}

\section{Data Construction}

In this appendix, we provide additional details on our
data construction.
The source population is Compustat-linked annual 10-K
filings. For each filing, we identify the income-tax rate reconciliation table,
extract the reconciliation rows, and convert the result to a common
firm-year-by-category panel.
We also include specific SIC codes used for industry classification.

Every U.S.-publicly-listed company must report a reconciliation table 
that decomposes accounting items detailing the gap between the statutory 
tax rate and their effective rate. We collect the SEC filing URLs for 
10K filings that can be merged with \citet{CompustatAnnual}. In Table
\ref{tab:StageOneCounts}, we report that the total number of URLs used in our
stage 1 extraction is  \CountNum{\InputUrlsTotal}.  Once we eliminate duplicates,
the total number of unique SEC filings is \CountNum{\ReconUniqueFilings}.
The parsing methods we use are designed for HTML formats. If we exclude
non-HTML formats and the negligible number of downloads that failed, we have
\CountNum{\PhaseOneStatusOK} usable filings to parse  and classify.

Table \ref{tab:StageOneClassifications} shows the results of the 
classification of the \CountNum{\PhaseOneStatusOK} usable filings.
Each one is parsed into text blocks and candidate tables. We use
key words, symbols and phrases from the reconciliation tables to
score each table. The score represents the likelihood that 
a reconciliation table will be extractable. We summarize categories
in the table.  Those scored high-confidence, likely but incomplete, or
ambiguous account for roughly seven-tenths of the unique SEC filings.
The remaining categories are filings with low scores or for which no
reconciliation table was found. Tables may be missing if the
information is not required under ASC 740. Every packet stage 1 produces is passed to stage 2
regardless of its bucket, and the lower-scoring categories still yield a
small number of usable reconciliations. The count of packets is slightly
below the count of filings that stage 1 completed:
\CountNum{\ArNoPacket}\ filings finished without a candidate table to
pass on, and a further \CountNum{\ArRescueOk}\ come from a later
rescue pass whose packets are parsed separately.

The output of stage 1 is an audited packet containing the best candidate
table or tables and nearby text. The goal of the second stage is to
turn each of these packets into structured firm-year observations on a common
schema. Table \ref{tab:StageOneToTwo} reports this transition: of the
\CountNum{\PacketsProcessed} packets parsed in stage 2,
\CountNum{\ProcessedWithRows} yield at least one reconciliation row and
\CountNum{\ProcessedNoRows} yield none.
In Table \ref{tab:StageTwoParsing}, we report the total count
in the first row and then split the total three ways. First, we report
the number of observations that pass or fail an arithmetic check, namely,  
that the sum of reconciliation values equals the gap between the effective 
and statutory tax rates, within a tolerance of 0.2 percentage points for
tables reported to a decimal and 1.0 point for those reported in whole
numbers or in amounts. Most pass.  Second, we report the number
of observations that do or do not require LLM assistance when parsing.
In other words, we deterministically parse the packet to first detect 
column headers that are consecutive
years, then locate rows starting with the statutory rate or amount, subsequent 
adjustments, and the effective rate or amount, and finally to do the
arithmetic check on row sums. If the parsing or arithmetic checks fail,
the stage 1 packet is reassessed by GPT-5.2. Third, we report
the number of reconciliations that use tax rates and the number
that use tax amounts.

Output from the second stage parsing is filtered to create the 
final firm panel dataset used in the regressions.  We keep 
observations in the period 1998 to 2024 that have a match
in \citet{CompustatAnnual} and pass the arithmetic check.
Since each 10K filing reports the tax reconciliation for the 
current year as well as, typically, two past years, we deduplicate the dataset
by the priority rules described in Section \ref{sec:Data}: own-year
before comparative, passing the arithmetic check before failing it,
original filings before amendments, higher-confidence extractions
before lower, and the earliest filing before later restatements. Table \ref{tab:StageTwoToFinal} shows row counts as
we filter the stage 2 data.
After the filtering is complete, we have a total of
\CountNum{\AnalysisPanelFirmYears} firm-year observations.
For some results, we also apply industry filters and use
SIC codes summarized in Table \ref{tab:SICcodes}.

\subsection{Coverage by Year\label{app:coverage}}

Figure \ref{fig:CoverageByYear} reports the share of the Compustat
population of Table \ref{tab:Coverage} for which we recover a
reconciliation, by year. Coverage rises with the spread of machine-readable
filings: it is \ShareNum{\ReconCovInit}\ percent in 1998, passes
two-thirds in 2004---which is why the descriptive figures of Section
\ref{sec:Results} begin in that year---and reaches
\ShareNum{\ReconCovGR}\ percent by 2008. Sales-weighted coverage is
higher throughout, because the firms we miss are small.

\begin{figure}[!tp]
\centering
\caption{\textsc{Coverage of the Analysis Panel by Year}}
\label{fig:CoverageByYear}
\begin{tikzpicture}
\begin{axis}[
    width=0.82\textwidth, height=0.42\textwidth,
    xmin=1997.5, xmax=2024.5, ymin=0, ymax=100,
    xlabel={}, ylabel={Coverage (percent)},
    xtick={1998,2002,2006,2010,2014,2018,2022},
    ytick={0,20,40,60,80,100},
    scaled x ticks=false,
    x tick label style={/pgf/number format/1000 sep=},
    tick label style={font=\footnotesize},
    label style={font=\footnotesize},
    legend style={font=\footnotesize, at={(0.98,0.05)}, anchor=south east, draw=none, fill=white},
    grid=major, grid style={gray!20},
    axis lines=left,
]
\addplot[thick, black, mark=*, mark size=1.1] table[col sep=comma, x=year, y=sales_pct]
    {\Inputs/coverage_time_tikz/coverage_by_year.csv};
\addlegendentry{Sales-weighted}
\addplot[thick, black, dashed, mark=o, mark size=1.1] table[col sep=comma, x=year, y=equal_pct]
    {\Inputs/coverage_time_tikz/coverage_by_year.csv};
\addlegendentry{Equal-weighted}
\draw[gray, dashed] (axis cs:2004,0) -- (axis cs:2004,100);
\end{axis}
\end{tikzpicture}

\bigskip
\begin{minipage}{0.95\textwidth}
\footnotesize{\textit{Notes:} Share of the Compustat population defined in
Table \ref{tab:Coverage} for which the analysis panel contains a
reconciliation. The equal-weighted series is the share of firm-years; the
sales-weighted series weights by sales. The vertical line marks 2004, the
first year in which equal-weighted coverage exceeds two-thirds.}
\end{minipage}

\bigskip
\bigskip
\end{figure}

\subsection{Assigning Reconciliation Rows to Categories\label{app:categories}}

Stage 2 returns each reconciliation row with the label its filer wrote,
and those labels are free text: \CountNum{\CatDistinctLabels} distinct
strings appear in the panel. ``Effect of foreign operations,''
``Jurisdictional rate differences,'' and ``Earnings of subsidiaries
taxed at other than U.S.\ statutory rate'' are all the foreign rate
differential, and none shares a common substring with the others. A
third stage therefore maps labels to the categories of Table
\ref{tab:Taxonomy}.

The mapping is a deterministic cascade of \CountNum{\CatNumRules}
ordered regular expressions applied to the lower-cased label. The first
rule that matches assigns the row to one of
\CountNum{\CatNumCategories} fine categories, which are then folded
into the seven categories of Table \ref{tab:Taxonomy}; a row matching
no rule is assigned to \emph{other}. The assignment is exhaustive and
one-to-one---every row receives exactly one category, and no row is
split---and it is applied after deduplication, so the retained filing
determines a firm-year's categories. Because the rules are ordered,
specific provisions are tested before general patterns: GILTI, FDII,
BEAT, and Subpart F are matched by name before the broader foreign-rate
patterns, and the transition tax and rate remeasurement before either,
so that a row naming a TCJA provision is not absorbed by a generic
``foreign'' rule.

Two boundaries are worth stating because they are judgment calls rather
than consequences of the rules. Foreign currency translation and
exchange gains and losses are excluded from \emph{international}: they
are not a differential between tax rates. Withholding taxes are
excluded when the label does not identify them as foreign, since
unqualified ``withholding taxes'' is also a domestic item. Subpart F
inclusions, by contrast, are assigned to \emph{international}, since
they are a current U.S.\ tax on foreign income and the pre-TCJA
counterpart of GILTI, as is BEAT, the third of TCJA's international
provisions.

Two diagnostics bear on whether the mapping leaks. First, over
2013--2024 the residual \emph{other} category carries
\ShareNum{\CatOtherSalesPct}\ percent of the sales-weighted absolute
rate gap, against \ShareNum{\CatIntlSalesPct}\ percent for
\emph{international}; of the rows in that residual, only
\ShareNum{\CatOtherForeignRowPct}\ percent carry a label naming a
foreign concept, and they account for
\ShareNum{\CatOtherForeignPpPct}\ percent of its absolute
contribution. The residual is large, but it is not mostly misfiled
international items: it is dominated by rows the filer itself labeled
``other.'' Second, \ShareNum{\CatNoIntlFyPct}\ percent of firm-years
have no international row at all. This is why a firm-year with no row in
a category is treated differently in different places. The
distributional tables count a category as present only when the
firm-year reports a row there, so a firm reporting no international
line is not a firm reporting an international contribution of zero, and
the number of firm-years underlying a category is smaller than the
panel and differs across categories. The regressions of Section
\ref{sec:Results} assign such a firm-year a contribution of zero, so
that the treatment variable and every outcome are defined on one common
sample.


% 
% APPENDIX TABLE 1: Stage 1 counts
%

\begin{table}[!tp]
\centering
\begin{threeparttable}
\caption{\textsc{Stage 1 Sample Counts}}
\label{tab:StageOneCounts}

   \begin{tabular}{lcr}
   \hline
   \hline
   \noalign{\vskip 5pt}
   \textsc{Samples} &\phantom{hellooo} &  \textsc{Counts} \tabularnewline
   \noalign{\vskip 5pt}
   \hline
   \noalign{\vskip 6pt}
   \textsc{Input URLs}              & &  \CountNum{\InputUrlsTotal} \tabularnewline
   \noalign{\vskip 3pt}
   \textsc{Less duplicate rows}     & &  \CountNum{\InputUrlsDuplicates}\tabularnewline
   \noalign{\vskip 3pt}
   \textsc{Equals unique filings}   & &  \CountNum{\ReconUniqueFilings}\tabularnewline
   \noalign{\vskip 3pt}
   \textsc{Less non-HTML links}     & &  \CountNum{\InputUrlsSkipNonHtml}\tabularnewline
   \noalign{\vskip 3pt}
   \textsc{Less download failures}  & &  \CountNum{\InputUrlsFail}\tabularnewline
   \noalign{\vskip 3pt}
   \textsc{Equals stage 1 extracts} & &  \CountNum{\PhaseOneStatusOK} \tabularnewline
   \noalign{\vskip 6pt}
   \hline
   \end{tabular}

\bigskip
\begin{tablenotes}[flushleft]
\item[] \footnotesize{\textit{Notes:} 
The table reports sample counts for stage 1 extracts of SEC 10K filings.
Input URLs include Compustat-linked 10Ks and manually audited links.
Non-HTML links and download failures are skipped in the stage 1 parsing.}
\end{tablenotes}
\end{threeparttable}

\bigskip
\bigskip
\end{table}



% 
% APPENDIX TABLE 2: Stage 1 Sample Classifications
%

\begin{table}[!tp]
\centering
\begin{threeparttable}
\caption{\textsc{Stage 1 Sample Classifications}}
\label{tab:StageOneClassifications}

   \begin{tabular}{lcr}
   \hline
   \hline
   \noalign{\vskip 5pt}
   \textsc{Classification} & \phantom{hellooo} &  \textsc{Counts} \tabularnewline
   \noalign{\vskip 5pt}
   \hline
   \noalign{\vskip 6pt}
   \textsc{All stage 1 extracts}               & &  \CountNum{\PhaseOneStatusOK} \tabularnewline
   \noalign{\vskip 6pt}
   \textsc{Reconciliation-like buckets}                & &  \CountNum{\PhaseOneExtractableSum}\tabularnewline
   \noalign{\vskip 3pt}
   \ \ \ \ \textsc{High-confidence extractable }       & &  \CountNum{\PhaseOneBucketHighConf} \tabularnewline
   \noalign{\vskip 3pt}
   \ \ \ \ \textsc{Likely-extractable but incomplete}  & &  \CountNum{\PhaseOneBucketLikely}\tabularnewline
   \noalign{\vskip 3pt}
   \ \ \ \ \textsc{Ambiguous}                          & &  \CountNum{\PhaseOneBucketAmbiguous} \tabularnewline
   \noalign{\vskip 6pt}
   \textsc{Other buckets}                              & &  \CountNum{\NoPhaseOneExtraction}\tabularnewline
   \noalign{\vskip 3pt}
   \ \ \ \ \textsc{Low-confidence extractable}         & &  \CountNum{\PhaseOneBucketNarrative}\tabularnewline
   \noalign{\vskip 3pt}
   \ \ \ \ \textsc{No table found}                     & &  \CountNum{\PhaseOneBucketNoTable}\tabularnewline
   \noalign{\vskip 3pt}
   \ \ \ \ \textsc{Wrong table type}                   & &  \CountNum{\PhaseOneBucketWrongTable}\tabularnewline
   \noalign{\vskip 6pt}
   \hline
   \end{tabular}

\bigskip
\begin{tablenotes}[flushleft]
\item[] \footnotesize{\textit{Notes:} 
The table reports sample counts for stage 1 extracts of SEC 10-K filings classified
in terms of potential for retrieving a firm-year reconciliation table in stage 2.
Every stage 1 extract is passed to stage 2 regardless of
its bucket, and the classification records stage 1's assessment.}
\end{tablenotes}
\end{threeparttable}

\bigskip
\bigskip
\end{table}


% 
% APPENDIX TABLE 3: Stage 1 to Stage 2
%

\begin{table}[!tp]
\centering
\begin{threeparttable}
\caption{\textsc{Stage 1 to Stage 2}}
\label{tab:StageOneToTwo}

   \begin{tabular}{lcr}
   \hline
   \hline
   \noalign{\vskip 5pt}
   \textsc{Sample} & \phantom{hellooo} &  \textsc{Counts} \tabularnewline
   \noalign{\vskip 5pt}
   \hline
   \noalign{\vskip 6pt}
   \textsc{Packets processed by stage 2}          & &  \CountNum{\PacketsProcessed}\tabularnewline
   \noalign{\vskip 3pt}
   \ \ \ \ \textsc{Without reconciliation rows}   & &  \CountNum{\ProcessedNoRows}\tabularnewline
   \noalign{\vskip 3pt}
   \ \ \ \ \textsc{With reconciliation rows}      & &  \CountNum{\ProcessedWithRows}\tabularnewline
   \noalign{\vskip 6pt}
   \textsc{Filings entering outside the main stage 2 run} & &  \CountNum{\FilingsOutsideMainRun}\tabularnewline
   \noalign{\vskip 6pt}
   \textsc{Filings with reconciliation rows}      & &  \CountNum{\PhaseTwoUniqueFilings}\tabularnewline
   \noalign{\vskip 6pt}
   \textsc{Total filing-year observations}        & &  \CountNum{\PhaseTwoTotalRows} \tabularnewline
   \noalign{\vskip 6pt}
   \hline
   \end{tabular}

\bigskip
\begin{tablenotes}[flushleft]
\item[] \footnotesize{\textit{Notes:} 
The table reports counts of stage 1 packets with and without at least one
reconciliation row and the number of filing-year observations derived from
packets that had a reconciliation row. Filings entering outside the main
stage 2 run include filings rescued from earlier runs and a targeted
re-extraction of filings whose fiscal year and period of report disagree.
Each filing-year observation is one filing and reconciliation year pair.}
\end{tablenotes}
\end{threeparttable}

\bigskip
\bigskip
\end{table}



% 
% APPENDIX TABLE 5: Stage 2 Parsing 
%

\begin{table}[!tp]
\centering
\begin{threeparttable}
\caption{\textsc{Stage 2 Parsing Results}}
\label{tab:StageTwoParsing}

   \begin{tabular}{lcr}
   \hline
   \hline
   \noalign{\vskip 5pt}
   \textsc{Output} & \phantom{hellooo} &  \textsc{Counts} \tabularnewline
   \noalign{\vskip 5pt}
   \hline
   \noalign{\vskip 6pt}
   \textsc{Total filing-year observations} & &  \CountNum{\PhaseTwoTotalRows} \tabularnewline
   \noalign{\vskip 6pt}
   \textsc{Arithmetic checks:}           & &  \tabularnewline
   \noalign{\vskip 2pt}
   \ \ \ \ \textsc{Successful}           & &  \CountNum{\PhaseTwoCheckTrue}\tabularnewline
   \noalign{\vskip 2pt}
   \ \ \ \ \textsc{Unsuccessful}         & &  \CountNum{\PhaseTwoCheckFalse}\tabularnewline
   \noalign{\vskip 6pt}
   \textsc{Parsing methods:}             & &  \tabularnewline
   \noalign{\vskip 2pt}
   \ \ \ \ \textsc{No LLM assistance}    & &  \CountNum{\PhaseTwoDeterministic}\tabularnewline
   \noalign{\vskip 2pt}
   \ \ \ \ \textsc{LLM assisted}         & &  \CountNum{\PhaseTwoLlmFallback}\tabularnewline
   \noalign{\vskip 6pt}
   \textsc{Reconciliation types:}        & &  \tabularnewline
   \noalign{\vskip 2pt}
   \ \ \ \ \textsc{Tax rates reported}   & &  \CountNum{\PhaseTwoRateForm}\tabularnewline
   \noalign{\vskip 2pt}
   \ \ \ \ \textsc{Tax amounts reported} & &  \CountNum{\PhaseTwoAmountForm}\tabularnewline
   \noalign{\vskip 6pt}
   \hline
   \end{tabular}

\bigskip
\begin{tablenotes}[flushleft]
\item[] \footnotesize{\textit{Notes:} 
The table reports counts of filing-year observations split by
arithmetic checks, parsing methods, and reconciliation types. The
arithmetic check is applied here to the raw extraction and again in
Table \ref{tab:StageTwoToFinal}, after filing-years are deduplicated
and placed on a common schema. The second application is the stricter
of the two and is the one the panel uses:
\CountNum{\ArPassThenFail}\ filing-years pass the check at extraction
and fail it on reassembly, and \CountNum{\ArFailThenPass}\ go the
other way, so the two tables report different counts of the same test.}
\end{tablenotes}
\end{threeparttable}

\bigskip
\bigskip
\end{table}





\begin{table}[!tp]
\centering
\begin{threeparttable}
\caption{\textsc{Stage 2 to Main Firm-level Sample}}
\label{tab:StageTwoToFinal}

   \begin{tabular}{lcr}
   \hline
   \hline
   \noalign{\vskip 5pt}
   \textsc{Sample} & \phantom{hellooo} &  \textsc{Counts} \tabularnewline
   \noalign{\vskip 5pt}
   \hline
   \noalign{\vskip 6pt}
   \textsc{Total filing-year observations}          & &  \CountNum{\PhaseTwoTotalRows} \tabularnewline
   \noalign{\vskip 3pt}
   \ \ \ \ \textsc{Net change in matching to firms} & &  \CountNum{\NetFilingMatch}\tabularnewline
   \noalign{\vskip 6pt}
   \textsc{Equals base panel filing-years}          & &  \CountNum{\BasePanelFilingYears}\tabularnewline
   \noalign{\vskip 3pt}
   \ \ \ \ \textsc{Less: Years before 1998 or after 2024} & &  \CountNum{\DropYearRange}\tabularnewline
   \noalign{\vskip 3pt}
   \ \ \ \ \textsc{Less: Outside expected filing window}  & &  \CountNum{\DropFilingWindow}\tabularnewline
   \noalign{\vskip 3pt}
   \ \ \ \ \textsc{Less: No firm-year link to Compustat}  & &  \CountNum{\DropNoCompustat}\tabularnewline
   \noalign{\vskip 3pt}
   \ \ \ \ \textsc{Less: Unsuccessful arithmetic checks}  & &  \CountNum{\DropCheckFail}\tabularnewline
   \noalign{\vskip 3pt}
   \ \ \ \ \textsc{Less: Duplicate firm-year rows}        & &  \CountNum{\DropDuplicateFirmYears}\tabularnewline
   \noalign{\vskip 6pt}
   \textsc{Equals final firm-year sample}           & &  \CountNum{\AnalysisPanelFirmYears}\tabularnewline
   \noalign{\vskip 6pt}
   \hline
   \end{tabular}

\bigskip
\begin{tablenotes}[flushleft]
\item[] \footnotesize{\textit{Notes:} 
The table reports counts for the final firm-year sample. Filings are matched to
\citet{CompustatAnnual} firms to form the base panel. A combined 10-K filed by
several co-registrants covers more than one firm; because stage 2 recovers a
single reconciliation from each document, the filing is assigned to the
registrant whose reported taxes match it and the remaining registrants are
dropped. Duplicate firm-year rows arise because each 10-K reports three
reconciliation years, so the same firm-year appears in several filings; we keep
one observation per firm and year. }
\end{tablenotes}
\end{threeparttable}

\bigskip
\bigskip
\end{table}



% 
% APPENDIX TABLE 7: Industry SIC Codes 
%

\begin{table}[!tp]
\centering
\begin{threeparttable}
\caption{\textsc{Industry SIC Codes}}
\label{tab:SICcodes}

   \begin{tabular}{lcr}
   \hline
   \hline
   \noalign{\vskip 5pt}
   \textsc{Industry Category} & \phantom{helloo} &  \textsc{SIC Codes} \tabularnewline
   \noalign{\vskip 5pt}
   \hline
   \noalign{\vskip 6pt}
   \textsc{Technology}       & &  3570--3579, 3670--3679, 7370--7379 \tabularnewline
   \noalign{\vskip 3pt}
   \textsc{Pharmaceuticals}  & &  2830--2836, 8731 \tabularnewline
   \noalign{\vskip 3pt}
   \textsc{Manufacturing }   & &  2000--3999 excl.~codes above \tabularnewline
   \noalign{\vskip 3pt}
   \textsc{Oil \& Gas}       & &  1310--1389, 2910--2999 \tabularnewline
   \noalign{\vskip 3pt}
   \textsc{Retail}           & &  5200--5999 \tabularnewline
   \noalign{\vskip 3pt}
   \textsc{Utilities}        & &  4900--4939 \tabularnewline
   \noalign{\vskip 3pt}
   \textsc{Other}            & &  not elsewhere categorized \tabularnewline
   \noalign{\vskip 6pt}
   \hline
   \end{tabular}


\bigskip
\begin{tablenotes}[flushleft]
\item[] \footnotesize{\textit{Notes:} 
For industry assignments, we use SIC codes from the \citet{CompustatAnnual}
database. For our full sample, we exclude heavily regulated financial
institutions, holding companies, and funds, namely,
6000--6799, except major group 61 other than 6111 and 6159, and except 
6211, 6282, 6411, 6512, 6531, and 6532. }
\end{tablenotes}
\end{threeparttable}

\bigskip
\bigskip
\end{table}



\section{Additional Results}

In this appendix, we report additional difference-in-differences results,
namely, estimated contributions of all reconciliation categories
and profit-location responses for different subsets of firms as
we vary characteristics.


\subsection{Contributions by Reconciliation Category\label{app:channels}}

Table \ref{tab:ChannelDiD} estimates equation (\ref{eq:PooledDiD}) with each
reconciliation category's contribution to the tax rate gap as the outcome, and
reports the sum of those coefficients against $\hat\beta_{\text{ETR}}$. The
effective tax rate is the statutory rate plus the sum of the category
contributions, and the industry-by-year effects absorb the statutory change, so
the two should agree up to the winsorization applied to each category.
Section \ref{sec:EffectiveTaxRate} reads the result.


%
% TABLE: Pooled DiD by reconciliation category
%
\begin{table}[!tp]
\centering
\begin{threeparttable}
\caption{\textsc{Pooled Difference-in-Differences Estimates \\ by Reconciliation Category}}
\label{tab:ChannelDiD}

   \begin{tabular}{clcrc}
   \hline
   \hline
   \noalign{\vskip 5pt}
   \phantom{h} & \textsc{Category} &\phantom{hello}
                    &\multicolumn{1}{c}{\textsc{Coefficient}} &\phantom{h}\tabularnewline
   \noalign{\vskip 5pt}
   \hline
   \noalign{\vskip 6pt}
   &\textsc{International}   & & \CoefNum{\ChIntlCoef}\StarBox{\ChIntlStars} &\tabularnewline
                       & & & \SENum{\ChIntlSe}\StarBox{} & \tabularnewline
    \noalign{\vskip 3pt}
   &\textsc{State \& local}   & & \CoefNum{\ChStateLocalCoef}\StarBox{\ChStateLocalStars} &\tabularnewline
                       & & & \SENum{\ChStateLocalSe}\StarBox{} & \tabularnewline
    \noalign{\vskip 3pt}
   &\textsc{Federal credits}   & & \CoefNum{\ChFedCredCoef}\StarBox{\ChFedCredStars} &\tabularnewline
                       & & & \SENum{\ChFedCredSe}\StarBox{} & \tabularnewline
    \noalign{\vskip 3pt}
   &\textsc{Book v.\ tax}   & & \CoefNum{\ChBookTaxCoef}\StarBox{\ChBookTaxStars} &\tabularnewline
                       & & & \SENum{\ChBookTaxSe}\StarBox{} & \tabularnewline
    \noalign{\vskip 3pt}
   &\textsc{Valuation}   & & \CoefNum{\ChValnCoef}\StarBox{\ChValnStars} &\tabularnewline
                       & & & \SENum{\ChValnSe}\StarBox{} & \tabularnewline
    \noalign{\vskip 3pt}
   &\textsc{One time}   & & \CoefNum{\ChOneTimeCoef}\StarBox{\ChOneTimeStars} &\tabularnewline
                       & & & \SENum{\ChOneTimeSe}\StarBox{} & \tabularnewline
    \noalign{\vskip 3pt}
   &\textsc{Other}   & & \CoefNum{\ChResidCoef}\StarBox{\ChResidStars} &\tabularnewline
                       & & & \SENum{\ChResidSe}\StarBox{} & \tabularnewline
    \noalign{\vskip 3pt}
   \noalign{\vskip 3pt}
   \hline
   \noalign{\vskip 5pt}
   &\textsc{Sum of categories}    & & \CoefNum{\ChSumCoef}\phantom{***}\phantom{h} &\tabularnewline
   &\textsc{Effective tax rate} & & \CoefNum{\EtrCoef}\StarBox{\EtrStars} &\tabularnewline
   \noalign{\vskip 3pt}
   &\textsc{Sum excluding international} & & \CoefNum{\ChNonIntlCoef}\phantom{***}\phantom{h} &\tabularnewline
   \noalign{\vskip 6pt}
   \hline
   \end{tabular}

\bigskip
\begin{tablenotes}[flushleft]
\item[] \footnotesize{\textit{Notes:} Each row reports the pooled difference-in-differences
estimate of equation (\ref{eq:PooledDiD}) with that category's contribution to the tax rate gap
as the outcome, on \CountNum{\ChFirms}\ firms and \CountNum{\ChObs}\ firm-years.  A firm-year
with no row in a category is assigned a contribution of zero, so all categories are estimated on one
common sample.  Because the effective tax rate is the statutory rate plus the sum of the category
contributions, and the industry-by-year effects absorb the statutory change, the sum of the
category coefficients equals $\hat\beta_{\text{ETR}}$ up to the winsorization applied to each
category; the two are reported together as a check.  Standard errors clustered at the firm level
are reported in parentheses. $^{*} p<0.10$, $^{**} p<0.05$, $^{***} p<0.01$.}
\end{tablenotes}
\end{threeparttable}

\bigskip
\bigskip
\end{table}

\subsection{Heterogeneity in the Profit-Location Response\label{app:heterogeneity}}

Table \ref{tab:Heterogeneity} reports the interaction coefficient $\hat\gamma$
from equation (\ref{eq:DiDvariant}) with the foreign profit share as the
outcome, one characteristic $Z_i$ at a time. Continuous characteristics are
measured over 2013--2016, winsorized at the 1st and 99th percentiles, and
standardized, so $\hat\gamma$ is the change in the profit-location response
associated with a one standard deviation higher $Z$. The two binary measures
are indicators, so $\hat\gamma$ is the difference between firms that have such
a subsidiary and firms that do not.

%
% TABLE: heterogeneity in the profit-location response
%
\begin{table}[!tp]
\centering
\begin{threeparttable}
\caption{\textsc{Heterogeneity in the Profit-Location Response}}
\label{tab:Heterogeneity}

   \begin{tabular}{clcrrc}
   \hline
   \hline
   \noalign{\vskip 5pt}
   \phantom{h} & \textsc{Characteristic} $Z_i$ &\phantom{hello}
                    &\multicolumn{1}{c}{$\hat\gamma$}
                    &\multicolumn{1}{c}{\textsc{Firm-years}} &\phantom{h}\tabularnewline
   \noalign{\vskip 5pt}
   \hline
   \noalign{\vskip 6pt}
   &\textsc{R\&D / assets} & & \CoefNum{\HetRd}\StarBox{\HetRdStars}
                          & \CountNum{\HetRdObs}\phantom{hell} &\tabularnewline
                       & & & \SENum{\HetRdSe}\StarBox{} & & \tabularnewline
    \noalign{\vskip 3pt}
   &\textsc{Cash / assets} & & \CoefNum{\HetCash}\StarBox{\HetCashStars}
                          & \CountNum{\HetCashObs}\phantom{hell} &\tabularnewline
                       & & & \SENum{\HetCashSe}\StarBox{} & & \tabularnewline
    \noalign{\vskip 3pt}
   &\textsc{Intangibles ex-goodwill / assets} & & \CoefNum{\HetIntan}\StarBox{\HetIntanStars}
                          & \CountNum{\HetIntanObs}\phantom{hell} &\tabularnewline
                       & & & \SENum{\HetIntanSe}\StarBox{} & & \tabularnewline
    \noalign{\vskip 3pt}
   &\textsc{Foreign revenue share} & & \CoefNum{\HetFgnRev}\StarBox{\HetFgnRevStars}
                          & \CountNum{\HetFgnRevObs}\phantom{hell} &\tabularnewline
                       & & & \SENum{\HetFgnRevSe}\StarBox{} & & \tabularnewline
    \noalign{\vskip 3pt}
   &\textsc{Any haven or conduit subsidiary} & & \CoefNum{\HetHaven}\StarBox{\HetHavenStars}
                          & \CountNum{\HetHavenObs}\phantom{hell} &\tabularnewline
                       & & & \SENum{\HetHavenSe}\StarBox{} & & \tabularnewline
    \noalign{\vskip 3pt}
   &\textsc{Any classic-haven subsidiary} & & \CoefNum{\HetClassicHaven}\StarBox{\HetClassicHavenStars}
                          & \CountNum{\HetClassicHavenObs}\phantom{hell} &\tabularnewline
                       & & & \SENum{\HetClassicHavenSe}\StarBox{} & & \tabularnewline
    \noalign{\vskip 3pt}
   &\textsc{Log haven/conduit subsidiary count} & & \CoefNum{\HetHavenCount}\StarBox{\HetHavenCountStars}
                          & \CountNum{\HetHavenCountObs}\phantom{hell} &\tabularnewline
                       & & & \SENum{\HetHavenCountSe}\StarBox{} & & \tabularnewline
    \noalign{\vskip 3pt}
   &\textsc{Log assets} & & \CoefNum{\HetLogAt}\StarBox{\HetLogAtStars}
                          & \CountNum{\HetLogAtObs}\phantom{hell} &\tabularnewline
                       & & & \SENum{\HetLogAtSe}\StarBox{} & & \tabularnewline
    \noalign{\vskip 3pt}
   \noalign{\vskip 3pt}
   \hline
   \end{tabular}

\bigskip
\begin{tablenotes}[flushleft]
\item[] \footnotesize{\textit{Notes:} Each row is a separate estimate of
equation (\ref{eq:DiDvariant}) with the foreign profit share as the outcome and
that characteristic as $Z_i$, including firm and industry-by-year fixed
effects. Continuous characteristics are measured pre-TCJA firm averages over 2013--2016, winsorized
at the 1st and 99th percentiles, and standardized. Coefficients are in
percentage points of the foreign profit share. Samples differ across rows
because the characteristics have different coverage. Standard errors clustered
at the firm level are reported in parentheses. $^{*} p<0.10$,
$^{**} p<0.05$, $^{***} p<0.01$.}
\end{tablenotes}
\end{threeparttable}

\bigskip
\bigskip
\end{table}








\end{document}


